% Chapter 4 -- Methodology

\section{Introduction}
\label{sec:ch4-intro}

Chapter~3 stated eight functional requirements that partition the query-level control flow of CAEM, from pre-routing confidence estimation through ablation instrumentation. This chapter operationalises those requirements into a working system. Each mechanism described here corresponds to a concrete block of the reference implementation on the Branch~C Qwen-2.5-3B-Instruct backbone, and each equation is the equation the code evaluates at run time. The chapter is constructed so that two kinds of reader are served: one who reads linearly and assembles a complete mental model of CAEM, and one who reads only the figures, tables, and captions and still recovers the architecture.

The CAEM pipeline executes in eight stages for every arriving query, bracketed by an outer self-improvement loop that re-verifies the episodic memory and fine-tunes the base model at fixed cycle boundaries. Stage~1 estimates, in two sequential forward passes of the decoder-only backbone (a short greedy decode for a token-probability aggregate and a non-generative pass for a layer-variance convergence ratio), how confidently the base model could answer the query using its internal parameters alone. Stage~2 retrieves the single nearest stored episode from the episodic memory. Stage~3 dispatches the query to one of three processing tiers on the basis of a blended score over pre-route confidence and the retrieved episode's stored composite, with a safety override that forces external grounding whenever pre-route confidence falls below a minimum floor. Stage~4 is the chosen tier. Tier~1 returns a stored answer untouched; Tier~2 generates from the base model without retrieved passages, using a one-shot scaffolded chain-of-thought prompt and a forced decoder prefix; Tier~3 generates under grounding against a preprocessed Wikipedia passage index, retrieving passages with dense-passage-retrieval embeddings, cross-encoder reranking to a top-3 context, and conditioning the same Qwen-2.5-3B-Instruct generator on those passages. Stage~5 subjects every Tier~2 and Tier~3 answer to a nine-signal verifier -- three internal, three sample-set, three grounding, plus a question--answer relevance signal that closes a specific off-topic failure mode and a contradiction-probability diagnostic -- whose composite score $u_{\text{stored}}$ is the memory's storage criterion. Grounding-family signals are evaluated under task-conditional hypothesis construction (directional for 3-way claims, option-text-substituted for multi-choice letters, literal text for factoid and long-form queries) so that the same entailment judge provides a calibrated $[0,1]$ signal across heterogeneous benchmarks. Stage~6 emits one of four categorical decisions: \textsc{Store}, \textsc{Deferred}, \textsc{Abstain}, or \textsc{Discard}. Stage~7 commits the episode to memory, flags it for training eligibility, or returns a principled refusal, as dictated by Stage~6. Stage~8 runs at cycle boundaries rather than per query: it re-verifies every stored entry under the freshly fine-tuned weights, prunes entries whose composite has fallen below a retroactive threshold, curates a training pool from the survivors, fine-tunes the base model under an $L_{2}$-anchor regulariser, and rolls the cycle back if a held-out MMLU retention probe falls below the forgetting bound. Between cycles, the thresholds that gate \textsc{Store} / \textsc{Deferred} and training eligibility are re-fit per cycle by an EMA-smoothed quantile rule on the most recent calibration fold, so that CAEM's decision surface tracks the generator's shifting distribution without requiring offline intervention.

The eight stages correspond one to one with Functional Requirements FR1 through FR7 of Chapter~3, with FR8 (ablation and evaluation instrumentation) realised by the experimental harness described in Chapter~5 rather than by a pipeline stage. Section~\ref{sec:system-overview} presents the complete architecture in a single figure, introduces the notation used throughout the rest of the chapter, and walks a single fabricated query end to end through all eight stages to ground the abstract description in a concrete trace. Sections~\ref{sec:pre-route} through~\ref{sec:self-improvement} then treat the stages one by one at implementation fidelity. Section~\ref{sec:theoretical-analysis} states the three core theorems (Data Purity, Monotonicity, Convergence) together with the Branch~C additions that refine them (Theorem~T4's benchmark-fixed asymptotic hallucination-elimination bound and Corollaries~C7--C10, which treat the parameter-bounded open-domain residual, multi-base-model extension, epistemic-floor framing, and the retrieval-bounded structural lower bound). Section~\ref{sec:implementation} gives the reproducibility-critical implementation details, including the cuDNN-SDPA attention backend and \texttt{torch.compile} on the generator forward pass that together yield a $2.4\times$ wall-clock speedup on Qwen-2.5-3B relative to eager execution. Chapter~5 lays out the experimental protocol under which that chapter reports its results, and Section~\ref{sec:ch4-summary} closes the chapter.

\section{System Overview}
\label{sec:system-overview}

Figure~\ref{fig:caem-pipeline} shows the per-query path through CAEM on a single page. The between-cycle self-improvement loop that re-verifies memory and fine-tunes the base model is described in Section~\ref{sec:self-improvement}. Every symbol used in the rest of the chapter is defined in Table~\ref{tab:notation}, and a single fabricated query is traced end to end through all eight stages in Example~\ref{ex:walkthrough} so that the reader acquires the mechanisms in concrete form before meeting them in the abstract.

% ------------------------------------------------------------------
% Figure 4.1 : Full CAEM pipeline (wide spacing, non-overlapping labels)
% ------------------------------------------------------------------
\begin{figure}[p]
\centering
\resizebox{1.08\textwidth}{!}{%
\begin{tikzpicture}[
    xshift=-0.9cm,
    font=\normalsize,
  % ---- Node styles ----
  stage/.style  = {rectangle, draw=teal!65!black, line width=0.9pt, rounded corners=5pt,
                   minimum height=1.5cm, minimum width=6.5cm, align=center, fill=teal!8},
  router/.style = {rectangle, draw=teal!65!black, line width=0.9pt, rounded corners=5pt,
                   minimum height=1.8cm, minimum width=8.5cm, align=center, fill=yellow!18},
  mcyl/.style   = {cylinder, draw=teal!65!black, line width=0.9pt,
                   shape border rotate=90, aspect=0.25,
                   minimum height=3.0cm, minimum width=3.2cm, align=center, fill=teal!12},
  tA/.style     = {rectangle, draw=green!55!black, line width=1pt, rounded corners=5pt,
                   minimum height=1.7cm, minimum width=4.0cm, align=center, fill=green!12},
  tB/.style     = {rectangle, draw=red!55!orange, line width=1pt, rounded corners=5pt,
                   minimum height=1.7cm, minimum width=4.0cm, align=center, fill=red!12!orange!10},
  tC/.style     = {rectangle, draw=red!60!black, line width=1pt, rounded corners=5pt,
                   minimum height=1.7cm, minimum width=4.0cm, align=center, fill=red!10},
  verify/.style = {rectangle, draw=teal!65!black, line width=1pt, rounded corners=5pt,
                   minimum height=4.2cm, minimum width=9.5cm, align=center, fill=teal!5},
  decide/.style = {rectangle, draw=teal!65!black, line width=0.9pt, rounded corners=5pt,
                   minimum height=1.8cm, minimum width=9.5cm, align=center, fill=yellow!18},
  oServe/.style = {rectangle, draw=green!60!black, line width=1pt, rounded corners=4pt,
                   minimum height=1.4cm, minimum width=3.6cm, align=center,
                   fill=green!30, font=\small\scshape},
  oStore/.style = {rectangle, draw=green!60!black, line width=0.9pt, rounded corners=3pt,
                   minimum height=1.2cm, minimum width=2.7cm, align=center,
                   fill=green!28, font=\small\scshape},
  oDefer/.style = {rectangle, draw=green!45!black, line width=0.9pt, rounded corners=3pt,
                   minimum height=1.2cm, minimum width=2.7cm, align=center,
                   fill=green!12, font=\small\scshape},
  oAbst/.style  = {rectangle, draw=orange!75!black, line width=0.9pt, rounded corners=3pt,
                   minimum height=1.2cm, minimum width=2.7cm, align=center,
                   fill=orange!28, font=\small\scshape},
  oDisc/.style  = {rectangle, draw=red!60!black, line width=0.9pt, rounded corners=3pt,
                   minimum height=1.2cm, minimum width=2.7cm, align=center,
                   fill=red!28, font=\small\scshape},
  simp/.style   = {rectangle, draw=violet!60!black, line width=1pt, rounded corners=5pt,
                   minimum height=3.2cm, minimum width=17.0cm, align=center, fill=violet!7},
  % ---- Arrow styles ----
  farr/.style   = {-Latex, line width=0.9pt},
  garr/.style   = {-Latex, line width=1.1pt, green!50!black},
  oarr/.style   = {-Latex, line width=1.1pt, color=red!55!orange},
  rarr/.style   = {-Latex, line width=1.1pt, red!60!black},
  darr/.style   = {<->, dashed, line width=0.9pt, teal!55!black},
  varr/.style   = {-Latex, line width=1pt, dashed, violet!60!black},
    elab/.style   = {font=\small\itshape, inner sep=2pt, fill=none, draw=none}
]

% ==================================================================
% NODES   (main trunk centred at x = 6; y compressed to fit A4 page)
% ==================================================================

\node[stage]  (q)   at (6,   0.0)  {\textbf{Query $q$}};

\node[stage]  (s1)  at (6,  -2.7)  {%
    \textbf{S1\quad Pre-route confidence}\\[4pt]
    $u_{\text{pre}} = 0.60\,u_{\text{token}} + 0.40\,(1{+}C_{\text{conv}})^{-1}$\\
    {\scriptsize two forward passes;\;$T$ re-fit per cycle}};

\node[stage]  (s2)  at (6,  -5.8)  {%
    \textbf{S2\quad Episodic memory retrieval}\\[4pt]
    SBERT 768-d cosine nearest-neighbour search};

% Episodic memory cylinder  (right of S2, connected by dashed bidirectional arrow)
\node[mcyl]   (mem) at (16.5,-5.8) {%
    \textbf{Episodic}\\[2pt]\textbf{Memory}\\[3pt]$K{=}10^{6}$};

\node[router] (s3)  at (6, -9.2)  {%
    \textbf{S3\quad Three-tier router}\\[4pt]
    $S(q)=0.70\cdot\mathrm{sim}(q,e^{*})+0.30\cdot u_{\text{stored}}(e^{*})$\\
    {\scriptsize safety override:\;if $u_{\text{pre}}<0.60$ the score is bypassed and Tier~3 is forced}};

% ---- Tier row  (T1 far-left / T2 centre / T3 far-right)
%  S5 min-width = 9.5 cm centred at x=6  ->  left edge  ~ x = 1.25
%  T1 placed at x=0 is therefore outside S5, allowing a clean bypass path.
\node[tA] (t1) at (0.0, -12.2)  {%
    \textbf{Tier 1}\\[3pt]return $e^{*}$ answer\\{\scriptsize(from verified memory)}};
\node[tB] (t2) at (6.0, -12.2)  {%
    \textbf{Tier 2}\\[3pt]one-shot CoT generation\\{\scriptsize(base model only)}};
\node[tC] (t3) at (12.0,-12.2)  {%
    \textbf{Tier 3}\\[3pt]RAG: DPR top-5\\{\scriptsize$\to$ cross-encoder rerank to top-3}};

% ---- S4: Tier execution grouping (compact left-side label; no wrapper) ----
\node[draw=none, inner sep=6pt, fit=(t1)(t2)(t3),
      label={[font=\small\bfseries, teal!65!black, anchor=east, xshift=-6pt]left:S4}]
     (s4box) {};

% Tier 1 direct-serve terminal  (left side, same y-band as S5)
\node[oServe] (serve) at (-3.5,-17.0) {%
    \textsc{Answer served}\\{\scriptsize from verified memory}};

% S5 centred at x=6, min-width=9.5 cm  ->  spans ~ x=1.25..10.75
\node[verify] (s5)  at (6, -17.0)  {%
    \textbf{S5\quad Nine-signal Verifier}\\[6pt]
    \textit{Internal:}\quad$u_{\text{token}},\;u_{\text{dropout}},\;u_{\text{internal}}$\\[3pt]
    \textit{Sample-set:}\quad$s_{\text{avg}},\;p_{\text{entail}},\;h_{\text{norm}}$\\[3pt]
    \textit{Grounding:}\quad$p_{\text{ground}}^{\max},\;p_{\text{ground}}^{\text{mean}},\;p_{\text{ground}}^{\text{atomic}}$\\[5pt]
    {\footnotesize\textbf{confabulation gate:}\;
     $u_{\text{internal}}{\geq}0.70\;\wedge\;p_{\text{ground}}^{\max}{\leq}0.20
     \;\Rightarrow\;$\textsc{Discard (early exit)}}};

\node[decide] (s6)  at (6, -21.0)  {%
    \textbf{S6\quad Four-outcome storage decision}\\[5pt]
    {\scriptsize
    $u_{\text{stored}}{\geq}0.65\Rightarrow\textsc{Store}$\quad
    $u_{\text{stored}}{\geq}0.45\Rightarrow\textsc{Deferred}$\quad
    $p_{\text{ground}}^{\max}{<}0.20\Rightarrow\textsc{Abstain}$\quad
    \textit{else:}\;\textsc{Discard}}};

% Terminal outcome boxes  (y = -24.5)
\node[oStore] (out1) at (2.0, -24.5)  {\textsc{Store}\\{\tiny write + train-flag}};
\node[oDefer] (out2) at (5.0, -24.5)  {\textsc{Deferred}\\{\tiny write, no train}};
\node[oAbst]  (out3) at (8.0, -24.5)  {\textsc{Abstain}\\{\tiny principled refusal}};
\node[oDisc]  (out4) at (11.0,-24.5)  {\textsc{Discard}\\{\tiny suppressed}};

% ---- S7: Memory write grouping (compact left-side label; no wrapper) ----
\node[draw=none, inner sep=6pt, fit=(out1)(out2)(out3)(out4),
      label={[font=\small\bfseries, teal!65!black, anchor=east, xshift=-6pt]left:S7}]
     (s7box) {};

% S8: between-cycle self-improvement  (wide band at bottom)
\node[simp] (s8) at (6, -27.5) {%
    \textbf{S8\quad Between-cycle self-improvement loop}
    \quad\textit{(runs once at each cycle boundary, not per query)}\\[6pt]
    (1)~Retroverify all stored entries
        ($\tau_{\text{retro}}{=}0.50$;\;upward-only: never demotes a previously verified entry)\\[3pt]
    (2)~Curate training pool ($\tau_{\text{train}}{=}0.75$;\;90:10 task-to-general mix)\quad
    (3)~Fine-tune:\;$\mathcal{L}(\theta)=\mathcal{L}_{\text{CE}}
        +\tfrac{\lambda}{2}\|\theta-\theta_{\text{prev}}\|^2$,\;$\lambda{=}0.01$\\[3pt]
    (4)~Abort guard:\;$\rho=A_{\text{post}}/A_{\text{pre}}\geq0.93$;\;
        roll back to $\theta_{\text{prev}}$ if violated};

% ==================================================================
% ARROWS
% ==================================================================

% ---- Main trunk ----
\draw[farr] (q.south)  -- (s1.north);
\draw[farr] (s1.south) -- node[elab,right]{$u_{\text{pre}}$} (s2.north);
\draw[farr] (s2.south) -- node[elab,right]{$u_{\text{pre}}$ + query embedding} (s3.north);

% ---- S2 <-> Episodic memory  (bidirectional dashed) ----
\draw[darr] (s2.east)
    -- node[elab,above,pos=0.55,align=center]{%
        retrieve $e^{*}$\;:\;$\mathrm{sim}(q,e^{*})$,\\$u_{\text{stored}}(e^{*})$}
    (mem.west);

% ---- S3 -> three tiers ----
\coordinate (br) at (6, -10.9);   % branch waypoint below S3
\draw[farr] (s3.south) -- (br);

% Tier 1 (green): branch LEFT, condition label above horizontal leg
\draw[garr]
    (br) --
    node[elab,above,pos=0.45]{\textbf{if}\;$S(q)\geq0.90$}
    (0.0,-10.9) -- (t1.north);

% Tier 2 (orange): straight down, condition label to the RIGHT
\draw[oarr]
    (br) --
    node[elab,right,pos=0.5]{\textbf{elif}\;$\mathrm{sim}(q,e^{*})\geq0.75$}
    (t2.north);

% Tier 3 (red): branch RIGHT, condition label above horizontal leg
\draw[rarr]
    (br) --
    node[elab,above,pos=0.75]{\textbf{else};\;or forced when\;$u_{\text{pre}}<0.60$}
    (12.0,-10.9) -- (t3.north);

% ---- Tier 1 bypass: route LEFT then DOWN, well outside S5 ----
% S5 left edge ≈ x=1.25; this path uses x=-3.5 to stay clear.
\coordinate (t1byp1) at (0.0, -14.0);    % just below T1
\coordinate (t1byp2) at (-3.5,-14.0);    % far left
\draw[garr]
    (t1.south) -- (t1byp1) --
    node[elab,below,pos=0.5]{\textit{bypass S5 and S6: episode already verified at storage time}}
    (t1byp2) -- (serve.north);

% ---- T2 -> S5 ----
\draw[oarr]
    (t2.south) --
    node[elab,right,pos=0.45]{answer + reasoning chain}
    (t2 |- s5.north);

% ---- T3 -> S5 routed on the right to avoid center-label overlap ----
\draw[rarr]
    (t3.south) --
    (12.0,-14.2) --
    (13.6,-14.2) --
    node[elab,right,pos=0.55]{answer + reasoning chain}
    (13.6,-16.9) --
    (s5.east);

% ---- S5 -> S6 ----
\draw[farr]
    (s5.south) --
    node[elab,right]{outputs:\;$u_{\text{stored}},\;p_{\text{contra}},\;p_{\text{ground}}^{\max}$}
    (s6.north);

% ---- S6 fan-out -> terminal boxes ----
\coordinate (br2) at (6,-22.8);
\draw[farr] (s6.south) -- (br2);
\draw[farr] (br2) -| (out1.north);
\draw[farr] (br2) -| (out2.north);
\draw[farr] (br2) -| (out3.north);
\draw[farr] (br2) -| (out4.north);

% ---- Store + Deferred accumulate into S8 at cycle end ----
% Drop only 0.3 cm so the horizontal segment stays ABOVE s8.north,
% then the vertical leg enters s8 cleanly from the top.
\draw[varr]
    (out1.south) -- ++(0,-0.3) -|
    node[elab,left,pos=0.05]{\textit{cycle end}}
    ($(s8.north)+(-2.5,0)$);
\draw[varr]
    (out2.south) -- ++(0,-0.3) -|
    ($(s8.north)+(-0.8,0)$);

% ---- S8 -> S1 : OUTER TEN-CYCLE LOOP ----
%
%  After S8 concludes, the updated weights θ_new (or θ_prev on rollback)
%  are loaded and Stage 1 of Cycle n+1 begins.  This arrow is the outer loop.
%
%  Route: S8.west  →  far left (x ≈ -5.5)  →  up to y of S1  →  S1.west
%
\draw[varr]
    (s8.west)
    -- ++(-5.5, 0) coordinate (ll)
    -- node[elab,left,pos=0.45,align=left]{%
        \textbf{Cycle $n{+}1$ begins}\\[3pt]
        S1 receives\\
        $\theta_{\text{new}}$ (fine-tuned)\\[1pt]
        or $\theta_{\text{prev}}$ (restored\\
        if abort guard fired)}
    (ll |- s1.west)
    -- (s1.west);

\end{tikzpicture}%
}
\caption{%
    CAEM pipeline overview (S1--S8): each query passes through confidence estimation, memory retrieval, routing, tier execution, verification, and storage decision.
    Stage~8 runs once per cycle boundary for retroverification and $L_2$-anchored self-improvement.%
}
\label{fig:caem-pipeline}
\end{figure}

% ------------------------------------------------------------------
% Table 4.1 : Notation used throughout Chapter 4
% ------------------------------------------------------------------
\begin{table}[tbp]
\centering
\caption{Notation used throughout Chapter~4. All probability-style signals lie in $[0,1]$ and are either native probabilities (NLI outputs) or normalised into $[0,1]$.}
\label{tab:notation}
\small
\begin{tabular}{@{}p{2.2cm}p{2.1cm}p{9.4cm}@{}}
\toprule
\textbf{Symbol} & \textbf{Range} & \textbf{Description} \\
\midrule
\multicolumn{3}{@{}l}{\textit{Confidence signals}} \\
\addlinespace[2pt]
$u_{\text{token}}$        & $[0,1]$ & Token-probability aggregate over greedy prefix. \\
$C_{\text{conv}}$         & $\mathbb{R}_{\geq 0}$ & Transformer decoder-stack layer convergence ratio. \\
$u_{\text{pre}}$          & $[0,1]$ & Pre-route confidence (Stage 1 output, temperature-scaled after Cycle~0). \\
$u_{\text{dropout}}$      & $[0,1]$ & Monte Carlo dropout disagreement over $K_D$ stochastic passes. \\
$u_{\text{internal}}$     & $[0,1]$ & Internal confidence fusion of $u_{\text{token}}$ and $u_{\text{dropout}}$. \\
$s_{\text{avg}}$          & $[0,1]$ & Mean pairwise cosine similarity over $M$ sampled reasoning chains. \\
$p_{\text{entail}}$       & $[0,1]$ & Mean NLI entailment probability chain-to-answer. \\
$h_{\text{norm}}$         & $[0,1]$ & Normalised semantic entropy over $K_S$ chains clustered by bidirectional entailment. \\
$p_{\text{ground}}^{\max}$ & $[0,1]$ & Maximum entailment answer-to-passage over top-3 reranked passages. \\
$p_{\text{ground}}^{\text{mean}}$ & $[0,1]$ & Mean entailment answer-to-passage over top-3 reranked passages. \\
$p_{\text{ground}}^{\text{atomic}}$ & $[0,1]$ & Weakest-link entailment over decomposed atomic facts. \\
$q\_a\_\text{relevance}$  & $[0,1]$ & Cross-encoder relevance on (question, answer). \\
$p_{\text{contra}}$       & $[0,1]$ & Contradiction probability (diagnostic; not read by Stage~6). \\
$u_{\text{stored}}$       & $[0,1]$ & Composite verifier confidence written to memory. \\
\addlinespace[2pt]
\multicolumn{3}{@{}l}{\textit{Routing}} \\
\addlinespace[2pt]
$\mathrm{sim}(q, e)$      & $[-1,1]$ & Cosine similarity in sentence-embedding space. \\
$e^{*}$                   & ---       & Single nearest stored episode to $q$. \\
$S(q)$                    & $[0,1]$   & Combined routing score $\alpha\,\mathrm{sim}(q,e^{*}) + (1-\alpha)\,u_{\text{stored}}(e^{*})$. \\
$\alpha$                  & $0.70$    & Routing similarity blend weight. \\
\addlinespace[2pt]
\multicolumn{3}{@{}l}{\textit{Thresholds}} \\
\addlinespace[2pt]
$\tau_{\text{store}}$     & $0.65$ & Composite threshold for \textsc{Store}. \\
$\tau_{\text{defer}}$     & $0.45$ & Composite threshold for \textsc{Deferred}. \\
$\tau_{\text{abs}}$       & $0.20$ & Grounding-floor threshold for \textsc{Abstain}. \\
$\tau_{\text{retro}}$     & $0.50$ & Retroactive re-verification prune threshold. \\
$\tau_{\text{train}}$     & $0.75$ & Training-eligibility threshold. \\
$\tau_{\text{forget}}$    & $0.93$ & Forgetting-bound (cycle retention). \\
\addlinespace[2pt]
\multicolumn{3}{@{}l}{\textit{Self-improvement and memory}} \\
\addlinespace[2pt]
$\theta,\ \theta_{\text{prev}}$ & ---       & Model parameters at current step and at cycle start (anchor). \\
$\lambda$                  & $0.01$    & L2-anchor coefficient. \\
$\rho$                     & $\mathbb{R}_{\geq 0}$ & Retention ratio $\mathrm{MMLU}_{\text{post}}/\mathrm{MMLU}_{\text{pre}}$. \\
$K_{\text{mem}}$           & $10^{6}$  & Episodic memory capacity. \\
$N_{\text{train}}$         & $20{,}000$ & FAISS exact-to-IVF-PQ promotion point. \\
$n_{\text{list}}$          & $4096$    & IVF coarse centroid count (memory index). \\
$K_D$                      & $5$       & MC dropout passes for $u_{\text{dropout}}$. \\
$K_S$                      & $10$      & Semantic entropy sample count for $h_{\text{norm}}$. \\
$M$                        & $3$       & Chain count for $s_{\text{avg}}$. \\
\bottomrule
\end{tabular}
\end{table}

\subsection{Stage-by-stage narrative}
\label{sec:stage-narrative}

\paragraph{Stage 1: Pre-route confidence.}
Two sequential forward passes of the decoder-only backbone compute the stage output. The first is a short greedy decode of at most $32$ tokens, whose per-step log-probabilities combine into the token-probability aggregate $u_{\text{token}}$. The second is a non-generative forward pass with hidden states exposed, whose per-layer outputs are split at the midpoint and combined into the layer-variance convergence ratio $\Cconv$. The two fuse into the raw pre-route confidence $\tilde{u}_{\text{pre}} = 0.60\,u_{\text{token}} + 0.40\,(1{+}\Cconv)^{-1}$ and are passed through a scalar temperature scaler in logit space. The temperature defaults to $T = 1$ at initialisation and is re-fit at every cycle boundary by L-BFGS on a 500-sample calibration slice disjoint from both the SIL training pool and the evaluation set; a production-mode flag disables the re-fit for deployments without labelled calibration data. $\upre$ is the only quantity Stage~3 consumes to evaluate the safety override, and one of two that enter the combined routing score.

\paragraph{Stage 2: Episodic memory retrieval.}
The query is encoded by Sentence-BERT \texttt{all-mpnet-base-v2} into a $768$-dimensional $L_{2}$-normalised vector and submitted to a FAISS index (a flat inner-product backend for memories below the promotion point $N_{\text{train}} = 20{,}000$, an \texttt{IndexIVFPQ} with $n_{\text{list}} = 4096$ coarse centroids and $n_{\text{probe}} = 32$ search-time probes once promoted). The top-$1$ neighbour $e^{*}$ is returned with its cosine similarity and a pointer into a metadata map that carries the entry's stored composite $u_{\text{stored}}(e^{*})$ and its other signals. No generation has yet occurred.

\paragraph{Stage 3: Three-tier router.}
A safety override is evaluated first: if $\upre < 0.60$, the router forces Tier~3 regardless of memory content. Otherwise, a combined score $S(q) = 0.70\,\mathrm{sim}(q, e^{*}) + 0.30\,u_{\text{stored}}(e^{*})$ is computed with a fixed blend weight. Tier~1 is taken if $S(q) \geq 0.90$; Tier~2 if that threshold is not met but $\mathrm{sim}(q, e^{*}) \geq 0.75$; Tier~3 otherwise. The two thresholds are separated so that a high stored composite cannot arithmetically compensate for a low similarity -- an intentional refusal to let memory quality substitute for actual match.

\paragraph{Stage 4: Tier execution.}
Tier~1 returns the stored answer attached to $e^{*}$ directly and emits a \textsc{Store}-equivalent outcome without passing through the verifier, because $u_{\text{stored}}(e^{*})$ already cleared the composite gate at the original storage event and the entry is re-verified at cycle boundaries. Tier~2 generates from the Qwen-2.5-3B-Instruct backbone using a ChatML-wrapped, one-shot scaffolded chain-of-thought prompt and a forced \texttt{Reasoning:} decoder prefix -- not literally zero-shot, but not retrieval-augmented either. Tier~3 retrieves the top-$5$ passages from a DPR-preprocessed Wikipedia index (21~million passages on Branch~C), cross-encoder reranks them to the top-$3$ with \texttt{bge-reranker-v2-m3}, and conditions the same Qwen-2.5-3B generator on those passages under the Tier-3 prompt template. The generator backbone is shared between tiers, so the cuDNN-SDPA attention backend and \texttt{torch.compile} on the forward pass (Section~\ref{sec:implementation}) apply uniformly to every tier that generates.

\paragraph{Stage 5: Unified nine-signal verification.}
Tier~2 and Tier~3 answers pass through a single verifier that computes nine signals in three families, plus a question--answer relevance signal and a contradiction diagnostic. The internal family $\{u_{\text{token}}, u_{\text{dropout}}, u_{\text{internal}}\}$ measures the model's self-assessed confidence on the specific answer. The sample-set family $\{s_{\text{avg}}, p_{\text{entail}}, h_{\text{norm}}\}$ measures cross-chain agreement over $M{=}3$ resampled reasoning chains plus $K{=}10$ semantic-entropy samples. The grounding family $\{p_{\text{ground}}^{\max}, p_{\text{ground}}^{\text{mean}}, p_{\text{ground}}^{\text{atomic}}\}$ measures how well the answer is entailed by external Wikipedia evidence; the atomic variant prompts the generator to decompose the answer into standalone factual claims and records the weakest-link entailment across atoms, catching partial-fabrication failures that the max and mean aggregates cannot distinguish. A tenth signal, $q\_a\_\text{relevance}$, is a cross-encoder score on the (question, answer) pair and closes the off-topic failure mode where a hallucinated answer scored high on entailment because the retrieved passage happened to match the hallucination rather than the question. The grounding signals are evaluated against a task-conditional hypothesis: $3$-way claims (FEVER-style SUPPORTS / REFUTES / NEI) use directional entailment with negation generation so that refutation separates from paraphrase noise; multi-choice letters (ARC-Challenge) are substituted with their option text before entailment scoring, turning the bare label into a propositional statement the judge can score; factoid and long-form queries use the literal answer as the hypothesis. The verifier also records $p_{\text{contra}}$ as a diagnostic, though under the shipped MiniCheck judge it is structurally zero and no decision branch reads it. A confabulation early-exit gate evaluated before the composite discards the answer immediately whenever internal confidence is high and grounding is effectively absent ($u_{\text{internal}} \geq 0.70$ and $p_{\text{ground}}^{\max} \leq 0.20$). Surviving answers are combined into the composite $u_{\text{stored}}$ by an eight-weight convex combination (Section~\ref{sec:verifier}).

\paragraph{Stage 6: Four-outcome storage decision.}
The decision tree reads only $u_{\text{stored}}$ and $p_{\text{ground}}^{\max}$: if $u_{\text{stored}} \geq \tau_{\text{store}}$ the outcome is \textsc{Store}; if $\tau_{\text{defer}} \leq u_{\text{stored}} < \tau_{\text{store}}$ it is \textsc{Deferred}; if $u_{\text{stored}} < \tau_{\text{defer}}$ and $p_{\text{ground}}^{\max} < \tau_{\text{abs}}$ it is \textsc{Abstain} -- no evidence either way, emit a principled refusal rather than a low-confidence answer; all other cases emit \textsc{Discard}. Grounding weakness is absorbed into $u_{\text{stored}}$ through the $p_{\text{ground}}^{\text{mean}}$ and $p_{\text{ground}}^{\max}$ terms rather than through a separate hard check, and $p_{\text{contra}}$ is recorded as a diagnostic but not read by the decision tree under the shipped MiniCheck judge. The thresholds $\tau_{\text{store}} = 0.65$, $\tau_{\text{defer}} = 0.45$, $\tau_{\text{abs}} = 0.20$ are fitted quantile estimates of the Cycle-0 calibration distribution and are re-fit at every cycle boundary by an EMA-smoothed quantile rule (Section~\ref{sec:self-improvement}).

\paragraph{Stage 7: Memory write and output.}
A \textsc{Store} decision writes the full episode to memory -- question, reasoning chain, answer, $768$-dimensional $L_{2}$-normalised SBERT embedding, source-benchmark tag (used by Stage~8 to gate training eligibility through the ID/OOD split), storage cycle index, timestamp, the full nine-signal verifier record plus $p_{\text{contra}}$ and $q\_a\_\text{relevance}$, the categorical decision, and the early-exit flag -- and flags it as training-eligible provided $u_{\text{stored}} \geq \tau_{\text{train}} = 0.75$ \emph{and} the source benchmark is in the training pool (not a held-out transfer benchmark). A \textsc{Deferred} decision writes the episode identically but withholds the training flag, pending promotion or pruning at the next cycle's retroactive re-verification pass. \textsc{Abstain} returns a principled refusal to the caller and writes nothing. \textsc{Discard} suppresses the generation entirely. Every write is gated by the novelty filter of Section~\ref{sec:memory}, which rejects near-duplicates at cosine similarity above $0.95$ so that retrieval-feedback updates on the existing entry replace redundant writes.

\paragraph{Stage 8: Between-cycle self-improvement.}
At every cycle boundary, the full verifier is re-run over every stored episode under the updated weights. Updates are applied asymmetrically by default: re-verification only \emph{raises} a stored composite and replaces the full signal record on promotion; it never writes back a score below the entry's original, so that a single bad cycle cannot silently demote a previously well-verified episode. A retroverify loop-filter pass runs first and drops entries whose stored reasoning chains trip the loop-repetition detector before any verifier forward pass is spent on them. Irrecoverable survivors (new $u_{\text{stored}} < \tau_{\text{retro}} = 0.50$) are pruned outright. A training pool is then curated from training-eligible stores (those with $u_{\text{stored}} \geq \tau_{\text{train}} = 0.75$ and a source benchmark in the training pool), mixed with general-domain examples in a $90{:}10$ ratio, and the base model is fine-tuned under an $L_{2}$-anchor regulariser $\mathcal{L}(\theta) = \mathcal{L}_{\text{CE}} + \tfrac{\lambda}{2}\|\theta - \theta_{\text{prev}}\|^{2}$ with $\lambda = 0.01$. A per-cycle MMLU evaluation computes the retention ratio $\rho = \mathrm{MMLU}_{\text{post}}/\mathrm{MMLU}_{\text{pre}}$; if $\rho < 0.93$ the post-cycle weights are discarded and $\theta_{\text{prev}}$ is restored (a hard rollback), otherwise committed. The decision-tree thresholds ($\tau_{\text{store}}, \tau_{\text{defer}}, \tau_{\text{train}}$) and the Stage-1 temperature scaler are re-fit on the new cycle's calibration slice in the same pass, so that the full decision surface tracks the generator's shifting distribution without offline intervention. The loop is scheduled for ten cycles but is permitted to early-stop at cycle~$\geq 5$ when an MMLU-plateau or rollback condition triggers; Chapter~5 records the actual completed cycle count for baseline-matched comparisons.

\subsection{Worked example}
\label{sec:walkthrough}

\refstepcounter{example}\label{ex:walkthrough}
\begin{tcolorbox}[
  enhanced,
  breakable,
  width=\textwidth,
  colback=teal!3,
  colframe=teal!55!black,
  colbacktitle=teal!55!black,
  coltitle=white,
  fonttitle=\bfseries,
  title={Example~\theexample\quad End-to-end trace through CAEM on two contrasting queries},
  boxrule=0.6pt,
  arc=3pt,
  left=14pt, right=14pt, top=10pt, bottom=12pt,
  titlerule=0pt,
  attach boxed title to top left={yshift=-2mm, xshift=10pt},
  boxed title style={arc=2pt, boxrule=0pt, colback=teal!55!black},
  toptitle=2pt, bottomtitle=2pt,
  enlarge top at break by=0mm,
  enlarge bottom at break by=0mm,
]

{\color{green!45!black}\rule{6pt}{6pt}}\ \textbf{Query~A\ \ (Tier~1 cheap path)}\par\medskip

Consider $q = $ ``\emph{Who discovered radium?}'' arriving at Cycle~3, with the memory already populated by roughly $17{,}000$ entries from earlier cycles.

\medskip
\noindent\textcolor{teal!55!black}{\textbf{Stage~1 \textbar{} Pre-route confidence.}}\ One encoder pass produces $u_{\text{token}}=0.82$ over the greedy prefix and $C_{\text{conv}}=0.41$, giving
\[
u_{\text{pre}} \;=\; 0.60 \cdot 0.82 \;+\; 0.40 \cdot \tfrac{1}{1.41} \;=\; 0.776
\]
after temperature scaling fit at the end of Cycle~0. The safety override is not triggered.

\medskip
\noindent\textcolor{teal!55!black}{\textbf{Stage~2 \textbar{} Memory retrieval.}}\ SBERT nearest neighbour returns $e^{*}$ with $\mathrm{sim}(q, e^{*}) = 0.93$ and $u_{\text{stored}}(e^{*}) = 0.86$.

\medskip
\noindent\textcolor{teal!55!black}{\textbf{Stage~3 \textbar{} Three-tier router.}}\ The combined score is
\[
S(q) \;=\; 0.70 \cdot 0.93 \;+\; 0.30 \cdot 0.86 \;=\; 0.909 \;\geq\; 0.90,
\]
so the router dispatches to \textbf{Tier~1}.

\medskip
\noindent\textcolor{green!45!black}{\textbf{Stage~4 \textbar{} Tier~1 execution.}}\ The stored answer attached to $e^{*}$, ``\emph{Marie Curie and Pierre Curie}'', is returned directly without invoking the verifier or any generation. This is the cheap path the verified-memory investment is designed to enable.

\medskip
\noindent\textcolor{green!45!black}{\textbf{Stages~5--7 \textbar{} Verifier bypassed.}}\ Tier~1 entries carry a composite from their original storage event and are re-verified between cycles; the outcome is \textsc{Store}-equivalent and the user is served immediately from verified memory.

\vspace{10pt}
\noindent\textcolor{teal!55!black}{\rule{\linewidth}{0.4pt}}
\vspace{6pt}

{\color{orange!75!black}\rule{6pt}{6pt}}\ \textbf{Query~B\ \ (Tier~2 full verify-and-decide path, at the end of Cycle~3)}\par\medskip

A later query $q'=$ ``\emph{Who first synthesised aspirin?}'' routes instead to Tier~2; the one-shot scaffolded-CoT generator produces the extracted answer ``\emph{Felix Hoffmann}''. The verifier's record is:

\medskip
\begin{center}
\renewcommand{\arraystretch}{1.2}
\begin{tabular}{@{}l l @{\hspace{18pt}} l l @{\hspace{18pt}} l l@{}}
\toprule
\multicolumn{2}{@{}l}{\textit{Internal}} & \multicolumn{2}{l}{\textit{Sample-set}} & \multicolumn{2}{l@{}}{\textit{Grounding}} \\
\midrule
$u_{\text{internal}}$ & $0.78$ & $s_{\text{avg}}$     & $0.74$ & $p_{\text{ground}}^{\max}$          & $0.88$ \\
                      &        & $p_{\text{entail}}$  & $0.81$ & $p_{\text{ground}}^{\text{mean}}$   & $0.71$ \\
                      &        & $h_{\text{norm}}$    & $0.22$ & $p_{\text{ground}}^{\text{atomic}}$ & $0.65$ \\
\bottomrule
\end{tabular}\quad
\begin{tabular}{@{}l l@{}}
\toprule
\multicolumn{2}{@{}l@{}}{\textit{Relevance / diag.}} \\
\midrule
$q\_a\_\text{relevance}$ & $0.55$ \\
$p_{\text{contra}}$       & $0.04$ \\
\bottomrule
\end{tabular}
\end{center}

\medskip
\noindent\textcolor{orange!75!black}{\textbf{Confabulation gate.}}\ Not triggered, since $p_{\text{ground}}^{\max}=0.88 \gg 0.20$.

\medskip
\noindent\textcolor{orange!75!black}{\textbf{Composite (eight-weight, Branch~C).}}
\begin{align*}
u_{\text{stored}}
&= 0.22\,(0.71) + 0.18\,(0.88) + 0.06\,(0.65) + 0.14\,(0.81) \\
&\quad + 0.20\,(0.55) + 0.10\,(0.74) + 0.08\,(0.78) + 0.02\,(1 - 0.22) \\
&= \mathbf{0.729}.
\end{align*}

\medskip
\noindent\textcolor{orange!75!black}{\textbf{Decision.}}\ $p_{\text{contra}}$ is recorded for diagnostic auditing but no decision branch reads it under the shipped MiniCheck judge. Because $u_{\text{stored}} = 0.729 \geq \tau_{\text{store}} = 0.65$, the decision is \textsc{Store}. Because $u_{\text{stored}} < \tau_{\text{train}} = 0.75$, it is \emph{not} flagged as training-eligible; it will be reconsidered at the next cycle's retroactive re-verification pass.

\end{tcolorbox}

The rest of Chapter~4 fills in each stage at the fidelity required to re-implement it from this text. Sections~\ref{sec:pre-route} through~\ref{sec:self-improvement} follow the eight-stage order of Figure~\ref{fig:caem-pipeline}.

\section{Pre-Routing Confidence Estimation}
\label{sec:pre-route}

Stage~1 produces a calibrated scalar $\upre \in [0,1]$ that estimates how confidently the base model could answer the incoming query from its internal parameters alone. The stage runs before any passage retrieval, any full-answer generation, and any storage decision, so that its cost is paid uniformly across all three downstream tiers -- including Tier~1 queries that will ultimately be served without producing a new answer. Two design commitments shape the computation. First, the estimate must be obtainable from signals that do not require sampling multiple completions, so that the per-query latency budget is bounded by a fixed small number of forward passes rather than by resample count. On the decoder-only Branch~C backbone (Qwen-2.5-3B-Instruct), Stage~1 uses a short greedy decode of at most $32$ tokens for the token-probability signal and one additional non-generative forward pass for the convergence signal, giving a cost floor of two sequential passes per query. Second, $\upre$ must discriminate reliably at its low end, since it is consumed both by the combined routing score of Stage~3 (Section~\ref{sec:router}) and by the safety override that forces Tier~3 retrieval whenever $\upre$ falls below $0.60$. Confidence signals that are well calibrated in the high region but noisy near zero are unsuitable for this role, because a safety override is a hard block on the cheap paths and its false positives cost nothing while its false negatives cost hallucinations.

Stage~1 combines two complementary signals produced by the two sequential passes. The first signal, $u_{\text{token}}$, is a token-level aggregate over a short greedy-decoded answer prefix; it captures lexical confidence in the surface form the model would have produced. The second signal, $\Cconv$, is a layer-variance ratio computed over the query's transformer-stack hidden states under a forward pass \emph{without} generation; it captures representational stability before any token is emitted. Both passes are preceded by a ChatML wrap that places the query inside a user turn and appends the assistant generation-prompt marker; without this envelope, per-token probabilities on an instruction-tuned decoder-only model scatter across the full vocabulary and collapse $u_{\text{token}}$ to numerically negligible values (empirically confirmed on Qwen-2.5-3B during pre-launch smoke testing). The two signals disagree systematically in the failure mode Stage~1 most needs to catch: a query that is superficially familiar in phrasing but substantively out of distribution yields a high $u_{\text{token}}$ over whatever surface form the greedy decoder commits to, but a hidden-state stack whose late layers have not yet converged on a stable representation yields an elevated $\Cconv$. Fusing both signals ensures that $\upre$ is suppressed in exactly this regime.

\paragraph{Token-probability aggregate.}
The first pass is a greedy decode of at most $32$ tokens. Decoding is deterministic (\texttt{do\_sample=False}) so that $u_{\text{token}}$ is reproducible for a given query and model state; the prefix cap is enforced by the \texttt{max\_new\_tokens} argument, and generation also terminates early at the EOS token. EOS itself, when emitted, is excluded from the aggregate. Writing $\ell_{1},\dots,\ell_{n}$ for the non-EOS log-probabilities of the emitted prefix, the aggregate is their geometric mean expressed as an exponentiated arithmetic mean of log-probabilities,
\begin{equation}
u_{\text{token}} \;=\; \exp\!\Bigl(\tfrac{1}{n}\sum_{i=1}^{n}\ell_{i}\Bigr)
\;=\; \Bigl(\prod_{i=1}^{n} p(y_{i}\mid y_{<i}, q)\Bigr)^{1/n},
\label{eq:utoken}
\end{equation}
so that $u_{\text{token}} \in (0,1]$ with the natural interpretation of mean per-token probability. The geometric-mean form is chosen over the arithmetic mean of probabilities because log-probabilities decompose additively across steps, making the aggregate insensitive to prefix length: a $10$-token and a $30$-token answer with the same mean per-token probability receive the same score. The prefix cap of $32$ tokens is imposed so that Stage~1 cost remains bounded even on queries whose greedy completion runs long; empirically the cap is reached only on enumerations, and the confidence signal derived from the first $32$ tokens is a faithful proxy for the full answer. Any exception raised by the generator -- for instance, a CUDA out-of-memory event or a degenerate input -- is caught and returns $u_{\text{token}} = 0$, which under the Stage~3 safety override forces the query onto Tier~3 and preserves a fail-safe contract that no Tier-1 hit is ever served from a failed estimate.

\paragraph{Layer convergence ratio.}
After the token pass, a second forward call -- this time without generation -- runs the base model over the same ChatML-wrapped query with \texttt{output\_hidden\_states=True}. For a decoder-only backbone, the returned tuple contains $L+1$ tensors of shape \texttt{(batch, seq\_len, hidden\_dim)}: index~$0$ is the input embedding layer, and indices~$1$ through~$L$ are the outputs of the $L$ transformer decoder layers ($L = 36$ on Qwen-2.5-3B). The embedding layer $h_{0}$ is excluded so that the ratio measures contextualisation rather than the tokenizer's inductive bias. The remaining stack is split at its midpoint into an early half $\{h_{1},\dots,h_{L/2}\}$ and a late half $\{h_{L/2+1},\dots,h_{L}\}$, and the convergence ratio is defined as the ratio of scalar activation variances,
\begin{equation}
\Cconv \;=\; \frac{\mathrm{Var}\bigl(h_{1},\dots,h_{L/2}\bigr)}{\mathrm{Var}\bigl(h_{L/2+1},\dots,h_{L}\bigr) \,+\, \varepsilon},
\qquad \varepsilon = 10^{-8},
\label{eq:cconv}
\end{equation}
where variance is computed over the flattened concatenation of each half's hidden-state tensors (no pooling across the sequence dimension, so long queries contribute proportionally more samples). A well-converged representation is one in which the late layers vary much less than the early layers as the stack progressively stabilises the meaning of the query, producing $\Cconv > 1$. A poorly-converged representation, the situation in which late layers are still reshaping the query, produces $\Cconv$ near unity or below. The midpoint split is a design choice motivated by the intended decomposition of the decoder's work: the first half accumulates lexical and syntactic structure while the second half resolves it against the model's parametric knowledge. Any failure mode in which the model does not expose hidden states, or the forward pass itself raises, is caught and returns $\Cconv = 0$, which under the conversion below maps to the maximum contribution of the $\Cconv$ signal to the weighted sum; this is the conservative choice because $u_{\text{token}}$ still acts as a hard floor through the other summand.

\paragraph{Fusion and temperature calibration.}
The two signals are fused into a raw pre-route confidence by the linear combination specified in Equation~\ref{eq:upre} of Chapter~3, reproduced here with the transformation that maps $\Cconv$ into the unit interval made explicit,
\begin{equation}
\tilde{u}_{\text{pre}} \;=\; w_{\text{tok}} \cdot u_{\text{token}} \;+\; w_{\text{conv}} \cdot \frac{1}{1 + \Cconv}, \qquad w_{\text{tok}} = 0.60,\; w_{\text{conv}} = 0.40,
\label{eq:upre-raw}
\end{equation}
with $\tilde{u}_{\text{pre}}$ then clipped to $[0,1]$ to absorb the floating-point tail of either component. The clipping is mechanically necessary because the finite precision of the softmax and the hidden-state variance can each produce values marginally outside the unit interval, but in practice it fires only on degenerate queries where it is already immaterial which tier executes.

The raw estimate $\tilde{u}_{\text{pre}}$ is a biased estimator of empirical correctness: language models are systematically over-confident, and a weighted sum of two such signals is over-confident in the same direction. To remove the bias, $\tilde{u}_{\text{pre}}$ is passed through a single-parameter temperature scaler. The scaler operates in logit space:
\begin{equation}
\upre \;=\; \sigma\!\Bigl(\frac{\mathrm{logit}(\tilde{u}_{\text{pre}})}{T}\Bigr),
\qquad \sigma(x) = \frac{1}{1+e^{-x}},
\label{eq:temp-scale}
\end{equation}
where $T$ is the scalar temperature stored on the pipeline configuration object. This is the standard post-hoc recalibration of Guo et al.~\cite{pmlr-v70-guo17a}, written in sigmoid form for a scalar probability rather than in softmax form for a multi-class logit vector; the two are equivalent and the choice is notational. The temperature defaults to $T = 1$ at pipeline initialisation, which makes the scaler the identity and leaves the raw estimate unchanged; this is the operating behaviour during Cycle~0, before any calibration split has been evaluated. $T$ is subsequently fit by minimising the binary negative log-likelihood on a disjoint $500$-query calibration slice that is strictly non-overlapping with the SIL training pool and with the held-out evaluation set (this three-way isolation is enforced by an assertion in \texttt{run\_experiment.assert\_disjoint\_calibration} before SIL begins), using empirical exact-match correctness as the binary target:
\begin{equation}
T^{*} \;=\; \arg\min_{T > 0} \; -\frac{1}{N}\sum_{i=1}^{N}\Bigl[ y_{i} \log \sigma\bigl(\mathrm{logit}(\tilde{u}_{\text{pre}}^{(i)})/T\bigr) + (1 - y_{i}) \log\bigl(1 - \sigma(\mathrm{logit}(\tilde{u}_{\text{pre}}^{(i)})/T)\bigr)\Bigr].
\label{eq:temp-fit}
\end{equation}
NLL is used as the training objective in place of the expected calibration error (ECE) reported in Chapter~3's acceptance criterion, because ECE is a binned statistic and not differentiable; minimising NLL produces a temperature that minimises ECE in practice while remaining tractable for L-BFGS. By default $T^{*}$ is re-fit at each cycle boundary, because SIL fine-tuning shifts the generator's logit scale and the Cycle-0 scalar becomes stale~\citep{ovadia2019trust,thulasidasan2019mixup}; the \texttt{skip\_per\_cycle\_temperature} flag on \texttt{CAEMConfig} disables the per-cycle re-fit for production deployments in which labelled calibration data is unavailable, in which case $T$ is pinned to the most recent fitted value. The $u_{\text{stored}}$ composite weights are fixed by design and are \emph{not} re-fit, avoiding the small-$n$ overfitting risk of re-training the eight-weight composite on a 500-sample slice. Chapter~3's NFR on calibration stability (ECE $\leq 0.05$ after Cycle~1) is an acceptance check on this fit rather than an inline regularisation.

Algorithm~\ref{alg:prerout} summarises the stage. The two forward passes are explicit as distinct steps; both are preceded by the shared ChatML tokenisation step. The algorithm assumes that the temperature $T$ has been fit; in the very first cycle the temperature defaults to unity and the sigmoid-of-logit identity leaves the raw estimate unchanged, so the calibration is a strict improvement applied from Cycle~1 onward rather than a correctness dependency.

\begin{algorithm}
\caption{Stage~1 pre-route confidence $\upre$.}
\label{alg:prerout}
\begin{algorithmic}[1]
\Require Query $q$; decoder-only base model $f_{\theta}$ with $L$ transformer layers; weights $w_{\text{tok}} = 0.60$, $w_{\text{conv}} = 0.40$; temperature $T$; prefix cap $M = 32$.
\Ensure Calibrated scalar $\upre \in [0,1]$.
\State $q_{\text{wrap}} \gets$ ChatML wrap of $q$ (user turn + assistant generation-prompt marker)
\State \Comment{Pass 1: greedy decode of a short prefix for $u_{\text{token}}$}
\State $(y_{1{:}n},\; \ell_{1{:}n}) \gets$ greedy decode of $f_{\theta}(q_{\text{wrap}})$ up to $M$ tokens, recording emitted-token log-probabilities $\ell_{i}$
\State exclude EOS from $\ell_{1{:}n}$; if empty, \Return $0$ \Comment{safety-override eligible}
\State $u_{\text{token}} \gets \exp\bigl(\tfrac{1}{n}\sum_{i}\ell_{i}\bigr)$ \Comment{geometric mean of token probabilities}
\State \Comment{Pass 2: non-generative forward pass for $\Cconv$}
\State $h_{0{:}L} \gets$ forward pass of $f_{\theta}(q_{\text{wrap}})$ with \texttt{output\_hidden\_states=True}
\State $\mathrm{Var}_{\text{early}} \gets \mathrm{Var}(h_{1{:}L/2})$; \; $\mathrm{Var}_{\text{late}} \gets \mathrm{Var}(h_{L/2+1{:}L})$ \Comment{flattened across seq\_len $\times$ hidden\_dim; embedding layer $h_{0}$ excluded}
\State $\Cconv \gets \mathrm{Var}_{\text{early}} \,/\, (\mathrm{Var}_{\text{late}} + 10^{-8})$
\State $\tilde{u}_{\text{pre}} \gets \mathrm{clip}\bigl(w_{\text{tok}}\, u_{\text{token}} + w_{\text{conv}}\,(1 + \Cconv)^{-1},\; 0,\; 1\bigr)$
\State $\upre \gets \sigma\bigl(\mathrm{logit}(\tilde{u}_{\text{pre}}) / T\bigr)$ \Comment{identity when $T = 1$}
\State \Return $\upre$
\end{algorithmic}
\end{algorithm}

The scalar $\upre$ produced by Algorithm~\ref{alg:prerout} is the sole quantity Stage~3 uses to evaluate the safety override, and one of two quantities (the other being the stored composite of the retrieved episode) that enter the combined routing score. Its production is the only computation in the entire pipeline that is paid on every query regardless of tier, which sets the cost floor of CAEM at two sequential forward passes per query and motivates the no-sampling, short-decode design of this section.

\section{Episodic Memory}
\label{sec:memory}

The episodic memory is the persistent store that CAEM consults at Stage~2 of every query, writes to at Stage~7 after a \textsc{Store} outcome (or, at a cycle boundary, when a deferred entry is promoted out of the buffer of Section~\ref{sec:deferred-reconsider}), and re-verifies wholesale at Stage~8 between cycles. It is the central artefact that accumulates across cycles: a well-designed memory makes the cheap Tier~1 path increasingly attractive, while a poorly curated memory either forces unnecessary retrieval (by failing to serve near-duplicates of previous queries) or pollutes the Tier~1 path with low-quality answers. This section specifies what is stored, how it is indexed, how novelty and capacity are managed, how retrieval feedback and cycle-boundary consolidation update stored confidence, and how the popularity signal feeds the Stage~8 retroverification queue.

\paragraph{Data schema.}
Each memory entry is an \textsc{EpisodicEntry} instance holding a block of immutable content fields, a mutable quality record written by the verifier, and a mutable usage record written by the retrieval and retroverification paths. Table~\ref{tab:entry-schema} gives the full field list. The immutable block carries the natural-language triple (question, reasoning chain, answer), a 768-dimensional $L_{2}$-normalised SBERT embedding, the storage cycle index, the Unix timestamp at write time, and a \texttt{source\_benchmark} tag identifying the benchmark the episode originated from. The source-benchmark tag is load-bearing for Stage~8: the SIL training-pool gate excludes episodes whose source benchmark is in the held-out transfer panel, so that CAEM never trains on data drawn from a benchmark it is evaluated on. Invariance across cycles is load-bearing for a different reason: Stage~8 derives the fine-tuning pool from these fields, and mutating them mid-cycle would mean training against a moving target with no reliable audit trail of which (query, chain, answer) triples actually updated the model.

The mutable quality record carries the nine verifier signals (three internal, three sample-set, three grounding), the Branch~C question--answer relevance signal $q\_a\_\text{relevance}$, the contradiction diagnostic $p_{\text{contra}}$, the composite $u_{\text{stored}}$, the categorical Stage~6 decision, and the early-exit flag from the confabulation gate. The mutable usage record carries the cumulative retrieval count, the running success rate, a \texttt{retroverified} flag set at each cycle boundary, and a Branch~C popularity counter \texttt{hit\_counter} that records Tier-1 serves since the last retroverify. The popularity counter is distinct from the cumulative retrieval count: it is reset to zero each time the entry is re-scored and is consumed by the Stage~8 ``popular-needs-recheck'' queue described later in this section. A seventh mutable field, \texttt{merged\_source\_benchmarks}, records the benchmark tags absorbed by the memory-consolidation pass when a cluster of paraphrase-equivalent episodes is collapsed onto a single representative; the SIL training-pool gate reads the effective benchmark set as $\{\texttt{source\_benchmark}\} \cup \texttt{merged\_source\_benchmarks}$ and excludes the entry if any element of that set is in the transfer panel, so consolidation cannot leak held-out content into training through the back door of ``the representative happened to be in-distribution''.

The composite $u_{\text{stored}}$ is defined by the eight-weight convex combination of Chapter~3, Equation~3.5, repeated here for reference,
\begin{equation}
\begin{aligned}
u_{\text{stored}} \;=\;\;&
 0.22\,p_{\text{ground}}^{\text{mean}}
\;+\; 0.18\,p_{\text{ground}}^{\max}
\;+\; 0.06\,p_{\text{ground}}^{\text{atomic}}
\;+\; 0.14\,p_{\text{entail}} \\ &
\;+\; 0.20\,q\_a\_\text{relevance}
\;+\; 0.10\,s_{\text{avg}}
\;+\; 0.08\,u_{\text{internal}}
\;+\; 0.02\,(1 - h_{\text{norm}}),
\end{aligned}
\label{eq:ustored}
\end{equation}
with $(1 - h_{\text{norm}})$ appearing because semantic entropy is a disagreement signal and the composite is a confidence, and the weights summing to $1.00$. The grounding family is the largest single block (weight mass $0.46$) because grounding is the one lane that can falsify a confidently-wrong but internally-consistent answer; the question--answer relevance signal carries the single largest weight ($0.20$) because the Branch~C Cycle-0 audit identified off-topic high-entailment hallucinations as the dominant early-cycle failure mode, and the relevance term is the one signal that can distinguish ``the passage entails the answer'' from ``the answer is on-topic for the question''. Storage happens only if $u_{\text{stored}} \geq \tau_{\text{defer}}$ at Stage~6, so every entry in memory carries a positive prior that survives a two-threshold screen; this prior is the quantity Stages~3 and~5 subsequently consume.

\begin{table}[tbp]
\centering
\caption{Fields of an \textsc{EpisodicEntry}. Immutable content feeds the Stage~8 fine-tuning pool; mutable quality is updated by retroactive re-verification; mutable usage is updated on each retrieval and on memory consolidation.}
\label{tab:entry-schema}
\small
\begin{tabular}{@{}p{4.0cm}p{10.2cm}@{}}
\toprule
\textbf{Field} & \textbf{Contents} \\
\midrule
\multicolumn{2}{@{}l}{\textit{Immutable content (never updated after storage)}} \\
\addlinespace[2pt]
Triple & question, reasoning chain, answer \\
Embedding & 768-d L2-normalised SBERT vector (\texttt{float32}) \\
Provenance & storage cycle, Unix timestamp \\
Source benchmark & origin tag; gates SIL training via the ID/OOD split \\
\addlinespace[4pt]
\multicolumn{2}{@{}l}{\textit{Mutable quality record (updated by retroactive re-verification)}} \\
\addlinespace[2pt]
Nine signals & $u_{\text{token}}, u_{\text{dropout}}, u_{\text{internal}}, s_{\text{avg}}, h_{\text{norm}}, p_{\text{entail}}, p_{\text{ground}}^{\max}, p_{\text{ground}}^{\text{mean}}, p_{\text{ground}}^{\text{atomic}}$ \\
Relevance & $q\_a\_\text{relevance}$ (Branch~C cross-encoder) \\
Contradiction & $p_{\text{contra}}$ (diagnostic; not read by Stage~6) \\
Composite & $u_{\text{stored}}$ (Equation~\ref{eq:ustored}) \\
Decision & \textsc{Store} / \textsc{Deferred} / \textsc{Abstain} / \textsc{Discard}; early-exit flag \\
\addlinespace[4pt]
\multicolumn{2}{@{}l}{\textit{Mutable usage record (retrieval, popularity, consolidation)}} \\
\addlinespace[2pt]
Retrieval count & cumulative integer $r$ \\
Success rate & running fraction of retrievals accepted downstream \\
Hit counter & Tier-1 serves since last retroverify (reset on re-score) \\
Retroverified flag & set at each cycle boundary; consumed by the Value score \\
Merged benchmarks & tags absorbed by consolidation; gates SIL training alongside source benchmark \\
\bottomrule
\end{tabular}
\end{table}

\paragraph{Embedding model and similarity metric.}
Embeddings are produced by the Sentence-BERT \texttt{all-mpnet-base-v2} model, which maps arbitrary queries into a 768-dimensional vector space. Every embedding is $L_{2}$-normalised to unit length at the moment of production, so that the inner product between two embeddings equals their cosine similarity,
\begin{equation}
\mathrm{sim}(q, e) \;=\; \hat{q}^{\top} \hat{e}, \qquad \hat{v} = v / \|v\|_{2},
\label{eq:cosine}
\end{equation}
which is the property FAISS's \texttt{METRIC\_INNER\_PRODUCT} relies on for cosine-ordered nearest-neighbour results. Normalisation is applied once at storage time and once at query time; it is never re-applied after retrieval, so similarities returned by FAISS are consumed directly by Stage~3's routing score without further rescaling. A defensive re-normalisation pass in the encoder guards against floating-point drift in specific \texttt{sentence-transformers} versions where internal \texttt{normalize\_embeddings=True} occasionally leaves vectors at norm slightly off unity.

\paragraph{Index structure and size-adaptive promotion.}
The memory is backed by a FAISS index that supports in-place removal during pruning. For the first phase of operation the backend is an \texttt{IndexIDMap} wrapping an \texttt{IndexFlatIP}, which performs exact cosine nearest-neighbour search at $O(N)$ per query. When the number of stored vectors first reaches a promotion threshold $N_{\text{train}} = 20{,}000$, the index is rebuilt in place as an \texttt{IndexIDMap2} wrapping an \texttt{IndexIVFPQ} with $n_{\text{list}} = 4096$ coarse centroids, product quantisation into $m = 64$ sub-vectors of $8$ bits each, and search-time probing of $n_{\text{probe}} = 32$ centroids; the switch from \texttt{IndexIDMap} to \texttt{IndexIDMap2} is required so that IVF-PQ entries remain individually removable under the pruning and consolidation paths. This promotion is the operationalisation of Chapter~3's memory-cost analysis: flat IP is cheap at $N \leq 20{,}000$ but scales poorly; IVF-PQ trades a small approximation error for sublinear retrieval and a $\sim 32{\times}$ reduction in embedding storage cost, which is what makes the $K_{\text{mem}} = 10^{6}$ target tractable on a single 24\,GB device. Before promotion the quantiser is trained on a uniformly sampled subset of up to $200{,}000$ existing vectors, which is sufficient because $n_{\text{list}}$ is the binding constraint (FAISS requires at least as many training points as centroids, satisfied at the $20{,}000$ threshold). After promotion the index returns approximate rather than exact nearest neighbours, but the approximation does not change the routing decision in practice: Stage~3's tier boundaries at $S(q) \geq 0.90$ and $\mathrm{sim}(q, e^{*}) \geq 0.75$ are far from the noise floor of IVF-PQ recall.

Per-query retrieval collapses to a single call: the query embedding $\hat{q}$ is submitted to the index, the top-1 neighbour $e^{*}$ is returned with its cosine similarity and the external integer identifier under which its mutable record is kept in the metadata map. Stage~3 reads $u_{\text{stored}}(e^{*})$ from the metadata map rather than from the index itself; storing composites inside the index is unnecessary because the identifier suffices and keeping quality metadata out of the FAISS payload avoids coupling the routing algorithm to the index's internal layout.

\paragraph{Novelty filter.}
Before an episode is committed at Stage~7 it is screened for novelty: the memory is searched for its single nearest neighbour $e^{*}$, and storage is rejected when
\begin{equation}
\mathrm{sim}(q, e^{*}) \;>\; \tau_{\text{nov}}, \qquad \tau_{\text{nov}} = 0.95,
\label{eq:novelty}
\end{equation}
so that near-duplicates do not accumulate as separate entries. The filter is deliberately one-sided: a query sufficiently similar to an existing entry is assumed to be covered by that entry, and the existing entry's mutable quality record is updated through the retrieval feedback loop rather than through a second store. The threshold is set high ($0.95$) because SBERT similarities above this level are semantically indistinguishable on the Branch-C Cycle-0 calibration set, while thresholds in the $0.85$ to $0.90$ range produced spurious rejections on paraphrases that legitimately tested a different aspect of the same topic. Without the filter, the training-pool curator of Stage~8 would repeatedly train on near-identical (query, answer) pairs, biasing the fine-tuning gradient toward whatever surface form happened to be written many times; the novelty filter is the mechanism that prevents this.

\paragraph{Retrieval feedback update.}
Each accepted retrieval at Stage~4 triggers two updates on the retrieved entry's mutable record. First, the retrieval count is incremented, the running success rate is updated by the standard incremental-mean formula, and (for Tier-1 serves) the hit counter is incremented,
\begin{equation}
r_{\text{new}} = r + 1, \qquad \bar{a}_{\text{new}} = \frac{r \cdot \bar{a} + \mathbb{1}[\text{accepted}]}{r + 1}.
\label{eq:retrieval-stats}
\end{equation}
Second, once an entry has accumulated at least $r_{\min} = 5$ retrievals, its stored composite is nudged by an exponential-moving-average rule with learning rate $\eta = 0.01$,
\begin{equation}
u_{\text{stored}} \;\gets\;
\begin{cases}
u_{\text{stored}} + \eta\,(1 - u_{\text{stored}}) & \text{if the retrieval was accepted,} \\
u_{\text{stored}} - \eta\,u_{\text{stored}} & \text{otherwise,}
\end{cases}
\label{eq:feedback-update}
\end{equation}
which draws the composite toward $1$ on success and toward $0$ on override. The floor of five retrievals before any update fires is important: the feedback rule is fundamentally noisy at low retrieval counts, and the cycle-boundary retroactive re-verification pass of Stage~8 is the authoritative quality update. The feedback rule therefore serves only as a fast-path adjustment between cycles, handling the case where an entry is retrieved heavily before the next cycle boundary arrives.

\paragraph{Memory consolidation.}
At each cycle boundary -- before retroverification, so that consolidation reduces the retroverify workload -- a consolidation pass clusters near-paraphrase episodes and keeps a single representative per cluster. Each stored episode queries its top-$k$ neighbours in the FAISS index (default $k = 20$); pairs with cosine similarity above a tighter consolidation threshold $\tau_{\text{cons}} = 0.92$ (tighter than the novelty threshold of $0.95$ only because consolidation operates on already-stored pairs whose $L_{2}$-normalised embeddings are stable) are joined through union-find into candidate clusters. Two safety guards apply before any cluster is merged: an \emph{answer-consistency} check requires every member's normalised answer string (lowercased, stripped, ASCII-punctuation removed) to equal the representative's, which catches paraphrase-equivalent queries that happen to share SBERT similarity but diverge on the correct answer (``Who directed X'' vs ``Who starred in X''); a \emph{$u_{\text{stored}}$ spread} check rejects clusters whose internal composite range exceeds $0.15$, on the grounds that a high within-cluster variance is evidence that the candidate cluster is actually heterogeneous rather than paraphrase-equivalent. Clusters that pass both guards are merged onto the maximum-$u_{\text{stored}}$ representative (never averaged, ties broken by lower entry identifier for checkpoint reproducibility), with retrieval counts summed, success rates recovered as a retrieval-weighted average, and benchmark tags from the absorbed members stored in the representative's \texttt{merged\_source\_benchmarks} tuple. The SIL training-pool gate of Stage~8 reads this merged set alongside \texttt{source\_benchmark} so that consolidation cannot back-door transfer-benchmark content into the training pool through a representative that happened to be in-distribution. A JSONL audit log records every merge decision -- accepted, skipped for answer inconsistency, or skipped for $u$-spread -- for post-hoc analysis in Chapter~5's memory-dynamics diagnostics.

\paragraph{Popularity and the retroverify queue.}
The hit counter introduced above exists specifically to surface high-traffic Tier-1 memories to the retroverify queue before their composites drift silently. Every Tier-1 serve increments the counter; the counter is reset to zero each time the entry is re-scored at a cycle boundary. At the start of every Stage~8 pass, entries with \texttt{hit\_counter} above a configurable floor (default $50$) are force-scheduled onto the retroverify queue regardless of cycle age. This protects CAEM against the failure mode in which a borderline-composite entry accumulates many Tier-1 serves during a cycle, continues to match future queries (so the retrieval-feedback rule keeps nudging its composite upward on weak evidence), and is never re-challenged by the full verifier because the general retroverify schedule prioritises recency. The popularity queue is the mechanism that makes the monotonicity claim of Theorem~\ref{thm:monotonicity} hold under heavy-tailed retrieval distributions, not just under uniform retrieval.

\paragraph{Capacity management by Value score.}
The store enforces a hard capacity bound $K_{\text{mem}} = 10^{6}$. When utilisation reaches a pruning trigger at $95\%$ of capacity, the bottom $20\%$ of entries by a Value score are removed and the FAISS index is rebuilt to reflect the deletions. The Value score is a convex combination of an Importance and a Recency component,
\begin{equation}
V(e) \;=\; w_{V,I}\,I(e) \;+\; w_{V,R}\,R(e), \qquad w_{V,I} = 0.70,\; w_{V,R} = 0.30,
\label{eq:value}
\end{equation}
with
\begin{align}
I(e) &= w_{I,\phi}\,\phi(e) \;+\; w_{I,r}\,\frac{r(e)}{r_{\max}} \;+\; w_{I,s}\,\bar{a}(e) \;+\; w_{I,u}\,u_{\text{stored}}(e),
\label{eq:importance}\\
R(e) &= \exp\bigl(-\lambda_{R}\,\mathrm{age}(e)\bigr), \qquad \lambda_{R} = 10^{-2},
\label{eq:recency}\\
\phi(e) &= \bigl(\mathrm{cycle}(e) + 1\bigr) / (N_{\text{cycles}} + 1),
\label{eq:phi}
\end{align}
and the Importance weights $(w_{I,\phi}, w_{I,r}, w_{I,s}, w_{I,u}) = (0.40, 0.30, 0.20, 0.10)$. The cycle proxy $\phi(e)$ rewards entries stored by later, better-trained cycles; the normalisation $(N_{\text{cycles}} + 1)$ in the denominator guarantees $\phi \leq 1$ at the final cycle exactly rather than exceeding unity. The retrieval count enters only after normalisation by the global maximum $r_{\max}$ so that a single frequently-hit entry cannot dominate the ranking. The recency term uses real wall-clock age with $\lambda_{R} = 10^{-2}$ per second, a slow decay chosen so that within a single cycle's duration the recency factor changes by at most a few percent; its role is to tie-break between entries of similar importance rather than to dominate them. Pruning is upper-bounded in frequency by the trigger condition, and because every retained entry survives a convex-combination selection over four quality signals, the surviving population is on average better-calibrated than the pre-pruning population.

\paragraph{Persistence.}
The full store -- FAISS index, metadata map, monotonic next-identifier counter, embedding dimensionality, and the \texttt{CAEMConfig} dataclass that parameterises the store -- is serialised to a pair of files at cycle boundaries. The index is written via FAISS's native binary format (\texttt{faiss.write\_index}); the rest is written as a pickle alongside. On reload, the metadata map and next-identifier counter are restored first, then the FAISS index is read; the index's IVF probing parameter is reset from the current configuration rather than from the saved index, so $n_{\text{probe}}$ can be tuned without rebuilding the memory. This is the mechanism by which a long-running experiment can be resumed at an arbitrary cycle boundary under the runner's \verb|--resume_from_cycle| path, and by which a retrieved store from one experimental configuration can be reused to seed another (the Phase-1a pre-main HuggingFace snapshot ships exactly this pair of files).

In summary, the memory module combines a normalised SBERT embedding space, a size-adaptive FAISS backend, a one-sided novelty filter, a per-retrieval feedback rule, a cycle-boundary paraphrase-cluster consolidation pass, a popularity-driven retroverify queue, and a Value-based pruning policy into a single persistent artefact whose query-time interface reduces to a single nearest-neighbour call plus a metadata lookup. All subsequent stages of CAEM are clients of this interface: Stage~3 consumes $(e^{*}, \mathrm{sim}(q, e^{*}), u_{\text{stored}}(e^{*}))$ to route, Stage~7 calls the novelty filter and the add operation to write, and Stage~8 iterates over every stored entry -- with consolidation and the popularity queue shaping which entries are re-verified first -- to keep the composite record aligned with the current generator.

\section{Three-Tier Router}
\label{sec:router}

Stage~3 converts the two upstream quantities -- the pre-route confidence $\upre$ from Stage~1 and the single nearest-neighbour retrieval $(e^{*}, \mathrm{sim}(q, e^{*}), u_{\text{stored}}(e^{*}))$ from Stage~2 -- into a categorical choice of processing tier. The tier choice determines the remainder of the per-query path: Tier~1 returns the stored answer attached to $e^{*}$ without any generation, Tier~2 generates from the Qwen-2.5-3B-Instruct backbone using a one-shot scaffolded chain-of-thought prompt and a forced decoder prefix but without retrieved passages, and Tier~3 generates under retrieval-augmented grounding against a preprocessed Wikipedia passage index. The router is the cost-control mechanism of the pipeline: every decision it makes to dispatch to Tier~1 rather than Tier~2 or~3 saves two to three seconds of decoder time, and every decision it makes to dispatch to Tier~3 rather than Tier~2 incurs the additional cost of passage retrieval and cross-encoder reranking but in exchange supplies the passages that the grounding family of the Stage-5 verifier consumes. The design of the router therefore has to trade cost against safety while keeping the trade explicit.

\paragraph{Two-mechanism design.}
The router applies two decision rules that are deliberately kept separate. The first rule is a safety override evaluated before any score is computed: if $\upre$ falls below a minimum floor $\tau_{\text{safe}} = 0.60$, the query is dispatched to Tier~3 regardless of how well it matched the memory. The second rule is a combined routing score over similarity and stored composite, thresholded to select between Tier~1 and Tier~2. The two rules are not merged into a single formula. The reason is arithmetic: if the safety condition entered additively alongside similarity and $u_{\text{stored}}$, a high stored composite from a very similar episode could arithmetically compensate for a dangerously low $\upre$, sending the query to Tier~1 exactly when the base model's internal state warns that it should not be trusted. Treating model readiness ($\upre$) and memory quality ($u_{\text{stored}}$) as two independent necessary conditions rather than as tradeable summands preserves the semantics that a low $\upre$ is a hard block on the cheap paths, not a discount to be offset.

\paragraph{Safety override.}
The first mechanism, evaluated before any score calculation, is the OR-condition:
\begin{equation}
\upre \;<\; \tau_{\text{safe}} \qquad \Rightarrow \qquad \text{Tier~3},
\label{eq:or-condition}
\end{equation}
with $\tau_{\text{safe}} = 0.60$. Three observations motivate the floor value. First, $\upre$ is temperature-calibrated by Section~\ref{sec:pre-route}'s per-cycle L-BFGS fit from Cycle~1 onward, so the threshold is referenced to a distribution of confidences that has been brought into alignment with empirical exact-match correctness; a floor of $0.60$ corresponds in that calibrated regime to roughly a $40\%$ confident-error risk on the cheap paths, which is unacceptable in an open-domain factual setting. Second, Tier~3 is the only path that introduces external evidence into the verifier, so forcing grounding is the operationally meaningful response to low model readiness: the verifier's grounding family $\{p_{\text{ground}}^{\max}, p_{\text{ground}}^{\text{mean}}, p_{\text{ground}}^{\text{atomic}}\}$ cannot fire on a Tier~1 or Tier~2 answer because no passages are retrieved. Third, the safety override sets a cost floor on queries the system distrusts: a Tier~3 pass takes a few seconds while a Tier~1 hallucination would propagate into memory (through Stages~5 and~7) and potentially into the next cycle's fine-tuning pool. The cost asymmetry makes Tier~3 the unambiguously correct default when the model itself indicates uncertainty.

\paragraph{Combined routing score.}
When the safety override does not fire, the router evaluates a combined score over the similarity and the stored composite of the retrieved episode,
\begin{equation}
S(q) \;=\; \alpha\,\mathrm{sim}(q, e^{*}) \;+\; (1 - \alpha)\,u_{\text{stored}}(e^{*}), \qquad \alpha = 0.70,
\label{eq:routing-score}
\end{equation}
so that similarity carries more weight than the stored composite by a ratio of $7$:$3$. The weighting is deliberate. Similarity is the direct relevance signal: it measures whether the stored episode is about the same question as the query, and no Tier~1 outcome is defensible without high similarity. The stored composite is an orthogonal quality signal that reflects how well the stored answer was verified at storage time and subsequently, but a very high $u_{\text{stored}}$ on a loosely related episode is not an argument for returning that episode's answer; it is an argument for not using Tier~1 at all. Up-weighting similarity therefore encodes the operational priority (relevance first, quality as a secondary gate), while retaining the stored composite in the score means that two episodes at equal similarity are separated by their verification history, with better-verified ones clearing the Tier~1 threshold sooner.

\paragraph{Tier dispatch rules.}
Given the combined score $S(q)$ and the raw similarity $\mathrm{sim}(q, e^{*})$, the router selects a tier according to Table~\ref{tab:tier-thresholds}. Tier~1 requires the combined score to clear $\tau_{\text{tier1}} = 0.90$, which is tight enough that Tier~1 is selected only when both signals are strong: setting $S(q) = 0.90$ and solving for the minimum composite at a given similarity gives $u_{\text{stored}}^{\min} = (0.90 - 0.70\,\mathrm{sim}(q, e^{*})) / 0.30$, so a similarity of $1.00$ requires $u_{\text{stored}} \geq 0.667$, a similarity of $0.95$ requires $u_{\text{stored}} \geq 0.783$, and a similarity of $0.90$ requires $u_{\text{stored}} = 0.900$ exactly. Tier~2 is selected when Tier~1 is not taken and the raw similarity alone exceeds $\tau_{\text{tier2}} = 0.75$; this captures the case where an episode is clearly topically relevant but the composite is not yet strong enough for cheap-path reuse, and one-shot scaffolded-CoT generation from the base model is expected to produce a defensible answer that the verifier can score. Tier~3 is selected in every remaining case: empty memory, low similarity, or borderline combined scores that do not clear the Tier~1 bar without also having the topical similarity to justify Tier~2.

\begin{table}[tbp]
\centering
\caption{Tier dispatch conditions at Stage~3. Conditions are evaluated in order. The safety override supersedes the combined-score rules; when it does not fire, the combined score selects Tier~1 or the raw similarity selects Tier~2, and Tier~3 is the default.}
\label{tab:tier-thresholds}
\small
\begin{tabular}{@{}p{1.4cm}p{7.8cm}p{5.2cm}@{}}
\toprule
\textbf{Tier} & \textbf{Condition (evaluated in order)} & \textbf{Action} \\
\midrule
3 (forced) & $\upre < \tau_{\text{safe}} = 0.60$ (safety override) & Forced RAG; skips score. \\
\addlinespace[2pt]
1          & $S(q) \geq \tau_{\text{tier1}} = 0.90$ and safety override not taken & Return $e^{*}$'s stored answer. \\
\addlinespace[2pt]
2          & $\mathrm{sim}(q, e^{*}) > \tau_{\text{tier2}} = 0.75$, Tier~1 not taken & One-shot scaffolded-CoT generation. \\
\addlinespace[2pt]
3 (default)& otherwise (including empty memory) & DPR retrieve + rerank + RAG. \\
\bottomrule
\end{tabular}
\end{table}

\paragraph{Behaviour on an empty memory.}
The router requires no special case for the cold-start regime in which the episodic store is empty. On an empty memory Stage~2 returns an empty result set, which the router interprets by substituting $\mathrm{sim}(q, e^{*}) = 0$ and $u_{\text{stored}}(e^{*}) = 0$ into Equation~\ref{eq:routing-score}, giving $S(q) = 0$. The combined score then falls below both tier thresholds, the raw similarity fails the Tier~2 bar, and the query dispatches to Tier~3 by default. The safety override rule is unaffected: it depends only on $\upre$, which is produced by Stage~1 regardless of memory state. As a consequence, the entire first cycle of CAEM (Cycle~0, during which the memory is initialised by the cold-start seed of Section~\ref{sec:self-improvement} and the calibration split is collected) runs exclusively on Tiers~2 and~3, which is the intended behaviour: the cheap Tier~1 path comes online organically as $u_{\text{stored}}$ values accumulate through Stage~7 writes and the retrieval-feedback loop of Section~\ref{sec:memory}.

Algorithm~\ref{alg:router} states the Stage~3 logic in its entirety. The strict evaluation order (override first, then score) is load-bearing and is the algorithmic encoding of the two-mechanism separation argued above.

\begin{algorithm}
\caption{Stage~3 three-tier routing decision.}
\label{alg:router}
\begin{algorithmic}[1]
\Require Pre-route confidence $\upre$; retrieved neighbour $(e^{*}, \mathrm{sim}(q, e^{*}), u_{\text{stored}}(e^{*}))$ or empty; weight $\alpha = 0.70$; thresholds $\tau_{\text{safe}} = 0.60$, $\tau_{\text{tier1}} = 0.90$, $\tau_{\text{tier2}} = 0.75$.
\Ensure Tier $t \in \{1, 2, 3\}$ with a safety-override flag.
\If{memory returned no result}
    \State $\mathrm{sim}(q, e^{*}) \gets 0$; $u_{\text{stored}}(e^{*}) \gets 0$ \Comment{empty-memory default}
\EndIf
\If{$\upre < \tau_{\text{safe}}$}
    \State \Return $(t = 3, \; \text{safety\_override} = \text{true})$ \Comment{mechanism 1: OR-condition}
\EndIf
\State $S(q) \gets \alpha\,\mathrm{sim}(q, e^{*}) + (1 - \alpha)\,u_{\text{stored}}(e^{*})$ \Comment{mechanism 2: combined score}
\If{$S(q) \geq \tau_{\text{tier1}}$}
    \State \Return $(t = 1, \; \text{safety\_override} = \text{false})$
\ElsIf{$\mathrm{sim}(q, e^{*}) > \tau_{\text{tier2}}$}
    \State \Return $(t = 2, \; \text{safety\_override} = \text{false})$
\Else
    \State \Return $(t = 3, \; \text{safety\_override} = \text{false})$
\EndIf
\end{algorithmic}
\end{algorithm}

The routing decision is logged on a per-query basis with all five upstream quantities ($\upre$, $\mathrm{sim}(q, e^{*})$, $u_{\text{stored}}(e^{*})$, $S(q)$, and the safety-override flag), because these form the substrate of the routing-distribution analyses reported in Chapter~5. In particular, the fraction of queries that triggered the safety override separately from the fraction that were routed to Tier~3 by the default rule are tracked as distinct metrics: the two populations differ in whether the model \emph{recognised} that it was uncertain (override) versus the memory simply not containing a relevant neighbour (default), and both are reported in the routing-audit table of Chapter~5.

\section{Unified Verifier}
\label{sec:verifier}

Every Tier~2 and Tier~3 answer passes through a single Stage~5 quality gate before the storage decision is made. The verifier consumes a (query, answer) pair and emits nine scalar signals in three families, plus a question--answer relevance signal and a contradiction diagnostic, together with a composite score $u_{\text{stored}}$ and an explicit flag for a confabulation early-exit gate. The signal set is designed around a single organising claim: any single confidence signal is exploitable in one direction or another, so the composite is built from complementary families that fail in different ways. Internal calibration signals are cheap and fine-grained but do not know whether the answer is true; sample-set signals detect disagreement across the model's own resamplings but cannot distinguish a confidently-wrong consensus from a correct one; external grounding signals test factual support against retrieved evidence but are vulnerable to partially-correct answers that are superficially supported; and the Branch~C question--answer relevance signal catches off-topic answers whose entailment to a wrongly-matched passage would otherwise inflate the grounding signals. The composite combines all four so that no single failure mode clears the storage threshold on its own. This section specifies the signals, their aggregation rules, the task-conditional hypothesis construction that makes the grounding signals calibrated across heterogeneous benchmark formats, and the early-exit gate. The four-outcome decision tree that consumes the composite is the subject of Section~\ref{sec:decision}.

\paragraph{Signal families and cost.}
Table~\ref{tab:verifier-signals} summarises the nine core signals plus $q\_a\_\text{relevance}$ and the $p_{\text{contra}}$ diagnostic, together with their source family, the number of model calls each contributes, and whether the signal is literature-derived (L) or a CAEM-specific design decision (D). The per-query cost of Stage~5 is dominated by the decoding passes required for the sample-set family: $M = 3$ chains at temperature $T_s = 0.7$ for $s_{\text{avg}}$ and $p_{\text{entail}}$ (shared), $K = 10$ independent samples at temperature $T_{\text{se}} = 1.0$ for $h_{\text{norm}}$, and $K_{D} = 5$ stochastic passes for $u_{\text{dropout}}$. The external-grounding family adds $20$ passages retrieved from the preprocessed Wikipedia index and reranked to the top $3$ by a cross-encoder, plus one claim-support-judge call per passage. The atomic-fact path additionally performs one greedy decomposition pass and one judge call per atom per passage. The relevance signal $q\_a\_\text{relevance}$ is a single cross-encoder call on the (question, extracted answer) pair, reusing the reranker's cross-encoder backbone so it adds only the forward pass cost, not a model load. All signals are computed once per Stage~5 invocation; there is no adaptive early stopping inside the signal computation itself (the only early exit is the composite-level confabulation gate described below).

\begin{table}[tbp]
\centering
\caption{The nine core verifier signals, the Branch~C relevance signal, and the contradiction diagnostic, grouped by family. Source column: [L] = adapted from literature, [D] = CAEM design. Cost column reports the dominant factor per query at Stage~5.}
\label{tab:verifier-signals}
\small
\begin{tabular}{@{}p{2.4cm}p{5.2cm}p{1.3cm}p{4.4cm}@{}}
\toprule
\textbf{Signal} & \textbf{Family / what it measures} & \textbf{Source} & \textbf{Dominant cost} \\
\midrule
$u_{\text{token}}$ & Internal: lexical confidence in the answer & [L] & 1 teacher-forced forward pass \\
$u_{\text{dropout}}$ & Internal: MC-dropout disagreement across stochastic passes & [L] & $K_{D} = 5$ decoding passes \\
$u_{\text{internal}}$ & Internal: fused scalar $0.5\,u_{\text{token}} + 0.5\,(1 - u_{\text{dropout}})$ & [D] & (reuses above) \\
\addlinespace[2pt]
$s_{\text{avg}}$ & Sample-set: mean pairwise SBERT cosine across $M$ chains & [L] & $M = 3$ decoding passes (shared) \\
$p_{\text{entail}}$ & Sample-set: entailment chain-to-answer, averaged over $M$ chains & [D] & $M$ judge calls (reuses chains) \\
$h_{\text{norm}}$ & Sample-set: normalised semantic entropy over $K$ samples, NLI-clustered & [L] & $K = 10$ decoding passes \\
\addlinespace[2pt]
$p_{\text{ground}}^{\max}$ & Grounding: max task-conditional entailment over top-3 passages & [D] & 3 judge calls \\
$p_{\text{ground}}^{\text{mean}}$ & Grounding: mean task-conditional entailment over top-3 passages & [D] & (reuses above) \\
$p_{\text{ground}}^{\text{atomic}}$ & Grounding: weakest-link entailment over decomposed atomic facts & [D] & $n_{\text{atoms}}$ judge calls + 1 decode \\
\addlinespace[2pt]
$q\_a\_\text{relevance}$ & Relevance: cross-encoder score on (question, extracted answer) & [D] & 1 cross-encoder call \\
\addlinespace[2pt]
$p_{\text{contra}}$ & Diagnostic: max contradiction probability (not read by Stage~6 under MiniCheck) & [D] & 3 judge calls (reuses ensemble) \\
\bottomrule
\end{tabular}
\end{table}

\paragraph{Internal calibration.}
The internal family measures the model's self-assessed confidence in the specific answer under evaluation, using no external information. The token aggregate $u_{\text{token}}$ here is defined analogously to Stage~1's (Equation~\ref{eq:utoken}) but is computed over the full emitted answer rather than over the greedy prefix and is evaluated under teacher forcing with the answer tokens as labels, so it measures the model's probability mass on the exact string being verified:
\begin{equation}
u_{\text{token}}^{\,\text{verify}} \;=\; \exp\!\Bigl(\tfrac{1}{n}\sum_{i=1}^{n} \log p_{\theta}(y_{i} \mid y_{<i}, q)\Bigr),
\label{eq:utoken-verify}
\end{equation}
with non-padding positions only contributing to the mean. The dropout signal follows the predictive-uncertainty framing of Gal and Ghahramani~\cite{pmlr-v48-gal16}: dropout is left active during inference and the generator is sampled $K_{D} = 5$ times at $T = 0.7$; a plurality answer is identified over the $K_{D}$ samples, and
\begin{equation}
u_{\text{dropout}} \;=\; 1 - \frac{\max_{c} \#\{i : a_{i} \equiv a_{c}\}}{K_{D}} \;\in\; [0, 1 - 1/K_{D}],
\label{eq:udropout}
\end{equation}
where $a_{i} \equiv a_{c}$ denotes surface-normalised string equality. Higher values indicate that the model disagrees with itself under perturbation, and the composite uses $(1 - u_{\text{dropout}})$ through the fused internal scalar. The fused internal scalar averages the two,
\begin{equation}
u_{\text{internal}} \;=\; \tfrac{1}{2}\,u_{\text{token}}^{\,\text{verify}} \;+\; \tfrac{1}{2}\,(1 - u_{\text{dropout}}),
\label{eq:uinternal}
\end{equation}
with equal weighting because the two signals capture qualitatively different failure modes: $u_{\text{token}}^{\,\text{verify}}$ is degraded by any token the model considers unlikely even when the overall answer is correct, and $u_{\text{dropout}}$ is degraded only when the model actively disagrees with itself; pairing them tightens the internal signal against either individual noise source. The fused scalar is the signal that enters the composite directly, and is also the signal checked by the confabulation gate.

\paragraph{Sample-set signals.}
The sample-set family examines what the model produces when asked the same query multiple times with non-deterministic sampling. Three signals share this family. For the first two, $M = 3$ chain-of-thought responses are generated at $T = 0.7$ and reused by both signals. The self-consistency signal $s_{\text{avg}}$ is the mean pairwise cosine similarity between chain embeddings produced by the same SBERT model used by the memory (Equation~\ref{eq:cosine}),
\begin{equation}
s_{\text{avg}} \;=\; \binom{M}{2}^{\!-1}\sum_{1 \leq i < j \leq M} \hat{c}_{i}^{\top} \hat{c}_{j},
\label{eq:savg}
\end{equation}
with the unique-pair $i < j$ convention so the diagonal and double counting are not included; this is the formulation of Wang et al.~\cite{wang2023selfconsistency}. The chain-to-answer entailment signal is the mean over the same $M$ chains of the NLI ensemble's entailment probability from chain to answer,
\begin{equation}
p_{\text{entail}} \;=\; \frac{1}{M}\sum_{i=1}^{M} \mathrm{NLI}_{\min}\bigl(c_{i} \Rightarrow a\bigr),
\label{eq:pentail}
\end{equation}
where $\mathrm{NLI}_{\min}$ denotes the minimum entailment probability across all NLI bundles in the ensemble (the conservative aggregation rule described below). The two signals combine differently: $s_{\text{avg}}$ rewards surface and semantic consistency across the chain space, while $p_{\text{entail}}$ rewards chains that individually support the final answer under entailment, so they decorrelate on cases where the model produces several internally-consistent but collectively-wrong chains.

The third sample-set signal is the normalised semantic entropy of Farquhar et al.~\cite{farquhar2024semantic}, computed over an independent set of $K = 10$ samples at $T = 1.0$. Samples are clustered by bidirectional NLI entailment: two samples $s_{i}$ and $s_{j}$ join the same cluster iff
\begin{equation}
\mathrm{NLI}\bigl(s_{i} \Rightarrow s_{j}\bigr) = \mathrm{ENT} \;\wedge\; \mathrm{NLI}\bigl(s_{j} \Rightarrow s_{i}\bigr) = \mathrm{ENT},
\label{eq:semantic-cluster}
\end{equation}
where $\mathrm{ENT}$ is the argmax-label under the NLI ensemble's majority-vote rule. The cluster assignment induces a distribution over cluster sizes, from which the Shannon entropy is computed and normalised by $\log_{2} K$,
\begin{equation}
h_{\text{norm}} \;=\; \frac{-\sum_{c} p_{c}\log_{2} p_{c}}{\log_{2} K}, \qquad p_{c} = \frac{|\{i : s_{i} \in c\}|}{K},
\label{eq:hnorm}
\end{equation}
so that $h_{\text{norm}} \in [0, 1]$. A value near zero indicates that almost all samples express semantically equivalent answers, while a value near one indicates that the samples partition into approximately $K$ distinct semantic classes. The composite treats $(1 - h_{\text{norm}})$ as a confidence: semantic agreement across resamples is a positive signal. When no NLI ensemble is configured, the fallback is surface-level entropy over normalised strings, which is strictly weaker because it cannot merge paraphrases; this fallback matters only in smoke-test configurations.

\paragraph{External grounding and task-conditional hypothesis construction.}
The grounding family anchors the verifier to external evidence rather than to the model's own output distribution. On every Tier~2 and Tier~3 verification, the preprocessed Wikipedia passage retriever is called with the query and asked for the top $K_{\text{retrieve}} = 20$ candidates; these are then reranked by a \texttt{bge-reranker-v2-m3} cross-encoder scored against the concatenated (query, answer) pair, and the top $K_{\text{rerank}} = 3$ are retained. Before any judge call is issued, Branch~C wraps the (passage, answer) pair in a \emph{task-conditional hypothesis} so that the same judge can score heterogeneous benchmark formats on a calibrated $[0,1]$ scale. Three hypothesis constructions apply:
\begin{itemize}
  \item \textbf{Factoid and long-form queries} (TriviaQA, Natural Questions, TruthfulQA, ASQA) use the literal extracted answer as the hypothesis: the judge scores $p \Rightarrow a$ directly.
  \item \textbf{Three-way claims} (FEVER-style SUPPORTS / REFUTES / NEI) use a directional rewriter. For a model claim $a$, a negator module generates $\bar{a}$ via a rule-based primary path (regex patterns for ``X is Y'' / ``X was Y'' / ``X verb(s) Y''), with a Qwen-3B one-shot fallback for rule misses and a self-validation check that rejects the negation when $p_{\text{entail}}(\bar{a} \Rightarrow a) \geq 0.30$. The judge is then called on both directions, and the directional grounding is the max-confident side $\max\{p_{\text{entail}}(p \Rightarrow a), p_{\text{entail}}(p \Rightarrow \bar{a})\}$; an NEI bidirectional certainty $\min\{p_{\text{entail}}(p \Rightarrow a), p_{\text{entail}}(p \Rightarrow \bar{a})\}$ distinguishes principled ``not enough info'' from defensive NEI by measuring how close both directions are to a mutual baseline (low on ambiguity, high on confident refutation).
  \item \textbf{Multi-choice letters} (ARC-Challenge) replace the bare letter with the text of the selected option: if the model answered ``C'', the hypothesis becomes ``The answer to the question is: $\langle$text of option C$\rangle$''. Without substitution the judge would score a bare letter token, which has no propositional content, and the MiniCheck default produced empirically $p_{\text{ground}}^{\text{mean}} \approx 0.08$ on correctly-answered ARC samples; option-text substitution recovers a signal that tracks correctness.
\end{itemize}
Two NLI-based signals are computed on the top-3 passages under the task-conditional hypothesis:
\begin{align}
p_{\text{ground}}^{\max} &\;=\; \max_{p \in P_{\text{top3}}} \mathrm{Judge}_{\min}\bigl(p \Rightarrow h(a)\bigr),
\label{eq:pgmax}\\
p_{\text{ground}}^{\text{mean}} &\;=\; \frac{1}{|P_{\text{top3}}|}\sum_{p \in P_{\text{top3}}} \mathrm{Judge}_{\min}\bigl(p \Rightarrow h(a)\bigr),
\label{eq:pgmean}
\end{align}
where $h(a)$ denotes the task-conditional hypothesis described above. A contradiction diagnostic
\begin{equation}
p_{\text{contra}} \;=\; \max_{p \in P_{\text{top3}}} \mathrm{NLI}_{\max}^{\,\text{contra}}\bigl(p \Rightarrow h(a)\bigr),
\label{eq:pcontra}
\end{equation}
is recorded for auditing. Under the shipping MiniCheck backend $p_{\text{contra}}$ is structurally zero because MiniCheck's binary supported / not-supported output does not separate active refutation from neutral silence; under the \texttt{roberta\_nli\_backend} ablation of Chapter~5 the signal is populated by a true three-class NLI ensemble and the same column becomes diagnostically meaningful. The max and mean aggregates have complementary failure modes. $p_{\text{ground}}^{\max}$ is robust to a partially-incorrect passage set because a single strongly-entailing passage suffices, but it is vulnerable to a superficially-matching passage that happens to mention the answer tokens without actually supporting the claim. $p_{\text{ground}}^{\text{mean}}$ is the reverse: less vulnerable to a single false-match passage, more sensitive to the quality of retrieval as a whole. Both enter the composite so that the resulting estimate is positive only when evidence is both present (max is high) and broadly supportive (mean is high).

\paragraph{Atomic-fact decomposition.}
The atomic signal is the one grounding mechanism that is not available through a single answer-level NLI call. It addresses a specific failure mode the max and mean aggregates cannot distinguish: an answer composed of several factual claims of which one is supported by the passages and another is silently fabricated, where the passage-level entailment is high enough on the supported claim to dominate the aggregates. The mechanism proceeds in three steps. First, the base model is prompted with a fixed instruction to decompose the answer into a numbered list of independent factual claims with pronouns resolved,
\begin{equation}
(a_{1},\dots,a_{n_{\text{atoms}}}) \;=\; \mathrm{Decompose}(a),
\label{eq:decompose}
\end{equation}
and the numbered output is parsed into a deduplicated list of claims. Second, each claim is individually tested against each of the top-3 reranked passages via the NLI ensemble's max-entail rule, giving a per-claim grounding score
\begin{equation}
g_{k} \;=\; \max_{p \in P_{\text{top3}}} \mathrm{NLI}_{\min}\bigl(p \Rightarrow a_{k}\bigr), \qquad k = 1, \dots, n_{\text{atoms}}.
\label{eq:atomic-per}
\end{equation}
Third, the atomic signal is defined as the weakest link across claims,
\begin{equation}
p_{\text{ground}}^{\text{atomic}} \;=\; \min_{k} g_{k},
\label{eq:pgatomic}
\end{equation}
so that the addition of a single unsupported claim suppresses the signal. When decomposition fails (parsing error, empty list, or the atomic path is disabled in configuration), the atomic signal falls back to $p_{\text{ground}}^{\text{mean}}$ rather than to zero; this prevents the composite from silently penalising answers whose decomposition simply happened to fail for tokenisation reasons, while ensuring that atomic decomposition adds information when it succeeds.

\paragraph{Question--answer relevance signal.}
The Branch~C relevance signal $q\_a\_\text{relevance}$ addresses an off-topic failure mode identified in the Phase-1a Cycle-0 audit: hallucinated answers that are syntactically consistent with the retrieved passage but semantically unrelated to the question can score high on $p_{\text{entail}}$ and $p_{\text{ground}}^{\max}$ because the retriever has matched on the hallucination's vocabulary rather than on the question's topic. The signal is a cross-encoder score on the (question, extracted answer) pair,
\begin{equation}
q\_a\_\text{relevance} \;=\; \sigma\bigl(\mathrm{CrossEnc}(q, a_{\text{display}})\bigr) \;\in\; [0, 1],
\label{eq:qarel}
\end{equation}
with $a_{\text{display}}$ the answer after CoT extraction (the same string that is scored for EM in evaluation) and the cross-encoder the same \texttt{bge-reranker-v2-m3} backbone used for passage reranking so that a second model load is not required. The signal has a default neutral value of $0.5$ when no cross-encoder is configured, which leaves the composite unchanged in configurations where the relevance cross-encoder is absent. It carries the largest single weight in the composite ($0.20$, Equation~\ref{eq:ustored-verify}) because the Phase-1a audit identified off-topic high-entailment hallucinations as the dominant early-cycle failure mode and because no other signal in the verifier tests question-to-answer topical match directly.

\paragraph{Claim-support judge: adaptive backend and ensemble fallback.}
The verifier's entailment and grounding signals are computed by a claim-support judge selected at pipeline construction via the \texttt{verifier\_backend} configuration setting. Branch~C ships an \texttt{AdaptiveNLIJudge} that routes between two underlying judges based on hypothesis length so that short propositional claims and long generative passages are each handled by the judge trained for that distribution. The default short-hypothesis judge is MiniCheck-Flan-T5-Large~\citep{tang2024minicheck}, a Flan-T5-Large checkpoint fine-tuned on synthetically decomposed language-model-generated claim-support data whose training distribution matches CAEM's runtime distribution on factoid and 3-way claims. MiniCheck is binary (supported versus not supported) rather than three-class NLI; its supported probability is used directly as the entailment signal, and under this backend $p_{\text{contra}}$ is structurally zero (Section~\ref{sec:decision}). The long-hypothesis judge is a FrozenQwenJudge instance of the same Qwen-2.5-3B-Instruct backbone used for generation, prompted with a one-shot claim-support template and queried for the \texttt{yes}/\texttt{no} token probability on $(p, h) \mapsto \text{``does this passage support this claim?''}$; it is selected when the hypothesis exceeds the MiniCheck context budget (512 tokens after wrapping), which is the operating regime for ASQA long-form answers. Because the two judges are trained on different distributions and emit probabilities on different calibration scales, Branch~C fits a Platt-scaling alignment of the Qwen-judge's logit-space score against MiniCheck's on a 500-sample overlap fold drawn from Cycle-0 eval JSONs; the alignment is accepted only when the Pearson correlation in logit space exceeds $0.70$, and runtime then applies the fitted $(a, b)$ to bring Qwen-judge scores into MiniCheck-calibrated units before the signal enters the composite. The generic-NLI backend, retained for the Chapter~5 ablation variant \texttt{roberta\_nli\_backend}, is an ensemble of two pretrained three-class MNLI models (RoBERTa-Large-MNLI~\cite{liu2019roberta} and DeBERTa-v3-Large-MNLI); under that backend, entailment probabilities are aggregated by minimum across the ensemble so that both models must agree before the signal is allowed to be high, and contradiction probabilities are aggregated by maximum so that any model flagging a contradiction raises the diagnostic. Under all backends, a single-model fallback is supported in configurations where only one judge bundle is available, in which case the min and max reduce to that bundle's own probability; no fallback signal is returned in the zero-model case, and the verifier's output on those signals defaults to $0.5$ for entailment and $0.0$ for contradiction so that the composite neither rewards nor punishes in the absence of evidence.

\paragraph{Composite score.}
The eight non-redundant signals from the ten (the three internal collapse to $u_{\text{internal}}$; the three sample-set contribute separately; all three grounding contribute separately because the max / mean / atomic aggregates have complementary failure modes; the relevance signal contributes separately; $p_{\text{contra}}$ is reserved for the decision tree and does not enter the composite) are combined into the composite $u_{\text{stored}}$ of Equation~\ref{eq:ustored}, reproduced here in the form the verifier actually evaluates:
\begin{equation}
\begin{aligned}
u_{\text{stored}} \;=\;\;&
 w_{\text{g,mean}}\,p_{\text{ground}}^{\text{mean}}
\;+\; w_{\text{g,max}}\,p_{\text{ground}}^{\max}
\;+\; w_{\text{g,atom}}\,p_{\text{ground}}^{\text{atomic}} \\ &
\;+\; w_{\text{entail}}\,p_{\text{entail}}
\;+\; w_{\text{rel}}\,q\_a\_\text{relevance}
\;+\; w_{\text{sc}}\,s_{\text{avg}}
\;+\; w_{\text{int}}\,u_{\text{internal}}
\;+\; w_{\text{se}}\,(1 - h_{\text{norm}}),
\end{aligned}
\label{eq:ustored-verify}
\end{equation}
with weights $(w_{\text{g,mean}}, w_{\text{g,max}}, w_{\text{g,atom}}, w_{\text{entail}}, w_{\text{rel}}, w_{\text{sc}}, w_{\text{int}}, w_{\text{se}}) = (0.22, 0.18, 0.06, 0.14, 0.20, 0.10, 0.08, 0.02)$ summing to unity so the composite remains in $[0, 1]$. The weight allocation is a deliberate reflection of the signal-family design, updated in Branch~C after the Phase-1a Cycle-0 audit identified off-topic hallucination and fragmented-atomic failure modes as dominant early-cycle failure sources. The three grounding terms jointly receive the largest block of mass ($0.46$) because external-evidence entailment is the one lane that can falsify a confidently-wrong but internally-consistent answer; within that block $p_{\text{ground}}^{\text{mean}}$ carries the largest share ($0.22$) for its broad test of factual correctness, $p_{\text{ground}}^{\max}$ carries $0.18$ (added explicitly in Branch~C to reward the robustness-to-partial-retrieval regime), and $p_{\text{ground}}^{\text{atomic}}$ carries a conservative floor of $0.06$ after the Cycle-0 audit showed its prior weight of $0.14$ was dominated by decomposition-failure noise on verbose reasoning chains rather than by true partial-fabrication signal. The relevance signal $q\_a\_\text{relevance}$ carries $0.20$ -- the second-largest single weight -- because it is the only signal that distinguishes ``the passage entails the answer'' from ``the answer is on-topic for the question'', and the audit showed off-topic high-entailment to be the dominant early-cycle failure mode. The remaining weight mass is divided among $p_{\text{entail}}$ ($0.14$), $s_{\text{avg}}$ ($0.10$), $u_{\text{internal}}$ ($0.08$), and $(1 - h_{\text{norm}})$ ($0.02$); the small semantic-entropy weight reflects its noisiness at the $K = 10$ sample budget where cluster-count artefacts dominate. The contradiction probability $p_{\text{contra}}$ does not enter the composite at all -- under MiniCheck it is structurally zero, and the decision tree of Section~\ref{sec:decision} does not read it.

\paragraph{Confabulation early-exit gate.}
Before the composite is formed, the verifier evaluates one hard safety check:
\begin{equation}
u_{\text{internal}} \;\geq\; \tau_{\text{ee,int}} \;\; \wedge \;\; p_{\text{ground}}^{\max} \;\leq\; \tau_{\text{ee,g}} \qquad \Rightarrow \qquad \text{early exit, \textsc{Discard}},
\label{eq:early-exit}
\end{equation}
with $\tau_{\text{ee,int}} = 0.70$ and $\tau_{\text{ee,g}} = 0.20$. The rule targets the specific failure mode where the base model is highly self-confident in an answer that has essentially no external grounding: not only does no retrieved passage broadly support the answer (this is what $p_{\text{ground}}^{\text{mean}}$ would catch), but not even the single most-supportive passage in the top three reaches a low entailment threshold of $0.20$. When both conditions hold, the composite is not even computed, the outcome is set to \textsc{Discard} with the early-exit flag set, and all nine signals are still recorded for analysis. The gate is conservative: it fires only when internal confidence is high and grounding is effectively absent, which is exactly the region of the signal space where the composite could otherwise be raised by the sample-set signals into a \textsc{Store}-range value on an answer the ensemble of NLI models cannot find support for. The precise thresholds were selected so that the gate fires on a manageable minority of Tier~2 and Tier~3 outputs at Cycle~0 and is tracked across cycles as a diagnostic in Chapter~5: a cycle-over-cycle rise would indicate that fine-tuning is introducing confabulations rather than resolving them.

Algorithm~\ref{alg:verifier} specifies the complete verifier, terminating with the composite $u_{\text{stored}}$ and the early-exit flag; the four-outcome decision tree that consumes these values is described in Section~\ref{sec:decision}.

\begin{algorithm}
\caption{Stage~5 unified verifier.}
\label{alg:verifier}
\begin{algorithmic}[1]
\Require Query $q$; answer $a$; base model $f_{\theta}$; SBERT encoder; claim-support judge (\texttt{AdaptiveNLIJudge}); cross-encoder reranker and relevance scorer; passage retriever; task-conditional hypothesis rewriter $h$; composite weights as in Equation~\ref{eq:ustored-verify}.
\Ensure Nine core signals, $q\_a\_\text{relevance}$, $p_{\text{contra}}$ diagnostic, composite $u_{\text{stored}}$, early-exit flag.
\State $u_{\text{token}}^{\,\text{verify}} \gets$ teacher-forced geometric mean log-prob of $a$ under $f_{\theta}(q)$
\State $\{a_{i}\}_{i=1}^{K_{D}} \gets$ generate $K_{D} = 5$ samples with dropout active at $T = 0.7$
\State $u_{\text{dropout}} \gets 1 - \max_{c} \#\{i : a_{i} \equiv a_{c}\} / K_{D}$
\State $u_{\text{internal}} \gets \tfrac{1}{2}\,u_{\text{token}}^{\,\text{verify}} + \tfrac{1}{2}\,(1 - u_{\text{dropout}})$
\State $\{c_{i}\}_{i=1}^{M} \gets$ generate $M = 3$ chains at $T = 0.7$
\State $s_{\text{avg}} \gets \binom{M}{2}^{-1}\sum_{i < j} \hat{c}_{i}^{\top}\hat{c}_{j}$
\State $p_{\text{entail}} \gets \tfrac{1}{M}\sum_{i} \mathrm{Judge}_{\min}(c_{i} \Rightarrow a)$
\State $\{s_{j}\}_{j=1}^{K} \gets$ generate $K = 10$ samples at $T = 1.0$
\State cluster $\{s_{j}\}$ by bidirectional entailment; $h_{\text{norm}} \gets -\sum_{c} p_{c}\log_{2} p_{c} \,/\, \log_{2} K$
\State $P_{\text{top20}} \gets \mathrm{Retrieve}(q, K_{\text{retrieve}} = 20)$; \; $P_{\text{top3}} \gets \mathrm{Rerank}(q, a, P_{\text{top20}}, K_{\text{rerank}} = 3)$
\State $h(a) \gets \mathrm{TaskConditionalHypothesis}(q, a)$ \Comment{directional / option-text / literal}
\State $p_{\text{ground}}^{\max}, p_{\text{ground}}^{\text{mean}} \gets$ max and mean of $\mathrm{Judge}_{\min}(p \Rightarrow h(a))$ over $P_{\text{top3}}$
\State $p_{\text{contra}} \gets \max_{p \in P_{\text{top3}}} \mathrm{NLI}_{\max}^{\,\text{contra}}(p \Rightarrow h(a))$ \Comment{diagnostic; not read by Stage~6}
\State $(a_{1},\dots,a_{n_{\text{atoms}}}) \gets \mathrm{Decompose}(a)$ \Comment{fallback to $p_{\text{ground}}^{\text{mean}}$ on failure}
\State $p_{\text{ground}}^{\text{atomic}} \gets \min_{k} \max_{p \in P_{\text{top3}}} \mathrm{Judge}_{\min}(p \Rightarrow a_{k})$
\State $q\_a\_\text{relevance} \gets \sigma\bigl(\mathrm{CrossEnc}(q, a_{\text{display}})\bigr)$
\If{$u_{\text{internal}} \geq \tau_{\text{ee,int}}$ \textbf{and} $p_{\text{ground}}^{\max} \leq \tau_{\text{ee,g}}$}
    \State \Return all signals, $u_{\text{stored}} \gets 0$, early-exit = true \Comment{confabulation gate}
\EndIf
\State $u_{\text{stored}} \gets$ Equation~\ref{eq:ustored-verify}
\State \Return all signals, $u_{\text{stored}}$, early-exit = false
\end{algorithmic}
\end{algorithm}

The verifier is the most expensive per-query computation in CAEM (several decoder invocations, claim-support-judge calls over up to four premise--hypothesis pairs, one cross-encoder rerank, one cross-encoder relevance call, and a bounded number of additional judge calls in the atomic path), which is why only Tier~2 and Tier~3 queries invoke it. Tier~1 queries are served from memory under the invariant that every stored entry cleared this verifier at its original storage event and has been re-verified at every cycle boundary since (Section~\ref{sec:retroverify}). The composite $u_{\text{stored}}$ and the max grounding probability $p_{\text{ground}}^{\max}$ are the two quantities consumed by the Stage~6 decision tree; the remaining signals, the contradiction diagnostic $p_{\text{contra}}$, and the early-exit flag are recorded on the episode for downstream calibration, ablation, and retroactive re-verification.

\section{Four-Outcome Storage Decision}
\label{sec:decision}

Stage~6 consumes two scalars from the verifier -- the composite $u_{\text{stored}}$ and the maximum grounding probability $p_{\text{ground}}^{\max}$ -- and emits one of four categorical outcomes: \textsc{Store}, \textsc{Deferred}, \textsc{Abstain}, or \textsc{Discard}. The contradiction probability $p_{\text{contra}}$ is also computed at Stage~5 but is recorded only as a diagnostic, because the shipped MiniCheck judge is a binary supported / not-supported classifier and returns $p_{\text{contra}} = 0$ by construction on every query; under the \texttt{roberta\_nli\_backend} ablation of Chapter~5 a three-class NLI ensemble does emit a calibrated contradiction probability, and that diagnostic surfaces in the per-sample record for post-hoc analysis. In the shipped pipeline, Stage~6 reads only $u_{\text{stored}}$ and $p_{\text{ground}}^{\max}$. Stage~7 is the mechanical consumer of the decision: it writes the episode to memory, withholds it from the training pool, returns a principled refusal, or suppresses the generation entirely. The two stages are treated together in this section because they form a single policy surface: the decision tree is meaningful only through the actions it triggers.

Most verification pipelines collapse Stage~6 into a binary store-or-discard rule on a thresholded composite. CAEM retains four outcomes because the composite alone is not sufficient to distinguish between confidently-wrong answers that must be suppressed, uncertain answers that the system should refuse rather than answer, borderline answers that merit another look at the next cycle boundary, and high-quality answers that are safe to commit. Collapsing the middle cases into either \textsc{Store} or \textsc{Discard} either pollutes the Tier~1 path with marginal answers or discards recoverable signal that the retroactive re-verification pass (Section~\ref{sec:retroverify}) would otherwise promote. The four-outcome structure is therefore a recoverability argument: Stage~6 allows a third state (\textsc{Deferred}) in which an answer's quality is good enough to keep around but not good enough to act on, and a fourth state (\textsc{Abstain}) in which the system owns up to uncertainty rather than silently storing or silently discarding.

\paragraph{On the semantics of the cutoffs.}
The three cutoffs $\tau_{\text{store}} = 0.65$, $\tau_{\text{defer}} = 0.45$, and $\tau_{\text{abs}} = 0.20$ below are heuristic bands on the \emph{unitless} composite $u_{\text{stored}}$ of Equation~\ref{eq:ustored}, not probabilities of correctness. A value of $0.65$ is therefore \emph{not} a claim that the answer has a $65\%$ chance of being right: the composite is a weighted sum of eight bounded signals (three grounding, chain-to-answer entailment, question--answer relevance, self-consistency, internal calibration, and the complement of normalised semantic entropy) whose individual scales are not probability-calibrated, and only the token-probability sub-signal passes through the temperature scaler. What is calibrated in the probability sense is $\upre$ at the pre-route stage (Equation~\ref{eq:temp-scale}); the post-generation composite here is a policy-defining aggregate, chosen so that the induced decision bands yield acceptable \textsc{Store} / \textsc{Deferred} / \textsc{Abstain} / \textsc{Discard} rates on the held-out calibration slice. The three cutoffs are moreover re-fit at every cycle boundary by an EMA-smoothed quantile rule (Section~\ref{sec:self-improvement}), so the operating bands track the generator's shifting composite distribution as fine-tuning progresses. Numeric values quoted in this section are therefore Cycle-0 fits, not invariants; the cycle-specific values used at inference time are logged per cycle in \texttt{calibrated\_thresholds.json}, and the Chapter~5 ablations sweep these cutoffs to confirm the policy is not brittle, but the numeric values themselves should be read as operating points on a dimensionless scale rather than as probabilistic guarantees.

\paragraph{Store and Deferred thresholds.}
The composite $u_{\text{stored}}$ is thresholded against two cutoffs in a cascade:
\begin{equation}
\textsc{Store} \text{ if } u_{\text{stored}} \geq \tau_{\text{store}}, \quad
\textsc{Deferred} \text{ if } \tau_{\text{defer}} \leq u_{\text{stored}} < \tau_{\text{store}},
\label{eq:store-defer}
\end{equation}
with $\tau_{\text{store}} = 0.65$ and $\tau_{\text{defer}} = 0.45$. The two thresholds partition the composite range into three bands. The upper band ($u_{\text{stored}} \geq 0.65$) authorises a commit: the composite is high enough that the answer is treated as verified knowledge, with all the downstream consequences (memory write, eligibility for the training pool, availability for Tier~1 retrieval on future queries). The middle band ($0.45 \leq u_{\text{stored}} < 0.65$) authorises a hold: the answer is plausibly correct (it cleared the $0.45$ floor, which is already well above a confidently-wrong answer's composite) but has not cleared the commit bar. The lower band falls through to the \textsc{Abstain} and \textsc{Discard} checks below. The gap between the two thresholds, and its midpoint value of $0.55$, is the single tuning knob with the most influence on the \textsc{Store}-to-\textsc{Deferred} ratio across the evaluation corpus; a tighter gap would collapse \textsc{Deferred} into \textsc{Store} at the risk of training on under-verified material, and a wider gap would shift recoverable answers into the \textsc{Discard} state and lose them permanently.

\paragraph{Abstain rule.}
Below $\tau_{\text{defer}}$, the decision tree evaluates a grounding-floor check before falling through to \textsc{Discard}:
\begin{equation}
u_{\text{stored}} < \tau_{\text{defer}} \;\wedge\; p_{\text{ground}}^{\max} < \tau_{\text{abs}} \qquad \Rightarrow \qquad \textsc{Abstain}, \qquad \tau_{\text{abs}} = 0.20.
\label{eq:abstain}
\end{equation}
The rule identifies the specific regime in which the composite is low and the single most-supportive retrieved passage does not meaningfully support the answer: there is neither internal confidence that the answer is right nor external evidence that it is right. In this regime, \textsc{Abstain} is preferred over \textsc{Discard} because the two outcomes differ in what the pipeline returns to the user. A \textsc{Discard} suppresses the generated answer and emits nothing downstream-usable; an \textsc{Abstain} returns a principled refusal (an explicit ``I do not know'' response), which is a useful output in factual QA settings where the failure mode of highest cost is a confident wrong answer rather than a refusal. The grounding-floor threshold $\tau_{\text{abs}} = 0.20$ is intentionally low: it fires only when the top-3 passages collectively fail to produce even a weakly-supporting entailment, which is the signal that no recoverable evidence exists in the current retrieval neighbourhood.

\paragraph{Discard as the default.}
Every case that does not match the preceding rules falls through to \textsc{Discard}. The default fires in exactly one scenario under the shipped pipeline: the intermediate regime where the composite is low ($u_{\text{stored}} < \tau_{\text{defer}}$) but retrieval did produce at least one supporting passage ($p_{\text{ground}}^{\max} \geq \tau_{\text{abs}}$). The answer is not good enough to keep, but the presence of some grounding signal means the failure was not of the recoverable-refusal kind (the query did have evidence in the retrieval index; the model simply failed to use it). Reporting a \textsc{Discard} rather than an \textsc{Abstain} in this case is the decision tree's way of distinguishing ``we had evidence and still got it wrong'' from ``we had no evidence at all'', because the two motivate different downstream responses in Chapter~5's error analysis: the first is evidence of a generation failure, the second of a retrieval failure. A separate \textsc{Discard} path is also taken when the Stage~5 confabulation early-exit gate fires (Equation~\ref{eq:early-exit}); in that case the composite is not even computed, the outcome is set to \textsc{Discard} with the early-exit flag, and the nine signals are still recorded on the sample for post-hoc analysis.

\begin{table}[tbp]
\centering
\caption{The four-outcome decision tree at Stage~6 and the Stage~7 action it triggers. Conditions are evaluated in order; the first matching condition determines the outcome.}
\label{tab:decision-tree}
\small
\begin{tabular}{@{}p{1.8cm}p{6.2cm}p{5.6cm}@{}}
\toprule
\textbf{Outcome} & \textbf{Condition (in order)} & \textbf{Stage~7 action} \\
\midrule
\textsc{Store}        & $u_{\text{stored}} \geq \tau_{\text{store}} = 0.65$ & Write episode; eligible for training if $u_{\text{stored}} \geq \tau_{\text{train}} = 0.75$ and the episode's source benchmark is in the training pool. \\
\addlinespace[2pt]
\textsc{Deferred}     & $\tau_{\text{defer}} \leq u_{\text{stored}} < \tau_{\text{store}}$, i.e.\ $0.45 \leq u_{\text{stored}} < 0.65$ & Pushed to the deferred-entry buffer (Section~\ref{sec:deferred-reconsider}) for cycle-boundary reconsideration; not committed to memory and not used for training at the current cycle. \\
\addlinespace[2pt]
\textsc{Abstain}      & $u_{\text{stored}} < \tau_{\text{defer}}$ \textbf{and} $p_{\text{ground}}^{\max} < \tau_{\text{abs}} = 0.20$ & Return principled refusal to user; no memory write. \\
\addlinespace[2pt]
\textsc{Discard} (default) & otherwise (including confabulation-gate early exit) & Suppress output; no memory write. \\
\bottomrule
\end{tabular}
\end{table}

\paragraph{Stage~7 consequences and training eligibility.}
Stage~7 is deterministic given the decision: \textsc{Store} writes a complete \textsc{EpisodicEntry} (Table~\ref{tab:entry-schema}) with the full nine-signal record, the Branch~C relevance signal $q\_a\_\text{relevance}$, the contradiction diagnostic $p_{\text{contra}}$, the source-benchmark tag, the composite, the early-exit flag, and the decision string itself; \textsc{Deferred} retains the episode for re-consideration at the next cycle boundary without committing it as a first-class memory entry; \textsc{Abstain} returns a refusal to the user; \textsc{Discard} suppresses the generation. Three additional invariants apply at write time. First, a \textsc{Store} decision is subject to the novelty filter of Equation~\ref{eq:novelty}: if the query is too similar to an existing memory entry, the store is rejected at Stage~7 even though Stage~6 green-lit it, and the existing entry's retrieval feedback fields are updated instead. Second, a Branch~C storage sanitiser rejects the write when the extracted answer matches a template-leak pattern (placeholder text that leaked from the prompt scaffold) or an evasive pattern (phrases such as ``the context does not mention''); these are explicit failure modes the Phase-1a Cycle-0 audit identified as early-cycle memory poisoners, and the sanitiser is the last line of defence before the episode enters the FAISS index. Third, an entry is flagged as eligible for the fine-tuning pool only if its composite further clears a training-eligibility threshold $\tau_{\text{train}} = 0.75$ \emph{and} its source benchmark (or any tag absorbed via consolidation; Section~\ref{sec:memory}) is in the training pool rather than the transfer panel; this is a stricter bar than the \textsc{Store} threshold because the training pool must be composed of only the most well-verified material, and the ID/OOD gate prevents transfer-benchmark content from leaking into training through memory. The gap between $\tau_{\text{store}} = 0.65$ and $\tau_{\text{train}} = 0.75$ is therefore an architectural commitment: Tier~1 retrieval and training data curation have different quality bars, and conflating them would either starve the training pool or over-serve Tier~1 from marginal memory.

Algorithm~\ref{alg:decision} states the Stage~6 and Stage~7 logic in the evaluation order the code uses. The order is load-bearing: the composite threshold is tested strictly before the \textsc{Abstain} check so that a borderline-high composite cannot be deflected into an \textsc{Abstain} by incidentally low grounding; the \textsc{Abstain} check precedes the default \textsc{Discard} so that the refusal path is taken whenever grounding is absent rather than merely weak; the novelty filter and storage sanitiser are applied only after the Stage~6 green-light so that a rejected novelty-check does not silently mutate Stage~6's recorded decision; and the training-eligibility check is applied only after the memory write so that a deferred entry does not accidentally enter the pool.

\begin{algorithm}
\caption{Stage~6 decision tree and Stage~7 action.}
\label{alg:decision}
\begin{algorithmic}[1]
\Require Composite $u_{\text{stored}}$; max grounding $p_{\text{ground}}^{\max}$; extracted answer $a_{\text{display}}$; source benchmark $b$; thresholds $\tau_{\text{store}} = 0.65$, $\tau_{\text{defer}} = 0.45$, $\tau_{\text{abs}} = 0.20$, $\tau_{\text{train}} = 0.75$ (Cycle-0 fits; re-fit per cycle).
\Ensure Outcome $\in \{\textsc{Store}, \textsc{Deferred}, \textsc{Abstain}, \textsc{Discard}\}$; Stage~7 action.
\If{$u_{\text{stored}} \geq \tau_{\text{store}}$}
    \If{novelty filter passes (Eq.~\ref{eq:novelty}) \textbf{and} $a_{\text{display}}$ is not a template-leak or evasive pattern}
        \State write episode to memory with source\_benchmark $= b$
        \If{$u_{\text{stored}} \geq \tau_{\text{train}}$ \textbf{and} $b \in$ training pool} flag as training-eligible \EndIf
    \Else
        \State update retrieval-feedback of existing neighbour (novelty reject) or skip write (sanitiser reject)
    \EndIf
    \State \Return $\textsc{Store}$
\EndIf
\If{$u_{\text{stored}} \geq \tau_{\text{defer}}$}
    \State push to deferred-entry buffer (Section~\ref{sec:deferred-reconsider}); do not commit to memory
    \State \Return $\textsc{Deferred}$
\EndIf
\If{$p_{\text{ground}}^{\max} < \tau_{\text{abs}}$}
    \State return a principled refusal to the caller
    \State \Return $\textsc{Abstain}$
\EndIf
\State suppress output
\State \Return $\textsc{Discard}$
\end{algorithmic}
\end{algorithm}

The per-cycle distribution over the four outcomes is one of the principal diagnostics reported in Chapter~5, because it reveals how the verifier is carving up the generation stream and how this carving evolves across cycles: a healthy trajectory shifts mass from \textsc{Deferred} into \textsc{Store} as the base model fine-tunes on the training pool (answers the model was borderline about in cycle~$n$ become commits in cycle~$n{+}1$), while a pathological trajectory shifts mass from \textsc{Store} into \textsc{Discard} (the model is drifting into confabulation and the early-exit gate is increasingly firing). The categorical structure of the Stage~6 decision is therefore not only a quality-control mechanism at the per-query level but a diagnostic at the cycle level, and Chapter~5 reports the full (cycle $\times$ outcome) distribution as a primary figure.

\section{Retroactive Re-Verification}
\label{sec:retroverify}

Each self-improvement cycle ends with a single sweep over the entire memory in which every stored episode is re-scored under the fine-tuned model and either promoted, pruned, or left untouched. This is Stage~8 in the pipeline and runs once per cycle boundary after fine-tuning (and is therefore distinct from the per-query retrieval-feedback update of Section~\ref{sec:memory}, which writes to $u_{\text{stored}}$ continuously as Tier~1 answers are consumed). The retroactive pass is the only place in the architecture where the memory's stored quality estimates are reconciled with the model that now consumes them: without it, a stored composite computed by the Cycle-0 verifier would continue to drive routing for episodes that are now more or less trustworthy under the Cycle-1 model.

\paragraph{Motivation: stored composites are snapshots, not invariants.}
Every $u_{\text{stored}}$ value written to memory is a snapshot of the verifier's view of a (question, answer) pair at the moment the episode was stored. Three of the nine signals (\,$u_{\text{token}}$, $u_{\text{dropout}}$, $u_{\text{internal}}$\,) read the model's internal state directly and therefore drift whenever the weights drift; two more (\,$s_{\text{avg}}$, $h_{\text{norm}}$\,) depend on stochastic decoding and drift with the model's output distribution. The remaining grounding signals depend on retrieval against a fixed passage corpus and are stable, but the NLI entailment scores between those passages and the generated answer still shift if the answer text itself is regenerated under the new model. Taken together, the stored composite drifts under every fine-tune. Retroactive re-verification is the pass that makes the drift visible and acts on it.

\paragraph{Loop-filter pre-prune.}
Before the verifier is invoked on any stored episode, Branch~C runs a cheap syntactic check on the episode's \texttt{reasoning\_chain}: a loop filter (\texttt{caem.training.loop\_filter.is\_repetitive\_loop}) scores the chain on a distinct-4/compression-ratio criterion and flags chains that collapse into a short repeating span or fail to produce four distinct 4-grams in their middle span. Entries whose stored chain trips the filter are pruned from both the FAISS index and the metadata map without a verifier forward pass, on the grounds that no amount of re-scoring rehabilitates a reasoning chain that is syntactically degenerate; the composite computed over such a chain is either spuriously low (the generator has collapsed and the verifier has nothing coherent to score) or spuriously high (pathological repetition can sometimes entail itself under NLI backbones). The loop filter is therefore a pure cost-saver and a quality floor in one: compute is reserved for episodes whose content is at least well-formed enough to evaluate. Loop-pruned counts are tracked separately from threshold-pruned counts in the per-cycle diagnostic log (\texttt{loop\_pruned} vs \texttt{threshold\_pruned}), and the two are summed into the single aggregate $n_{\text{rm}}$ that Chapter~5 reports, so that the aggregate retains its API shape while the breakdown is available for post-hoc analysis.

\paragraph{Three-way branch.}
For each stored episode $e$ that survives the loop-filter pre-prune, with current composite $u_{\text{stored}}^{(t)}$, the pass computes a new composite $u_{\text{stored}}^{\text{new}}$ using the updated verifier and assigns the episode to one of three branches:
\begin{equation}
e \;\longmapsto\;
\begin{cases}
\text{prune} & \text{if } u_{\text{stored}}^{\text{new}} \;<\; \tau_{\text{retro}}, \\[3pt]
\text{update} & \text{if } u_{\text{stored}}^{\text{new}} \;\neq\; u_{\text{stored}}^{(t)} \text{ and } u_{\text{stored}}^{\text{new}} \geq \tau_{\text{retro}}, \\[3pt]
\text{keep} & \text{otherwise},
\end{cases}
\qquad \tau_{\text{retro}} = 0.50.
\label{eq:retro-branch}
\end{equation}
A pruned episode is hard-deleted from both the FAISS index and the metadata dictionary. An updated episode has its composite replaced with $u_{\text{stored}}^{\text{new}}$ and its entire signal record -- the nine core signals, $p_{\text{contra}}$, the Branch~C relevance signal $q\_a\_\text{relevance}$, the \textsc{Store}/\textsc{Deferred}/\textsc{Abstain}/\textsc{Discard} decision string, and the early-exit flag -- overwritten with the new verifier output; its popularity counter \texttt{hit\_counter} is reset to zero (Section~\ref{sec:memory}) to mark the post-retroverify baseline for the next cycle's popular-needs-recheck queue, and its \texttt{retroverified} flag is set to \textsc{true}. A kept episode retains its stored composite unchanged, has its \texttt{hit\_counter} reset to zero, and has its \texttt{retroverified} flag set to \textsc{true} to record that the pass did see it this cycle.

\paragraph{Bidirectional update rationale.}
The update branch is deliberately bidirectional: if the new composite exceeds the prune threshold, the stored value is replaced regardless of whether $u_{\text{stored}}^{\text{new}}$ is higher or lower than $u_{\text{stored}}^{(t)}$. An earlier design (Session 42) made retroverify strictly upward-only on the argument that noisy single-pass verifier variance could demote a stable, high-quality episode below the Tier~1 retrieval band on a coin flip. The Branch~C Cycle-0 audit (2026-04-22) falsified that worry in the other direction: the upward-only rule was leaving stale confident entries in memory, because entries whose updated score fell in the $[\tau_{\text{retro}}, u_{\text{stored}}^{(t)})$ band were being silently preserved at their original composite rather than downgraded to the score the current verifier actually produced. The fix -- gated behind \texttt{retroverify\_allow\_downgrade}, which defaults to \textsc{true} in the shipped configuration -- is to write whichever value the current verifier returns, clipped only at the categorical prune floor $\tau_{\text{retro}}$ so that entries the verifier no longer supports are removed rather than quietly demoted to a borderline stored composite. The stochastic-variance concern is absorbed by the popular-needs-recheck queue below: entries that Tier~1 is actually serving heavily (and whose Tier-1 serves the retrieval-feedback rule is already updating) are force-scheduled for re-verification before the regular sweep, so a single unlucky score is rapidly re-tested under the next cycle's weights rather than persisting as an artefact. The threshold $\tau_{\text{retro}} = 0.50$ sits strictly above the original \textsc{Deferred} threshold $\tau_{\text{defer}} = 0.45$, so an entry that survives re-verification maintains a margin above what would merely have deferred it at first commit: re-verification is allowed to be stricter than initial storage because the system had already chosen to keep the episode once and should not pay full storage cost for an entry that re-scoring under current weights would not even have deferred.

\paragraph{Full-record refresh on update.}
When an episode is updated, the replacement is not limited to the scalar composite. The full verifier record is overwritten: all nine core signals, $p_{\text{contra}}$, $q\_a\_\text{relevance}$, the decision string, and the early-exit flag. Downstream consumers of the memory (the calibration diagnostics of Chapter~5, the atomic-fact-distribution analyses, any future ablation that needs to re-partition the memory by verifier decision) therefore always see the latest verifier view of an updated episode, not a stale mixture of the storage-time record and a later-computed composite. Kept episodes, in contrast, do not refresh their signal record; they carry their storage-time record unchanged, with only the \texttt{retroverified} flag and \texttt{hit\_counter} mutated. The asymmetry matches the composite's: only an update is a substantive refresh.

\paragraph{Popular-needs-recheck queue.}
Between cycle boundaries, Branch~C exposes a \texttt{force\_retroverify\_queue(hit\_threshold)} method on the memory store that returns every entry whose \texttt{hit\_counter} exceeds a configurable floor (default \texttt{hit\_counter\_force\_retroverify} = $10$), sorted by descending hit count with ties broken by entry identifier for deterministic iteration. The queue is consumed by two callers: the cycle-boundary retroverify scheduler, which prioritises these entries at the start of the regular sweep so that heavy-traffic Tier~1 memories are re-scored before cheaper ones; and an optional between-cycles fast-path script that can force re-verification of a specific high-pressure entry without waiting for the next cycle boundary. The design addresses a compounding failure mode the regular sweep cannot catch on its own: an entry that is Tier~1-served many times in a single cycle accumulates retrieval-feedback nudges on its composite based on usage (Section~\ref{sec:memory}, Equation~\ref{eq:feedback-update}) without having its underlying verifier signals re-examined, so a borderline-composite entry serving many queries can drift upward on Tier~1 acceptances alone while remaining poorly grounded. The popularity queue is the mechanism that makes the monotonicity claim of Theorem~\ref{thm:monotonicity} hold under heavy-tailed retrieval distributions, not just under uniform retrieval; it also establishes the ID-level correspondence between ``heavy user-side impact'' and ``priority model-side re-scrutiny'' that supports the purity and convergence arguments of Section~\ref{sec:theoretical-analysis}.

\paragraph{Retroverify versus retrieval feedback.}
The two memory-update mechanisms have disjoint responsibilities and should not be conflated. Retrieval feedback (Section~\ref{sec:memory}, Equation~\ref{eq:feedback-update}) is a continuous, per-query EMA update that tracks how often an episode's Tier~1 answer was accepted by the user or by downstream checks, and is therefore about \emph{usage quality}. Retroactive re-verification is a discrete, per-cycle pass that recomputes the composite from first principles under the current model weights, and is about \emph{model-relative quality}. The two signals can disagree: an episode can be consistently accepted at Tier~1 (high retrieval-feedback $u_{\text{stored}}$ movement) while simultaneously drifting below the retroverify threshold when the fine-tuned model now produces a different internal-state reading. In such a case retroverify wins, because the system's correctness argument rests on what the \emph{current} verifier says about the (question, answer) pair, not on historical usage alone. The code enforces this precedence by running retroverify after fine-tuning and allowing it to remove entries that retrieval feedback alone would have kept.

\paragraph{Abort handling.}
Retroactive re-verification is skipped on any cycle whose fine-tuning step was rolled back by the forgetting-ratio guard (Section~\ref{sec:self-improvement}). The reason is operational rather than conceptual: if the weights were restored to the previous cycle's checkpoint, re-scoring with unchanged weights produces the same composites that were already stored, up to stochastic-decoding noise, and the pass would therefore pay full compute cost for a no-op update with a small probability of injecting noise through the promotion path. The orchestrator in \texttt{self\_improvement.py} therefore guards the retroverify call with \texttt{if verify\_fn is not None and not aborted}, and records the counters $n_{\text{retroverified}} = n_{\text{retropruned}} = 0$ on aborted cycles so that Chapter~5's retroverify-dynamics figures can distinguish aborted cycles from cycles in which no episode moved.

\begin{algorithm}
\caption{Stage~8 retroactive re-verification (runs once per cycle boundary after successful fine-tuning).}
\label{alg:retroverify}
\begin{algorithmic}[1]
\Require Memory store $\mathcal{M}$; updated verifier $V^{(t+1)}$; loop filter \texttt{is\_repetitive\_loop}; prune threshold $\tau_{\text{retro}} = 0.50$; flags \texttt{retroverify\_prune\_loops} = \textsc{true}, \texttt{retroverify\_allow\_downgrade} = \textsc{true}.
\Ensure Memory $\mathcal{M}$ updated in place; counters $n_{\text{up}}, n_{\text{rm}}$ with $n_{\text{rm}}$ split internally into loop-pruned and threshold-pruned.
\State $n_{\text{up}} \gets 0$; $n_{\text{loop}} \gets 0$; $n_{\text{thr}} \gets 0$
\State snapshot $\gets$ list of entry IDs in $\mathcal{M}$, popular-needs-recheck queue first \Comment{Section~\ref{sec:memory}}
\ForAll{$\text{eid} \in \text{snapshot}$}
    \State $e \gets \mathcal{M}.\text{lookup}(\text{eid})$; \textbf{if} $e = \bot$ \textbf{continue}
    \If{\texttt{retroverify\_prune\_loops} \textbf{and} \texttt{is\_repetitive\_loop}$(e.\text{reasoning\_chain})$}
        \State $\mathcal{M}.\text{remove}(\text{eid})$; $n_{\text{loop}} \gets n_{\text{loop}} + 1$; \textbf{continue} \Comment{no verifier pass}
    \EndIf
    \State \texttt{try} $o \gets V^{(t+1)}.\text{verify}(e.\text{question}, e.\text{answer})$; \texttt{except} \textbf{continue}
    \State $u^{\text{new}} \gets o.u_{\text{stored}}$
    \If{$u^{\text{new}} < \tau_{\text{retro}}$}
        \State $\mathcal{M}.\text{remove}(\text{eid})$; $n_{\text{thr}} \gets n_{\text{thr}} + 1$
    \ElsIf{$u^{\text{new}} \neq e.u_{\text{stored}}$ \textbf{and} (\texttt{retroverify\_allow\_downgrade} \textbf{or} $u^{\text{new}} > e.u_{\text{stored}}$)}
        \State overwrite $e$'s nine signals, $p_{\text{contra}}$, $q\_a\_\text{relevance}$, decision, early-exit flag from $o$
        \State $\mathcal{M}.\text{update}\_u_{\text{stored}}(\text{eid}, u^{\text{new}})$
        \State $e.\text{retroverified} \gets \textsc{true}$; $e.\text{hit\_counter} \gets 0$; $n_{\text{up}} \gets n_{\text{up}} + 1$
    \Else
        \State $e.\text{retroverified} \gets \textsc{true}$; $e.\text{hit\_counter} \gets 0$ \Comment{checked this cycle; no update}
    \EndIf
\EndFor
\State $n_{\text{rm}} \gets n_{\text{loop}} + n_{\text{thr}}$
\State \Return $(n_{\text{up}}, n_{\text{rm}})$ with breakdown $(n_{\text{loop}}, n_{\text{thr}})$ in diagnostic log
\end{algorithmic}
\end{algorithm}

The per-cycle pair $(n_{\text{up}}, n_{\text{rm}})$ is the load-bearing diagnostic of the retroactive pass and is reported directly in Chapter~5. The dynamics claim underlying CAEM rests on the expected pattern that $n_{\text{up}}$ grows and $n_{\text{rm}}$ stays small across cycles: stored answers that the cycle-0 verifier rated marginally are re-rated more confidently by a better-calibrated cycle-1 verifier and by a better-grounded cycle-1 model, which is exactly the compounding signal that distinguishes confidence-aware episodic memory from a static retrieval cache. A trajectory in which $n_{\text{rm}}$ grows instead, especially if concentrated on episodes stored in a specific earlier cycle, is the counter-signal: the memory is accumulating material that the current model no longer regards as verified, and the system is in a regime where retroactive pruning is doing more work than retroactive promotion.

\section{Deferred-Entry Reconsideration}
\label{sec:deferred-reconsider}

Retroactive re-verification operates on episodes that have already been committed to memory; it is the mechanism by which the stored record is kept in agreement with the current verifier. It has no purchase, however, on queries whose Stage~6 outcome was \textsc{Deferred}: a deferred answer is plausibly correct but below the commit bar, and by construction it is not yet in memory for retroverify to see. Without a complementary mechanism, a deferred answer would be silently discarded at the end of Stage~7 and would have no second chance to enter memory even if the fine-tuned cycle-$(c{+}1)$ verifier, evaluating the same (question, answer) pair, would now rate it above $\tau_{\text{store}}$. This section specifies the deferred-entry buffer that holds such answers across cycle boundaries and the reconsideration pass that decides, at each boundary, which of them are promoted to memory, which are discarded as no-longer-plausible, and which are kept for one more cycle. Together with the retroverify sweep of Section~\ref{sec:retroverify}, reconsideration forms the second half of CAEM's cycle-boundary memory-update step: retroverify reconciles committed entries against the fine-tuned weights, reconsideration adjudicates the held-but-uncommitted tail.

\paragraph{Why the buffer exists.}
The thesis's decision-tree design carves the composite axis into four bands precisely to distinguish answers the system is confident in (\textsc{Store}) from answers whose confidence is insufficient to commit \emph{now} but whose quality is not low enough to throw away (\textsc{Deferred}). The value of the \textsc{Deferred} band is entirely conditional on whether those entries get a genuine second look: if the deferral decision is behaviourally identical to a discard, the band collapses to three outcomes and the composite axis is under-utilised. Retroverify alone cannot play this role, because the full-memory sweep of Section~\ref{sec:retroverify} iterates only over committed episodes. The buffer closes this gap by making deferral a temporally extended decision: Stage~7 pushes the (question, answer, embedding, storage-cycle, verifier-record, source-benchmark) tuple into a side structure separate from the main FAISS memory, and the next cycle's reconsideration pass re-scores it under the updated model.

\paragraph{Bounded FIFO with time-to-live.}
The buffer is a bounded first-in-first-out queue with a configured maximum size $N_{\text{buf}}$ (the implementation uses $N_{\text{buf}} = 10{,}000$) and a per-entry time-to-live $T$ measured in cycles (the implementation uses $T = 2$). The FIFO discipline is load-bearing for two reasons. First, deferred output at a rate that exceeds the rate of promotion would otherwise grow the buffer without bound, at which point its re-scoring cost per cycle becomes a denial-of-service against the rest of the cycle-boundary work; the cap makes the per-cycle reconsideration cost $O(N_{\text{buf}})$ in the worst case regardless of how much deferred traffic the verifier has produced. Second, eviction-on-overflow means that the buffer's content is always the most recently deferred material, which is the material whose verifier record reflects weights closest to the current weights and is therefore the material for which reconsideration is most likely to move the composite. Under capacity pressure the buffer is self-cleaning in the direction of relevance, not in the direction of age. Each entry additionally carries an $\mathrm{age}$ counter incremented on every cycle in which reconsideration did not yet promote it, and the TTL bounds this counter: once $\mathrm{age} + 1 \geq T$ the entry is dropped, which prevents a stream of persistently borderline (question, answer) pairs from occupying slots that newer deferrals could use more productively. The cap and the TTL together bound both the memory footprint of the buffer and the temporal horizon over which a deferred answer can be resurrected. Alongside these two discipline fields the entry also carries a \texttt{source\_benchmark} tag captured at push time, so that the training-pool ID/OOD gate in Section~\ref{sec:self-improvement} sees the same origin signal whether an episode entered main memory through the direct Stage~7 \textsc{Store} path or through a later promotion from the deferred buffer.

\paragraph{Three-way rule.}
At each cycle boundary, after retroverify completes (and only if the cycle was not aborted by the forgetting-ratio guard), every entry $b$ in the buffer is re-scored under the updated verifier $V^{(c+1)}$. Writing $u^{\text{new}}$ for the new composite and $o^{\text{new}}$ for the new decision string, each entry is routed by
\begin{equation}
b \;\longmapsto\;
\begin{cases}
\text{promote} & \text{if } o^{\text{new}} = \textsc{Store} \text{ and } u^{\text{new}} \;\geq\; \tau_{\text{store}}, \\[3pt]
\text{drop}    & \text{if } \mathrm{age}(b) + 1 \;\geq\; T, \\[3pt]
\text{keep}    & \text{otherwise (requeue with } \mathrm{age}(b) \;\mathrel{+}= 1\text{)}.
\end{cases}
\qquad T = 2.
\label{eq:defer-branch}
\end{equation}
The promotion branch fires exactly when the new verifier would have chosen \textsc{Store} for the same (question, answer) pair at the current cycle: the rule does not bypass the Stage~6 decision tree but replays it, and the promotion threshold is the same $\tau_{\text{store}} = 0.65$ as at first verification rather than the lower retroverify prune threshold. An entry whose new composite is in the \textsc{Abstain} or \textsc{Discard} bands is not partially rescued; the only path into memory is a full \textsc{Store} decision under the current verifier. Conversely, an entry whose new composite is still in the deferred band is not thrown away on that basis alone; only the TTL rule can drop it. The three-way branch mirrors retroverify's (prune, promote, keep) but with substantively different semantics: retroverify operates on already-committed memory and can afford upward-only drift with a stricter prune bar, whereas reconsideration operates on held-but-uncommitted entries and treats the absence of promotion as a vote to keep one more cycle rather than to delete. The asymmetry is intentional and makes the two passes composable rather than redundant.

\paragraph{Promotion preserves the full verifier record.}
When an entry is promoted, the resulting \texttt{EpisodicEntry} is not a partial object carrying only $u_{\text{stored}}$. It is constructed from the fresh verifier output and retains the full signal set of Section~\ref{sec:verifier}: the four uncertainty signals $(u_{\text{stored}}, u_{\text{token}}, u_{\text{dropout}}, u_{\text{internal}})$, the semantic similarity $s_{\text{avg}}$ and normalised entropy $h_{\text{norm}}$, the three grounding signals $(p_{\text{ground,mean}}, p_{\text{ground,max}}, p_{\text{ground,atomic}})$, the NLI signals $p_{\text{entail}}$ and $p_{\text{contra}}$, and the confabulation early-exit flag. This makes promotion behaviourally identical to a Stage~7 \textsc{Store} commit under the new verifier: the promoted entry is eligible for the same training-pool selection, retroverify-queue prioritisation, and ablation reporting as any other committed episode. Two audit-trail fields are preserved verbatim from push time rather than recomputed: the episode's \texttt{storage\_cycle} is held at the cycle at which it was first deferred (so the cycle-of-origin filter in the training pool and the audit logs see a stable origin), and the \texttt{source\_benchmark} tag is carried through so the ID/OOD gate does not degrade the entry to ``unknown origin'' purely because its commit was delayed by one cycle.

\paragraph{Novelty filter applies on promotion.}
Before the promoted entry is written, it passes through the same novelty filter of Equation~\ref{eq:novelty} that governs first-commit Stage~7 storage. Promotion is therefore not an unconditional insert: if the buffered entry is too similar to an existing memory entry (for instance because a later cycle's \textsc{Store} decision committed a near-duplicate in the interim) the insertion is suppressed. The suppressed case is still counted as a promotion rather than as a drop, because the buffer has accomplished its purpose --- the question is now covered by memory --- and re-queueing the entry would only force the reconsideration pass to re-decide the same duplication on the next cycle. This keeps the novelty guarantee of Section~\ref{sec:memory} invariant under the two entry points to the memory (direct Stage~7 commit and delayed reconsideration promotion): the memory never contains two entries within the similarity cap regardless of which path added them, while the promotion counter reported to Chapter~5 captures the full number of buffered entries that the fine-tuned verifier endorsed rather than only the subset for which a write actually occurred.

\paragraph{Scheduling and abort handling.}
Reconsideration runs as Step~8b of the cycle loop, strictly after retroverify (Step~8a), so that the novelty check at promotion time sees a memory that has already been reconciled under the new weights and so that the two passes share a single cycle-boundary abort gate. The pass is guarded by the same condition as retroverify: on cycles where the forgetting-ratio guard of Section~\ref{sec:self-improvement} rejects the fine-tune and rolls the weights back to $\theta^{(c-1)}$, reconsideration is skipped and the buffer is left untouched --- its entries retain their current $\mathrm{age}$ and are re-evaluated at the next successful cycle boundary. The rationale is identical to retroverify's: re-scoring under rolled-back weights reproduces storage-time composites, so the pass would trade compute for a no-op update. A verifier failure on an individual entry during reconsideration (exception, \texttt{None} return from the verify closure) is handled separately from the three-way rule: the entry is re-queued \emph{as-is} without an age bump and is counted toward $n_{\text{kept}}$, because the failure is attributable to the verifier and not to the entry. The orchestrator in \texttt{self\_improvement.py} exposes three counters: $n_{\text{promoted}}$ (entries for which the fine-tuned verifier endorsed \textsc{Store}, inclusive of novelty-suppressed writes), $n_{\text{dropped}}$ (entries evicted by the TTL rule), and $n_{\text{kept}}$ (entries re-queued for the next cycle, inclusive of verifier-failure re-queues).

\begin{algorithm}
\caption{Stage~8b deferred-entry reconsideration (runs once per cycle boundary, after retroverify, on non-aborted cycles).}
\label{alg:deferred-reconsider}
\begin{algorithmic}[1]
\Require Deferred buffer $\mathcal{B}$; memory store $\mathcal{M}$; fine-tuned verifier $V^{(c+1)}$; promote threshold $\tau_{\text{store}} = 0.65$; TTL $T = 2$.
\Ensure Buffer $\mathcal{B}$ and memory $\mathcal{M}$ updated in place; counters $n_{\text{pr}}, n_{\text{dr}}, n_{\text{kp}}$.
\State $n_{\text{pr}} \gets 0$; $n_{\text{dr}} \gets 0$; $n_{\text{kp}} \gets 0$
\State snapshot $\gets$ drain $\mathcal{B}$ into a list; clear $\mathcal{B}$ \Comment{rebuild from kept entries}
\ForAll{$b \in \text{snapshot}$}
    \State \texttt{try} $o^{\text{new}} \gets V^{(c+1)}.\text{verify}(b)$; \textbf{except} requeue $b$ into $\mathcal{B}$ without age bump; $n_{\text{kp}} \mathrel{+}= 1$; \textbf{continue}
    \If{$o^{\text{new}}$ is \texttt{None}}
        \State requeue $b$ into $\mathcal{B}$ without age bump; $n_{\text{kp}} \mathrel{+}= 1$; \textbf{continue} \Comment{verifier declined}
    \EndIf
    \If{$o^{\text{new}}.\text{decision} = \textsc{Store}$ \textbf{and} $o^{\text{new}}.u_{\text{stored}} \geq \tau_{\text{store}}$}
        \State $e^{\star} \gets \mathtt{EpisodicEntry}$ from $(b.\text{question}, b.\text{answer}, b.\text{embedding})$
        \State \hspace{1.2em} with full 11-signal record $\{u_\star, s_{\text{avg}}, h_{\text{norm}}, p_\star, \text{early\_exit}\}$ from $o^{\text{new}}$
        \State \hspace{1.2em} and audit fields $\text{storage\_cycle} \gets b.\text{storage\_cycle\_pushed}$, $\text{source\_benchmark} \gets b.\text{source\_benchmark}$, $\text{decision} \gets \textsc{Store}$
        \If{$\mathcal{M}.\text{is\_novel}(b.\text{embedding})$}
            \State $\mathcal{M}.\text{add}(e^{\star})$ \Comment{FAISS insert under novelty filter of Section~\ref{sec:memory}}
        \EndIf
        \State $n_{\text{pr}} \mathrel{+}= 1$ \Comment{near-duplicate write is still counted as a promotion}
    \ElsIf{$b.\text{age} + 1 \geq T$}
        \State $n_{\text{dr}} \mathrel{+}= 1$ \Comment{TTL expired; not re-queued}
    \Else
        \State $b.\text{age} \mathrel{+}= 1$; requeue $b$ into $\mathcal{B}$; $n_{\text{kp}} \mathrel{+}= 1$
    \EndIf
\EndFor
\State \Return $(n_{\text{pr}}, n_{\text{dr}}, n_{\text{kp}})$
\end{algorithmic}
\end{algorithm}

The triple $(n_{\text{pr}}, n_{\text{dr}}, n_{\text{kp}})$ is the per-cycle diagnostic of the reconsideration pass and complements the retroverify pair reported in the previous section. The expected trajectory under a healthy CAEM run is that $n_{\text{pr}}$ grows modestly across early cycles (answers the cycle-0 verifier was borderline about clear $\tau_{\text{store}}$ after one or two rounds of fine-tuning) while $n_{\text{dr}}$ stays comparable to the buffer's natural turnover under the TTL, and the buffer's age histogram concentrates at ages~0--1 rather than at $T-1$. A pathological trajectory in which $n_{\text{dr}}$ dominates is the counter-signal: the deferred band is accumulating answers the updated model still cannot commit to, and the composite axis's four-band design is not yielding the lifecycle it was designed for. Chapter~5's cycle-level tables report the per-cycle reconsideration triple alongside the retroverify pair so that the two cycle-boundary memory-update passes can be read as a coupled diagnostic rather than as independent counters.

\section{Self-Improvement Loop}
\label{sec:self-improvement}

Stage~8 is the mechanism that closes the outer loop of Figure~\ref{fig:caem-pipeline} and turns CAEM from a confidence-aware retrieval system into a self-improving one. At the end of each cycle it fine-tunes the Qwen-2.5-3B-Instruct base generator on the reasoning chains of the high-quality episodes in memory, regularised against a snapshot of the previous cycle's weights so that new knowledge accumulates without erasing the capabilities the model already had. The stage is deliberately conservative at three layers simultaneously: the training pool is strictly smaller than the Tier~1 retrieval pool (by design, not by accident), the regulariser is a uniform $L_{2}$ anchor rather than a full Fisher-weighted penalty, and the post-training capability check is out-of-distribution and measured against a \emph{pristine} Cycle-0 MMLU baseline that never updates across the run. If the check fails, the cycle is aborted, the weights are rolled back to the start-of-cycle snapshot, and all three cycle-boundary memory passes (retroverify, deferred-entry reconsideration, and consolidation) are skipped for that cycle. These choices together form the safety envelope inside which the compounding dynamics of Section~\ref{sec:retroverify} are allowed to operate.

\paragraph{Cold-start memory initialisation.}
Cycle~0 does not begin with an empty episodic memory. A one-time cold-start pass pre-populates the store by running the Qwen-2.5-3B-Instruct generator across the three ID benchmarks (FEVER, TriviaQA, Natural Questions) at a target of roughly $200$ verified episodes per benchmark, admitting the verifier's output under a \emph{loosened} cold-start override of the Stage-6 \textsc{Store} threshold, $\tau_{\text{store}}^{\text{cold}} = 0.50$ (versus the production $\tau_{\text{store}} = 0.65$). The pre-population runs under the same nine-signal verifier and four-outcome decision tree used online; only the admission threshold differs, and it is deliberately more permissive because the pre-SIL $\ustored$ distribution on a freshly-loaded backbone peaks lower than the post-calibration distribution for which $\tau_{\text{store}} = 0.65$ was chosen: at the production threshold the Cycle-0 pass would admit near-zero episodes and starve Tier~1 retrieval of an initial index. The $0.50$ value is itself a tightening from an earlier $0.45$ seed threshold (recorded in \texttt{caem/config.py} as the $\ustored$ shift that followed the prompt-redesign pass) chosen to reduce early-cycle noise in the seeded memory. The role of the seeded store is to give Tier~1 retrieval a non-empty index on the very first query of Cycle~0 so the routing decision can be exercised from the first sample rather than starting from an inevitable Tier-2 or Tier-3 fallback; once Cycle~0's online eval pass begins, the seed's contribution is immediately diluted by the cycle's own committed episodes, and from Cycle~1 onward the seed is indistinguishable from any other accumulated memory. The same seeded store feeds both the main CAEM run and every ablation variant, so cold-start state is a shared constant across comparisons rather than a confounder.

\paragraph{Training-pool curation.}
At the start of each cycle, the training pool is drawn from memory by filtering on the composite:
\begin{equation}
\mathcal{D}_{\text{train}}^{(c)} \;=\; \{\,e \in \mathcal{M}^{(c)} \;|\; u_{\text{stored}}(e) \;\geq\; \tau_{\text{train}}\,\}, \qquad \tau_{\text{train}} = 0.75.
\label{eq:train-pool}
\end{equation}
The training threshold $\tau_{\text{train}} = 0.75$ is stricter than the \textsc{Store} threshold $\tau_{\text{store}} = 0.65$ from Section~\ref{sec:decision}. The 0.10 gap between them is an architectural commitment: a stored episode is allowed to serve as a Tier~1 retrieval candidate with a lower composite than it needs to serve as training data, because retrieval acts on the model's output but training acts on the model's weights and can cause persistent harm if the data is under-verified. The gap cleanly separates the Tier~1 quality bar from the training-pool quality bar. Within the eligible set, the training target is the episode's \texttt{reasoning\_chain} field rather than the \texttt{answer} field: fine-tuning on the full chain teaches the model \emph{how} to reason toward the correct answer, not merely to memorise answer strings, and it is the mechanism by which Tier~3 retrieval-augmented reasoning is internalised across cycles. Two implementation-level defences further guard the pool. A multi-source ID/OOD gate excludes any episode whose effective benchmark set $\{\text{source\_benchmark}\} \cup \text{merged\_source\_benchmarks}$ intersects the held-out transfer split (ASQA and TruthfulQA are transfer-only by Chapter~3's protocol); without this guard, a post-consolidation entry merged from FEVER and TruthfulQA would leak held-out content into training. A repetitive-loop filter further rejects degenerate chains that fail \texttt{is\_repetitive\_loop}, a test calibrated on pre-Branch-C Flan-T5 pilot data where roughly a third of \textsc{Store}d chains were ``yes yes yes \ldots''-style loops that would have propagated into the weights as verbatim training targets; such chains fall back to the episode's \texttt{answer} string alone, matching the empty-chain and too-short-chain fallback branches.

\paragraph{Episode mix with a rehearsal buffer.}
The verified-episode pool $\mathcal{D}_{\text{train}}^{(c)}$ is combined with a small fraction of raw $(\text{question}, \text{answer})$ pairs from a static \emph{rehearsal buffer} $\mathcal{D}_{\text{gen}}^{\text{avail}}$ before fine-tuning. The mixing rule uses a fixed ratio $r_{\text{gen}} = 0.10$:
\begin{equation}
|\mathcal{D}_{\text{gen}}^{(c)}| \;=\; \min\!\left(\; \left\lceil \frac{|\mathcal{D}_{\text{train}}^{(c)}| \cdot r_{\text{gen}}}{1 - r_{\text{gen}}} \right\rceil,\; |\mathcal{D}_{\text{gen}}^{\text{avail}}| \right),
\label{eq:mix-size}
\end{equation}
so that episodes remain at $(1 - r_{\text{gen}}) = 90\%$ of the final mix regardless of how large the episode pool has grown by the current cycle. The rehearsal buffer is constructed once at run start by \texttt{load\_general\_data} in \texttt{scripts/run\_experiment.py}: it loads $\lvert \mathcal{D}_{\text{gen}}^{\text{avail}}\rvert = 1000$ raw QA pairs from TriviaQA's train split (HuggingFace \texttt{trivia\_qa/rc.nocontext/train}) at a row-offset past the end of the SIL training pool, so the buffer is strictly disjoint from both the SIL pool and from every benchmark's evaluation fold; the rows carry only $(\text{question}, \text{answer})$ and no scaffolded reasoning chain, so the mix exposes the backbone to an un-scaffolded answer-only format alongside the scaffolded chain-of-thought format of the verified episodes. The buffer is a rehearsal pool in the continual-learning sense of Scialom et al.~\cite{scialom2022fine} and Jin et al.~\cite{jin2024forget} rather than experience replay of past CAEM outputs: the episodes themselves carry the self-improvement signal; the rehearsal buffer is the static anti-forgetting anchor that prevents each cycle's parameter update from drifting away from the backbone's pre-training answer-format distribution.

Drawing the rehearsal buffer from TriviaQA's train split is a practical choice: TriviaQA is format-matched to CAEM's factoid-QA output style, the train split is large enough ($\sim\!87{,}000$ rows) to leave an ample disjoint tail past the $\sim\!50{,}000$-row SIL pool, and no additional dataset dependency is introduced. The cost is that the rehearsal distribution is \emph{not} truly general-purpose: it overlaps the answer-format distribution of one of the three ID benchmarks, and a cycle-over-cycle rehearsal of TriviaQA-style factoids may mildly favour TriviaQA-family distributions over FEVER's claim-verification labels or ARC-Challenge's multiple-choice format. The effect is bounded by the $10\%$ mix ratio, but Chapter~6 reports the per-benchmark decomposition so that any TriviaQA-favourable bias in the headline numbers is surfaced rather than absorbed. A stricter experimental design would draw the rehearsal buffer from a held-out dataset outside the evaluation panel (SQuAD, Natural Instructions, or a mixture) to make the anti-forgetting anchor distributionally neutral; that change is scheduled as future work rather than applied retroactively to this thesis's reported numbers.

\paragraph{Decoder-only training format.}
The prompt-and-target format mirrors the inference-time Tier~2 and Tier~3 generation contract of Section~\ref{sec:router}: the training path is the direct inverse of the inference path under the same backbone. For every training pair the question is wrapped in Qwen's ChatML envelope via the shared \texttt{build\_tier2\_prompt} builder used at inference, the forced ``\texttt{Reasoning:}'' prefix is appended as a prefill, and the target is the reasoning chain (or the answer fallback on the empty/too-short/loop-filter branches) with any leading ``\texttt{Reasoning:}'' token stripped so it does not duplicate the prefix. Cross-entropy is computed on a label tensor that masks all prompt and prefix positions to the ignore index: the loss fires only on the reasoning-chain-plus-answer span, not on the ChatML scaffolding or the forced prefix. This single choice is what lets the same pretrained backbone serve both roles without requiring a separate instruction-tuning phase, and it is the mechanism by which Tier~3 retrieval-augmented reasoning, once verified and committed to memory, is internalised into the base weights in a form the Tier~2 scaffold can re-emit.

\paragraph{$L_{2}$ anchor regularisation.}
The fine-tuning loss is the standard token-level cross-entropy on the label-masked span, plus an $L_{2}$ anchor on the parameter distance from the previous cycle:
\begin{equation}
\mathcal{L}(\theta) \;=\; \mathcal{L}_{\text{CE}}\!\left(f_{\theta}(q),\, y\right) \;+\; \frac{\lambda}{2} \,\lVert \theta - \theta_{\text{prev}} \rVert_{2}^{2}, \qquad \lambda = 0.01.
\label{eq:finetune-loss}
\end{equation}
Here $\theta_{\text{prev}}$ are the weights snapshotted at the start of the current cycle (equivalently, the accepted checkpoint of the previous cycle), and the anchor penalises every parameter uniformly for drifting away from that snapshot. This is the $L_{2}$ simplification of elastic weight consolidation \citep{doi:10.1073/pnas.1611835114}: full EWC weights the per-parameter squared deviation by the diagonal of the Fisher information matrix, which requires an extra forward pass over the training set to estimate and roughly doubles the cost of a cycle. The uniform-weight simplification is a principled approximation when the Fisher matrix is close to uniform across parameters, a regime that holds reasonably well for an instruction-tuned decoder-only model fine-tuned on diverse QA chains. The thesis reports this choice as a documented design decision and the full-EWC comparison as a scheduled ablation in Chapter~5.

The design is additionally supported by direct empirical evidence at
the 3B instruction-tuned scale we use. \citet{mehta2023empirical}
report that pre-training flattens the loss landscape to the point that
the Fisher-weighting advantage of EWC, SI, and MAS over plain
$L_{2}$ becomes minimal or negative at BERT-Large scale, and
\citet{wu2022pretrained} reproduce the same pattern across five
pre-trained backbones. The most directly comparable continual-fine-tune
study, \citet{jin2024forget}, measures Flan-T5-Large: vanilla sequential
fine-tuning over thirty-six P3 tasks produces only a four-to-six-percent
EM drop, and mini-batch replay of eight examples every ten steps drives
forgetting below one percent. Our ten-percent replay share in
Equation~\ref{eq:mix-size} sits an order of magnitude above that
proven-sufficient point, and \citet{scialom2022fine} additionally
demonstrate that one-percent rehearsal retains 99.8 percent of the
multi-task upper bound on T0pp-11B. Taken together, these studies
establish that plain $L_{2}$ anchored to the previous cycle plus
ten-percent replay is expected to hold catastrophic forgetting below
one percent per cycle at the 3B instruction-tuned scale, which is
comfortably within the $\rho_{\min} = 0.93$ cycle-level retention bound
of Equation~\ref{eq:forgetting-guard}.

\paragraph{Adaptive threshold calibration.}
The three decision-tree thresholds ($\tau_{\text{store}}$, $\tau_{\text{defer}}$, $\tau_{\text{train}}$) are not fixed hyperparameters. Earlier design documents gave canonical values $(0.65, 0.45, 0.75)$ read off a RoBERTa-MNLI composite distribution; under the MiniCheck backend adopted in this thesis the composite peaks closer to $0.55$ rather than $0.85$, so inherited threshold values would admit approximately zero storage events in Cycle 0 and starve the self-improvement loop. The thesis therefore calibrates the three thresholds from data rather than fixing them by literature prior. The design-time hyperparameters are the target quantiles of the $\ustored$ distribution:
\begin{equation}
\tau_{\text{store}} = Q_{\alpha_S}(\ustored), \quad
\tau_{\text{defer}} = Q_{\alpha_D}(\ustored), \quad
\tau_{\text{train}} = Q_{\alpha_T}(\ustored),
\label{eq:threshold-calibration}
\end{equation}
where $(\alpha_S, \alpha_D, \alpha_T) = (0.70, 0.40, 0.90)$ encode the decision-tree's design commitments: \textsc{Store} catches the top thirty percent of the observed distribution, \textsc{Deferred} catches the next thirty percent, training-pool admission is reserved for the top ten percent, and \textsc{Abstain}/\textsc{Discard} absorb the remaining forty percent. The quantiles are fit on the $n_{\text{cal}} = 500$ calibration fold (drawn from the training splits of FEVER, TriviaQA, and Natural Questions, disjoint from the evaluation fold by the content-hash protocol of Chapter~5), the same fold that fits the temperature scalar $T^{*}$ of Section~\ref{sec:pre-route}. Fitting thresholds on the calibration fold rather than on the evaluation fold is a deliberate disjointness guarantee: a quantile tuned on the fold that later evaluates metrics would constitute peeking at the test set and would invalidate the headline comparisons of Chapter~5. An assertion in the calibration script enforces $\tau_{\text{train}} > \tau_{\text{store}} > \tau_{\text{defer}}$; if the observed distribution violates this ordering the run is aborted rather than silently shipping a broken decision tree.

The three thresholds are re-fitted \emph{every} cycle, not just at Cycle~0, because fine-tuning shifts the $\ustored$ distribution along with the generator's logit scale and holding the thresholds fixed would desynchronise the decision bands from the signal they are quantiles of. Concretely (\texttt{adaptive\_thresholds\_per\_cycle = True} in \texttt{caem.config}), at every cycle boundary the fresh quantile fit $\tau_{\alpha}^{\text{fit}}$ is blended with the previous cycle's accepted threshold under an exponential-moving-average:
\[
\tau_{\alpha}^{(c)} \;=\; \alpha_{\text{EMA}} \cdot \tau_{\alpha}^{(c-1)} \;+\; (1 - \alpha_{\text{EMA}}) \cdot \tau_{\alpha}^{\text{fit}}, \qquad \alpha_{\text{EMA}} = 0.70,
\]
so that a single noisy cycle cannot swing the decision tree discontinuously and the stationary-distribution assumption is relaxed to a slowly-evolving one rather than strictly fixed. The remaining architectural thresholds --- the routing weight $\lambda_{\text{route}} = 0.70$, the pre-route safety floor $u_{\text{pre,min}} = 0.60$, the forgetting-tolerance floor $\rho_{\min} = 0.93$, and the $L_{2}$ anchor strength $\lambda = 0.01$ --- remain fixed by design and are not re-fitted. Chapter~5 reports the per-cycle fitted threshold values, the observed $\ustored$ quantile summary statistics, and the implied \textsc{Store}/\textsc{Deferred}/\textsc{Abstain}/\textsc{Discard} admission rates alongside the verifier-calibration diagnostics.

The P70 / P40 / P90 quantile targets in
Equation~\ref{eq:threshold-calibration} are design choices motivated
by the decision tree's semantics (storage reserved for the top third
of the observed distribution, deferred buffering for the next third,
training-pool admission for the top decile), not derivations from the
purity theorem's $\alpha > \tfrac{1}{2}$ premise. Chapter~5 closes
this gap empirically: \cref{fig:alpha-vs-tau} plots the measured
$\alpha$ against a grid of candidate $\tau_{\text{store}}$ values at
Cycle~0 (the calibration-matched distribution the threshold is fit on)
and at Cycle~$N$ (the converged distribution after self-improvement),
with the fitted P70 value marked. If $\alpha$ stays above
$\tfrac{1}{2}$ across $\tau_{\text{store}} \in [P_{60}, P_{80}]$ at
both endpoints the quantile choice is empirically robust; a crossing
within that band flags a sensitivity caveat for the narrative. The
analysis replays the decision tree at candidate $\tau_{\text{store}}$
values over per-sample scalars already recorded during the
purity-validation pass that populates Table~\ref{tab:purity}, so the
robustness check adds no inference compute beyond what the purity
validation already performs.

\paragraph{Cycle-2 retention diagnostic and structured fallback.}
The Jin 2024 one-percent forgetting expectation is a prior, not a
guarantee: CAEM's training pool is drawn from the model's own verified
outputs rather than from P3 and may induce a different drift profile.
We therefore commit, at design time, to a Cycle-2 hold-out diagnostic
as a conservative safety layer on top of the $\rho_{\min}$ forgetting
guard: a 500-example slice of Cycle-0 evaluation questions is frozen
before the first cycle begins and is re-evaluated on $\theta^{(2)}$ at
the end of Cycle 2. The slice is stratified across all seven benchmarks
(FEVER, TriviaQA, NQ, TruthfulQA, StrategyQA, ARC-Challenge, ASQA)
and is never used for training. If the absolute exact-match drop
against the Cycle-0 baseline on this slice exceeds three percentage
points, the pipeline logs a \emph{structured-fallback advisory}: the
remaining cycles should proceed under O-LoRA-style stacked orthogonal
adapters \citep{wang2023olora,biderman2024lora} with rank-sixteen
matrices on every attention and MLP module, merged every two cycles
to preserve some base-weight drift for the self-improvement
framing. \citet{wang2023olora} specifically test this on T5-Large
continual streams and report O-LoRA at sixty-nine percent of the
fifteen-task average accuracy against multi-task upper bound of
seventy-seven percent, cleanly Pareto-dominating both EWC
(forty-five percent) and replay (fifty-four percent) alone. If the
Cycle-2 drop is within three percentage points, plain $L_{2}$ plus
ten-percent replay is retained as the parsimonious choice, consistent
with the Mehta 2023 and Jin 2024 empirical regime. The three-percent
threshold is set against the combined Jin 2024 / Scialom 2022
upper envelope (both papers report sub-one-percent forgetting on
comparable continual streams), and is therefore a strict tightening
of the $\rho_{\min} = 0.93$ (seven-percent) cycle-level bound.

\paragraph{Training-path cascade.}
The implementation realises the Cycle-2 diagnostic as an explicit three-tier cascade over the training strategy rather than as a binary switch at cycle~2. The primary path is Full-parameter fine-tuning of the Qwen-2.5-3B-Instruct backbone with the $L_{2}$ anchor, using \texttt{bitsandbytes}'s 8-bit AdamW optimiser, gradient checkpointing, and bf16 mixed-precision autocast; the combination fits the weights, activations, optimiser state, and the resident $\theta_{\text{prev}}$ anchor within the 32\,GB envelope of a single RTX~5090. If a Cycle-2 advisory fires, or if a later cycle hits a Full-FT failure mode such as out-of-memory, retention below $\rho_{\min}$, or non-convergent loss, the harness flips \texttt{cfg.use\_lora\_training} and the next cycle switches to rank-sixteen LoRA adapters on every attention and MLP module with the $L_{2}$ anchor still applied against the (frozen) base weights, matching the merge-every-two-cycles schedule of \citet{wang2023olora,biderman2024lora}. If both Full-FT and LoRA fail on the same cycle, the outermost fallback is a memory-only cycle: the weights are left untouched, the cycle metadata records the fallback, and the self-consistent memory passes (retroverify, deferred-entry reconsideration, and consolidation) still run on the unchanged weights because their update is orthogonal to weight training. The cascade is not speculative machinery: each tier has a concrete trigger in the harness, and Chapter~5 reports which tier each cycle invoked alongside the per-cycle retention trajectory.

\paragraph{Out-of-distribution forgetting guard.}
After fine-tuning completes, the system evaluates the updated weights on a fixed 200-item MMLU split and computes a retention ratio against a \emph{pristine} Cycle-0 baseline measured once at the first cycle's start-of-cycle call and held unchanged for the rest of the run:
\begin{equation}
\rho_{\text{ret}}^{(c)} \;=\; \frac{\operatorname{MMLU}(\theta^{(c)};\, n{=}200)}{\operatorname{MMLU}(\theta^{(0)};\, n{=}200)}, \qquad \text{abort iff } \rho_{\text{ret}}^{(c)} \;<\; \rho_{\min}, \qquad \rho_{\min} = 0.93.
\label{eq:forgetting-guard}
\end{equation}
The denominator is the pristine Cycle-0 baseline, not the previous cycle's MMLU. Anchoring the denominator to the pristine baseline has a specific safety consequence: if drift accumulates gradually across several cycles while each individual cycle passes a naive ``vs.~previous cycle'' check, the cumulative erosion of general capability would be invisible to the guard. Anchoring against Cycle~0 makes the guard \emph{cumulative} --- the ratio captures drift from the pre-CAEM starting point, so a cycle that is individually small but that pushes the running trajectory below 93\% of the original MMLU triggers an abort even when its per-cycle delta looks innocuous. The guard is relative, not absolute, so a model that enters the run with a weak MMLU baseline does not trigger unwarranted aborts merely because the baseline was already low. The benchmark is deliberately out-of-distribution: MMLU's four-choice academic-subject questions are never in the training pool and share no topical overlap with the memory's accumulated QA distribution. This choice follows directly from the continual-learning framing of \citet{doi:10.1073/pnas.1611835114}: catastrophic forgetting in the continual-learning sense manifests as loss of \emph{unrelated} general capability, not as erosion of the target task, so an in-distribution probe (TriviaQA, Natural Questions) would be systematically insensitive to the failure mode the guard is supposed to detect. The same MMLU pass also feeds the RET axis of the Composite Evaluation Score (Chapter~3), so the guard incurs no additional evaluation compute: one benchmark pass serves simultaneously as the rollback trigger and as the reported general-capability retention metric.

\paragraph{Rollback and abort semantics.}
If $\rho_{\text{ret}}^{(c)} < \rho_{\min}$, the weights are restored to $\theta_{\text{prev}}$, the cycle is flagged aborted in the cycle record, and all three cycle-boundary memory passes --- retroactive re-verification (Section~\ref{sec:retroverify}), deferred-entry reconsideration (Section~\ref{sec:deferred-reconsider}), and the memory consolidation pass introduced next --- are skipped. The memory is not rolled back: episodes collected during the cycle remain stored, because their Stage~5 verification was performed under the previous-cycle weights and is therefore unaffected by the abort. An aborted cycle is not a failure of the overall pipeline but a commitment to the more conservative of the two available checkpoints; the next cycle restarts from $\theta_{\text{prev}}$ and sees an enlarged memory that includes the episodes just stored. This asymmetry (memory accumulates across aborts; weights do not) is deliberate: the abort path ensures that weight-level drift cannot compound across cycles, while memory-level knowledge persists and is available to future cycles even if the immediately-following fine-tune failed.

\paragraph{Memory consolidation (Step~8c).}
Retroverify and deferred reconsideration together reconcile \emph{individual} memory entries with the fine-tuned verifier. A third cycle-boundary pass, memory consolidation, operates on \emph{clusters} of near-duplicate entries that have accumulated across cycles and collapses each cluster into a single representative. The pass uses an SBERT cosine threshold of $\tau_{\text{cons}} = 0.92$ on the query-embedding space (the same embedding that drives the novelty filter of Section~\ref{sec:memory}) together with a $u_{\text{stored}}$-spread guard: a candidate cluster is only merged when all members are within $\Delta u_{\text{stored}} \leq 0.15$, so that a near-duplicate with materially different verification confidence is left as a distinct entry rather than collapsed into the representative. The merge elects the highest-$u_{\text{stored}}$ member as the representative, unions the \texttt{source\_benchmark} tags of all members into the representative's \texttt{merged\_source\_benchmarks} field (the tuple consumed by the multi-source ID/OOD gate defined above), and drops the non-representative members. Consolidation runs as Step~8c strictly after retroverify (Step~8a) and reconsideration (Step~8b), so that representative selection sees the freshest $u_{\text{stored}}$ values and avoids wasted work on entries that would have been pruned anyway. It is skipped on aborted cycles (nothing has changed since the previous cycle's consolidation) and is a no-op on Cycle~0 (cold-start memory is not collapsed on the first pass). The pass returns the pair $(n_{\text{cl}}, n_{\text{rv}})$ --- clusters merged and entries removed --- and Chapter~5 reports this pair alongside the retroverify counters and the reconsideration triple as a single cycle-boundary diagnostic block.

\paragraph{Checkpointing and multi-cycle trajectories.}
Each cycle saves a complete, independently-loadable checkpoint to a dedicated directory, including the model weights (or the restored $\theta_{\text{prev}}$ if the cycle aborted), the cycle metadata, and the random seed. The isolation has two operational consequences: a failed cycle cannot overwrite a successful predecessor, and any cycle's state can be reloaded in isolation for post-hoc comparison, ablation replay, or reproducibility audits. The outer loop is configured for $C = 10$ cycles in the code, all ten of which are evaluated end-to-end in Chapter~5; the per-cycle trajectory reported there is the full ten-cycle sweep. Across the evaluated trajectory, six quantities jointly characterise the cycle and are reported per-cycle in Chapter~5: the size of the training pool $|\mathcal{D}_{\text{train}}^{(c)}|$ (grows as the memory accumulates), the final training loss, the MMLU retention ratio $\rho_{\text{ret}}^{(c)}$ against the pristine Cycle-0 baseline, the retroverify counters $(n_{\text{up}}^{(c)}, n_{\text{rm}}^{(c)})$ from Section~\ref{sec:retroverify}, the reconsideration triple $(n_{\text{pr}}^{(c)}, n_{\text{dr}}^{(c)}, n_{\text{kp}}^{(c)})$ from Section~\ref{sec:deferred-reconsider}, and the consolidation pair $(n_{\text{cl}}^{(c)}, n_{\text{rv}}^{(c)})$ from the preceding paragraph. A healthy trajectory shows the training pool growing monotonically, the loss decreasing or stable, the retention ratio safely above $\rho_{\min}$, $n_{\text{up}}$ outpacing $n_{\text{rm}}$, and the consolidation counts tapering as the memory approaches a stable clustering.

\begin{algorithm}
\caption{Stage~8 self-improvement cycle (one invocation per cycle boundary).}
\label{alg:self-improvement}
\begin{algorithmic}[1]
\Require Current model $\theta^{(c-1)}$; memory store $\mathcal{M}^{(c)}$; deferred buffer $\mathcal{B}^{(c)}$; general pool $\mathcal{D}_{\text{gen}}^{\text{avail}}$; fine-tuned-verifier closure $V^{(c+1)}$; thresholds $\tau_{\text{train}} = 0.75$, $\rho_{\min} = 0.93$, $\tau_{\text{cons}} = 0.92$; hyperparameters $\lambda = 0.01$, $r_{\text{gen}} = 0.10$, epochs $E = 3$, learning rate $\eta = 10^{-5}$, batch $B = 16$; pristine anchor $m_{0}$ (persistent across cycles).
\Ensure Updated model $\theta^{(c)}$ (or restored $\theta^{(c-1)}$ on abort); cycle record.
\State $\mathcal{D}_{\text{train}}^{(c)} \gets \{e \in \mathcal{M}^{(c)} : u_{\text{stored}}(e) \geq \tau_{\text{train}}\}$ under OOD multi-source gate; target $= e.\text{reasoning\_chain}$, with \texttt{is\_repetitive\_loop}/empty/short fallback to $e.\text{answer}$
\If{$|\mathcal{D}_{\text{train}}^{(c)}| = 0$}
    \State \Return skip cycle (no eligible episodes)
\EndIf
\State sample $\mathcal{D}_{\text{gen}}^{(c)}$ from $\mathcal{D}_{\text{gen}}^{\text{avail}}$ of size given by Eq.~\ref{eq:mix-size}
\State $\mathcal{D}^{(c)} \gets \text{shuffle}(\mathcal{D}_{\text{train}}^{(c)} \cup \mathcal{D}_{\text{gen}}^{(c)})$ with ChatML-wrapped prompts, forced ``\texttt{Reasoning:}'' prefix, and label-masked targets
\State $\theta_{\text{prev}} \gets \theta^{(c-1)}$ \Comment{snapshot for $L_{2}$ anchor and rollback}
\State \textbf{if} $m_{0}$ unset \textbf{then} $m_{0} \gets \operatorname{MMLU}(\theta_{\text{prev}};\, n{=}200)$ \Comment{pristine Cycle-0 anchor, set once}
\State select training path: Full-FT + 8-bit AdamW (primary) $\to$ LoRA-16 (fallback~1) $\to$ memory-only (fallback~2)
\For{$k = 1$ to $E$}
    \ForAll{minibatch $\subset \mathcal{D}^{(c)}$}
        \State compute $\mathcal{L}(\theta)$ by Eq.~\ref{eq:finetune-loss}; AdamW step with grad-norm clip at $1.0$
    \EndFor
\EndFor
\State $m_{\text{post}} \gets \operatorname{MMLU}(\theta;\, n{=}200)$; $\rho_{\text{ret}} \gets m_{\text{post}} / m_{0}$ \Comment{denominator is pristine $m_{0}$, not $\theta_{\text{prev}}$}
\If{$\rho_{\text{ret}} < \rho_{\min}$}
    \State $\theta \gets \theta_{\text{prev}}$; save checkpoint; mark $\texttt{aborted} \gets \textsc{true}$; \Return \Comment{three memory passes skipped}
\EndIf
\State save checkpoint
\State $(n_{\text{up}}, n_{\text{rm}}) \gets \mathcal{M}^{(c)}.\text{retroverify}(V^{(c+1)})$ \Comment{Step 8a, Algorithm~\ref{alg:retroverify}}
\State $(n_{\text{pr}}, n_{\text{dr}}, n_{\text{kp}}) \gets \mathcal{B}^{(c)}.\text{reconsider}(V^{(c+1)}, \mathcal{M}^{(c)})$ \Comment{Step 8b, Algorithm~\ref{alg:deferred-reconsider}}
\If{$c > 0$ \textbf{and} \texttt{enable\_consolidation}}
    \State $(n_{\text{cl}}, n_{\text{rv}}) \gets \mathcal{M}^{(c)}.\text{consolidate}(\tau_{\text{cons}})$ \Comment{Step 8c, see \S\,Memory consolidation}
\EndIf
\State \Return updated $\theta^{(c)}$
\end{algorithmic}
\end{algorithm}

The compounding claim that motivates the cycle structure is a prediction, not a construction: nothing in the algorithm forces the model to improve cycle-over-cycle. What the algorithm does force is that improvement, when it happens, is attributable to the three signals the architecture controls --- the verified training pool curated from memory, the retroactive re-verification that follows fine-tuning, and the consolidation pass that collapses near-duplicate entries once both prior passes have settled --- rather than to any signal in isolation. The trajectory Chapter~5 reports therefore disaggregates the signals: an ablation in which fine-tuning is disabled isolates the memory's direct effect through Tier~1 retrieval, an ablation in which retroverify is disabled isolates the fine-tuning's direct effect on new generations, and an ablation in which consolidation is disabled isolates the clean-settled-set contribution to the training-pool curation. The joint configuration with all three channels active is the CAEM-complete condition, and its cycle trajectory is the headline figure of the Chapter~5 results.

\section{Theoretical Analysis}
\label{sec:theoretical-analysis}

The preceding sections specified CAEM as a set of decision rules, threshold gates, and update procedures. This section provides the formal underpinning for the central engineering claim those rules are designed to support: that storing only verified episodes and training the model on them produces a pool whose correctness rate is strictly higher than the model's unconditional accuracy, that the improvement is monotone in the base-model accuracy over repeated cycles, that the recurrence converges rather than diverging, and that the user-facing hallucination rate of the composed system inherits an architectural bound whose asymptotic floor is pinned to the universal epistemic limit imposed by corpus coverage and base-generator capacity. Five theoretical results, arranged in two groups, formalise this claim. The first group (Theorems~\ref{thm:purity}, \ref{thm:monotonicity}, and \ref{thm:convergence}) bounds the \emph{training-pool purity} --- the correctness rate of the stored-and-accepted set --- as a static identity, as a trajectory claim across cycles, and as a long-run recurrence result. The second group (Theorem~\ref{thm:asymptotic-elimination} together with Corollaries~\ref{cor:joint-optimum}--\ref{cor:self-correction}) lifts the pool-purity result to the \emph{user-facing hallucination rate} of the full CAEM pipeline: it decomposes that rate by routing tier, names its asymptotic floor in terms of quantities external to CAEM's design (the query distribution, the retrieval corpus, and the base-generator parameter count), and proves that hallucinated episodes which briefly evade the storage gates at a given cycle are self-pruned by the improving fine-tuned verifier in finite cycles. The two groups share the same two input quantities $p$ and $\alpha$ defined below, so we state the definition once and refer back to it throughout.

\begin{definition}[Base-model accuracy and verifier balanced accuracy]
\label{def:p-alpha}
Fix a cycle $c$ and a finite evaluation set $\mathcal{Q}$ of questions drawn i.i.d.\ from the task distribution. Let $y(q)$ denote the gold label for question $q \in \mathcal{Q}$ and let $\hat{y}_\theta(q)$ denote the model's greedy answer under parameters $\theta$. The base-model accuracy is
\[
p \;=\; \frac{1}{|\mathcal{Q}|}\sum_{q \in \mathcal{Q}} \mathbf{1}\!\left[\hat{y}_\theta(q) = y(q)\right].
\]
Write $A(q) = \mathbf{1}[\hat{y}_\theta(q) = y(q)]$ for the per-query correctness indicator. Let $V(q) \in \{0, 1\}$ denote the binarised verifier output, where $V(q) = 1$ means the four-outcome decision emitted \textsc{Store} (the three non-\textsc{Store} outcomes \textsc{Deferred}, \textsc{Abstain}, and \textsc{Discard} all map to $V(q) = 0$, since only \textsc{Store} entries enter the training pool). The verifier balanced accuracy on $\mathcal{Q}$ is
\[
\alpha \;=\; \tfrac{1}{2}\Pr\!\big[V=1 \mid A=1\big] + \tfrac{1}{2}\Pr\!\big[V=0 \mid A=0\big].
\]
\end{definition}

\begin{remark}[Operating point]
\label{rem:operating-point}
The pool-purity identity in Theorem~\ref{thm:purity} is derived under the symmetric operating-point assumption $\Pr[V=1 \mid A=1] = \alpha$ and $\Pr[V=1 \mid A=0] = 1 - \alpha$. In practice this is enforced by choosing the \texttt{retroverify\_prune\_threshold} so that the verifier's true-positive and true-negative rates are approximately equal on the labelled purity-validation set; departures from symmetry generalise the identity to the asymmetric form $P = p\cdot\mathrm{TPR} / (p\cdot\mathrm{TPR} + (1-p)\cdot\mathrm{FPR})$, which is the form the validation script evaluates when the operating-point check fails. The symmetric form is retained below as the primary statement because the threshold is calibrated to realise it.
\end{remark}

\subsection{Data Purity Theorem}
\label{sec:purity-theorem}

\begin{theorem}[Pool purity]
\label{thm:purity}
Under Definition~\ref{def:p-alpha} and the symmetric operating-point assumption of Remark~\ref{rem:operating-point}, the probability that an accepted episode is correct is
\[
P \;=\; \Pr\!\big[A = 1 \mid V = 1\big] \;=\; \frac{p\,\alpha}{p\,\alpha + (1 - p)(1 - \alpha)}.
\]
The pool is strictly purifying, meaning $P > p$, if and only if $\alpha > \tfrac{1}{2}$.
\end{theorem}

\begin{proof}
By Bayes' rule,
\[
    \Pr\!\big[A = 1 \mid V = 1\big] \;=\; \frac{\Pr\!\big[V = 1 \mid A = 1\big]\,\Pr\!\big[A = 1\big]}{\Pr\!\big[V = 1\big]}.
\]
The marginal probability of acceptance decomposes by the law of total probability as
\[
    \Pr[V = 1]
    \;=\; \Pr[V = 1 \mid A = 1]\,\Pr[A = 1] \;+\; \Pr[V = 1 \mid A = 0]\,\Pr[A = 0]
    \;=\; p\alpha + (1 - p)(1 - \alpha),
\]
where the second equality substitutes $\Pr[A=1] = p$ together with the symmetric operating-point values $\Pr[V = 1 \mid A = 1] = \alpha$ and $\Pr[V = 1 \mid A = 0] = 1 - \alpha$. Substituting back into Bayes' rule yields the stated identity for $P$.

For the purification claim, a direct algebraic simplification gives
\begin{align*}
    P - p
    &\;=\; \frac{p\alpha}{p\alpha + (1-p)(1-\alpha)} \;-\; p \\[2pt]
    &\;=\; \frac{p\alpha \;-\; p\bigl[p\alpha + (1-p)(1-\alpha)\bigr]}{p\alpha + (1-p)(1-\alpha)} \\[2pt]
    &\;=\; \frac{p(1-p)(2\alpha - 1)}{p\alpha + (1-p)(1-\alpha)}.
\end{align*}
Since $p \in (0, 1)$ and the denominator is strictly positive, the sign of $P - p$ is controlled by the factor $2\alpha - 1$. Hence $P > p$ if and only if $\alpha > \tfrac{1}{2}$, with equality $P = p$ precisely at the chance threshold $\alpha = \tfrac{1}{2}$, where a verifier that is no better than random contributes no purification.
\end{proof}

The identity has a direct engineering reading. The storage gate is a filter whose output distribution is a function of two inputs the pipeline controls independently: $p$ through whatever mechanism raises the base-model accuracy (Tier-1 retrieval hits, Tier-3 grounded generation, or accumulated fine-tuning), and $\alpha$ through the verifier's threshold and signal weights. The theorem says these two contributions compose deterministically to determine pool purity, which means the Chapter~5 empirical test reduces to three measurements per cycle rather than a joint sweep. Chapter~5 reports the comparison $|P_{\text{theory}} - P_{\text{obs}}|$ cycle by cycle in Table~\ref{tab:purity}.

\begin{remark}[Worked values]
At $p = 0.60$ and $\alpha = 0.80$, Theorem~\ref{thm:purity} predicts $P = 0.60 \cdot 0.80 / (0.60 \cdot 0.80 + 0.40 \cdot 0.20) = 0.48 / 0.56 \approx 0.857$. Raising $\alpha$ to $0.85$ at the same $p$ raises $P$ to $\approx 0.895$; raising $p$ to $0.70$ at $\alpha = 0.80$ raises $P$ to $\approx 0.903$. The two inputs are complements rather than substitutes, which is what justifies investing in both a stronger base model and a stronger verifier rather than one or the other.
\end{remark}

\begin{remark}[Empirical scope of $\alpha$]
\label{rem:alpha-empirical-scope}
Theorem~\ref{thm:purity} is a statement about two empirical quantities, $p$ and $\alpha$, both of which are measured rather than assumed. The theorem is therefore robust to any particular choice of verifier mechanism in the sense that its conclusion follows whenever $\alpha > \tfrac{1}{2}$ holds on the claim distribution actually scored. The engineering question is whether that inequality is in fact achievable on language model generated claims, which is a different distribution from the human written sentence pairs on which generic natural language inference models such as \texttt{roberta-large-mnli} were trained. Recent work on hallucination detection \citep{semanticillusion2025} reports that a generic NLI model with near state of the art held out accuracy on MultiNLI can nonetheless attain a one hundred percent false positive rate at ninety five percent recall on language model hallucinations, which corresponds to $\alpha \le 0.475$ under the symmetric operating point and therefore falsifies the assumption of Theorem~\ref{thm:purity}. The implementation response is to measure $\alpha$ directly on a stratified labelled set each cycle and to require $\alpha > \tfrac{1}{2}$ empirically before the storage policy is permitted to write new entries; Section~\ref{sec:implementation} describes the calibration diagnostic and the configuration switch that selects between a generic NLI backend and MiniCheck-Flan-T5-Large \citep{tang2024minicheck}, which is trained on language model generated claim support data and is the default backend from Cycle 0 onward. Chapter 5 reports the measured $\alpha$ per cycle alongside $|P_{\text{theory}} - P_{\text{obs}}|$, so that the theorem is audited end to end rather than accepted on the basis of the training distribution of any particular NLI model. The binary pass criterion for the theorem under the ten-cycle trajectory is $\alpha_c > \tfrac{1}{2}$ for \emph{every} cycle $c \in \{0, 1, \ldots, 10\}$, not merely on the mean: a single cycle falling to $\alpha_c \le \tfrac{1}{2}$ indicates that the verifier has degraded under fine-tuning on that cycle and that the pool is no longer purifying at that step, even if adjacent cycles retain $\alpha > \tfrac{1}{2}$. The reporting convention for the per-cycle trajectory is fixed in Chapter 5.
\end{remark}

\subsection{Monotonicity of Purity}
\label{sec:monotonicity-theorem}

\begin{theorem}[Monotone purification across cycles]
\label{thm:monotonicity}
Let $p_c$ denote the base-model accuracy at cycle $c$ and $\alpha_c$ the verifier balanced accuracy at the same cycle, and let $P_c = p_c \alpha_c / (p_c \alpha_c + (1 - p_c)(1 - \alpha_c))$ denote the purity predicted by Theorem~\ref{thm:purity}. Fix $\alpha_c = \alpha > \tfrac{1}{2}$ across cycles. Then $P_c$ is strictly increasing in $p_c$ on $(0, 1)$.
\end{theorem}

\begin{proof}
Write $P(p) = p\alpha \bigl/ \bigl[p\alpha + (1-p)(1-\alpha)\bigr]$ and let $D(p) = p\alpha + (1-p)(1-\alpha)$ denote the denominator. Differentiating with respect to $p$ at fixed $\alpha$ via the quotient rule,
\begin{align*}
    \frac{\mathrm{d}P}{\mathrm{d}p}
    &\;=\; \frac{\alpha \cdot D(p) \;-\; p\alpha \cdot D'(p)}{D(p)^2}
    \;=\; \frac{\alpha\bigl[p\alpha + (1-p)(1-\alpha)\bigr] \;-\; p\alpha\bigl[\alpha - (1-\alpha)\bigr]}{D(p)^2} \\[2pt]
    &\;=\; \frac{\alpha(1-\alpha)}{\bigl[p\alpha + (1-p)(1-\alpha)\bigr]^2}.
\end{align*}
For $\alpha \in (0, 1)$ the numerator $\alpha(1-\alpha)$ is strictly positive, and for $p \in [0, 1]$ the denominator $D(p)^2$ is strictly positive. Therefore $\mathrm{d}P / \mathrm{d}p > 0$ on $(0, 1)$, so $P_c$ is strictly increasing in $p_c$.
\end{proof}

The theorem's role in the cycle structure is to convert a per-step statement (Theorem~\ref{thm:purity}) into a trajectory claim. Whenever fine-tuning on the previous cycle's pool raises $p_c > p_{c-1}$ without degrading $\alpha_c$ below $\tfrac{1}{2}$, the next pool's purity is guaranteed to rise. The fixed-$\alpha$ assumption is not a requirement on the implementation but a simplification for the proof: the actual pipeline also runs retroactive re-verification under the updated model, which tends to raise $\alpha_c$ as well, so the observed trajectory should compound the two effects rather than isolate one. The Chapter~5 validation of this theorem is the purity trajectory across cycles $0\ldots N$ in Figure~\ref{fig:purity-trajectory}, together with a one-sided Spearman's rank correlation against the null hypothesis of no trend.

\subsection{Convergence of the Self-Improvement Recurrence}
\label{sec:convergence-theorem}

\begin{theorem}[Convergence of the purity recurrence]
\label{thm:convergence}
Consider the idealised recurrence in which the base-model accuracy at cycle $c + 1$ is bounded above by the previous cycle's pool purity,
\[
    p_{c+1} \;\le\; P_c,
\]
under the assumption that fine-tuning transmits the pool's correctness rate directly to the next-cycle generations. Fix $\alpha > \tfrac{1}{2}$ across cycles, and consider the map
\[
    f : [0, 1] \to [0, 1], \qquad f(p) \;=\; \frac{p\alpha}{p\alpha + (1 - p)(1 - \alpha)}.
\]
Then $f$ has exactly two fixed points, located at the endpoints $\{0, 1\}$, and for every starting point $p_0 \in (0, 1)$ the sequence $(p_c)_{c \ge 0}$ defined by $p_{c+1} = f(p_c)$ converges monotonically to $1$. The sequence of purity increments $\Delta P_c = P_c - P_{c-1}$ therefore satisfies $\Delta P_c \to 0$ as $c \to \infty$.
\end{theorem}

\begin{proof}
The argument has three parts: the fixed points, monotonicity and strict advancement, and convergence.

\emph{Fixed points.} Direct substitution gives $f(0) = 0$ and $f(1) = 1$, so both endpoints are fixed points. Conversely, solving $f(p) = p$ reduces to
\[
    p\alpha \;=\; p\bigl[p\alpha + (1-p)(1-\alpha)\bigr],
\]
which factors as $p(1-p)(2\alpha - 1) = 0$. Since $2\alpha - 1 \ne 0$ for $\alpha > \tfrac{1}{2}$, the only solutions in $[0, 1]$ are $p \in \{0, 1\}$.

\emph{Strict advancement on the open interval.} Theorem~\ref{thm:monotonicity} shows that $f$ is strictly increasing on $(0, 1)$. The identity established in the proof of Theorem~\ref{thm:purity},
\[
    f(p) - p \;=\; \frac{p(1-p)(2\alpha - 1)}{p\alpha + (1-p)(1-\alpha)},
\]
is strictly positive for $p \in (0, 1)$ whenever $\alpha > \tfrac{1}{2}$. Therefore $f(p) > p$ for every interior point.

\emph{Convergence.} Fix $p_0 \in (0, 1)$. The sequence $p_{c+1} = f(p_c)$ is strictly increasing (by the previous step) and bounded above by $1$ (since $f$ maps into $[0, 1]$). By the monotone convergence theorem, $p_c \to p^*$ for some $p^* \in [0, 1]$. Passing to the limit in $p_{c+1} = f(p_c)$ and using continuity of $f$ gives $p^* = f(p^*)$, so $p^*$ is a fixed point. The only fixed points are $\{0, 1\}$, and because the sequence is strictly increasing with $p_c > p_0 > 0$ for all $c$, the limit cannot be $0$. Therefore $p^* = 1$.

Since $P_c = f(p_c)$ and $f$ is continuous, convergence of $p_c$ to $1$ implies $P_c \to f(1) = 1$, and hence $\Delta P_c = P_c - P_{c-1} \to 0$.
\end{proof}

\begin{remark}[Ceiling in practice]
The idealisation $p_{c+1} \le P_c$ is an upper bound, not an equality: real fine-tuning transmits only a fraction of the pool's correctness rate to next-cycle generations, because the training set is finite, regularisation pulls the weights toward the base model, and the retention guard rolls weights back when MMLU drops. The effective ceiling observed in Chapter~5 is therefore strictly below $1$ and is set jointly by the training-pool purity threshold $\tau_{\text{train}}$ (which determines which of the stored episodes are used for training), the $L_2$ anchor strength (which determines how far weights can move), and the retention floor (which aborts training before diminishing returns would otherwise take over). These three factors are all properties of the training procedure and are in principle tunable. A fourth ceiling is external: the base generator's parameter capacity bounds base-model accuracy at $x(|\theta|, C) \leq 1$, the asymptotic correctness rate under retrieval corpus $C$, so fine-tuning on any training pool, however pure, cannot raise $p_{c}$ above this capacity-set bound. Corollary~\ref{cor:joint-optimum} below treats the parameter-capacity ceiling as the external constraint that distinguishes CAEM's architectural residual from a capacity-limited residual. Theorem~\ref{thm:convergence} guarantees the \emph{shape} of the trajectory, namely a monotone sequence with shrinking increments, but not the ceiling value itself. Chapter~5 reports both the trajectory (Table~\ref{tab:purity-convergence}) and the fitted ceiling in its convergence-check paragraph.
\end{remark}

Theorems~\ref{thm:purity}--\ref{thm:convergence} settle the behaviour of the \emph{training pool}, the set of stored-and-accepted episodes. What a user of CAEM actually experiences, however, is the system's answer to each arriving query. The next two subsections lift the pool-purity result to an architectural bound on the user-facing hallucination rate, decompose that bound by routing tier, and prove that hallucinated episodes do not persist in memory across cycles.

\subsection{Asymptotic Hallucination Elimination}
\label{sec:asymptotic-elimination}

Fix a query distribution $D$ and write $H_{t} = H(\text{CAEM}, D, t)$ for the user-facing hallucination rate of CAEM on $D$ at cycle $t$. We adopt the convention that a Stage-6 \textsc{Abstain} response does not count as a hallucination event: when the gate refuses to answer, the user sees a refusal, not a false claim. Under this convention, $H_{t}$ decomposes by routing tier as
\begin{equation}
H_{t} \;=\; \bigl(1 - \tau_{1}(t)\bigr) \cdot H_{\text{gen}}(t) \;+\; \tau_{1}(t) \cdot H_{\text{mem}}(t),
\label{eq:halluc-decomp}
\end{equation}
where $\tau_{1}(t) = \Pr_{q \sim D}[\text{Stage-3 routes } q \text{ to Tier~1}]$ is the Tier-1 hit rate, $H_{\text{gen}}(t)$ is the hallucination rate on Tier-2 and Tier-3 generated answers (conditioned on the gate not abstaining), and $H_{\text{mem}}(t)$ is the hallucination rate on Tier-1 retrieved answers.

\begin{theorem}[Asymptotic hallucination elimination]
\label{thm:asymptotic-elimination}
Assume Theorems~\ref{thm:purity}--\ref{thm:convergence} hold across cycles with $\alpha > \tfrac{1}{2}$, and assume the self-correction property of Corollary~\ref{cor:self-correction} below. Let $\tau_{1}^{*}(D) = \lim_{t \to \infty} \tau_{1}(t)$ denote the asymptotic Tier-1 hit rate of CAEM on $D$, and let $\varepsilon_{\text{arch}}$ denote the architectural Tier-2/3 hallucination ceiling imposed by the four cascaded storage gates of Section~\ref{sec:decision} (the Stage-6 confabulation early-exit, the composite $u_{\text{stored}}$ gate, the novelty filter of Section~\ref{sec:memory}, and the retroverify prune of Section~\ref{sec:retroverify}). Then
\begin{equation}
\lim_{t \to \infty} H_{t} \;\leq\; \bigl(1 - \tau_{1}^{*}(D)\bigr) \cdot \varepsilon_{\text{arch}}.
\label{eq:asymptotic-elimination}
\end{equation}
For any bounded query distribution with $\tau_{1}^{*}(D) = 1$ the limit is zero; for an open-domain distribution with $\tau_{1}^{*}(D) < 1$ the limit is a small architectural constant times the residual coverage gap.
\end{theorem}

\begin{proof}[Proof sketch]
Start from the routing decomposition of Equation~\ref{eq:halluc-decomp}. Three observations deliver the bound.

\emph{Memory term vanishes.} Corollary~\ref{cor:self-correction} shows that any hallucinated episode that briefly evades the storage gates at a given cycle is pruned by retroverify within a finite number of cycles under the improving fine-tuned verifier. Therefore $H_{\text{mem}}(t) \to 0$ as $t \to \infty$.

\emph{Generated term is architecturally bounded.} Each Tier-2 or Tier-3 answer shown to the user is scored by the four cascaded storage gates of Section~\ref{sec:decision}: the confabulation early-exit (driven by $u_{\text{internal}}$ and $p_{\text{ground,max}}$), the composite $u_{\text{stored}}$ gate (balanced accuracy $\alpha > \tfrac{1}{2}$ by Theorem~\ref{thm:purity}, so per-gate false-negative rate at most $1-\alpha$), the novelty filter (driven by the memory-side SBERT cosine), and the retroverify prune under the next-cycle verifier. The joint false-negative rate over the four gates is bounded above by a small architectural constant $\varepsilon_{\text{arch}}$ whose numerical value Chapter~5 estimates empirically from the Stage-6 disposition counts and the per-cycle retroverify outcomes in Table~\ref{tab:purity}.

\emph{Tier-1 hit rate saturates.} By construction, $\tau_{1}(t) \in [0, 1]$ is non-decreasing across cycles because each cycle commits new verified episodes to memory, retroverify prunes only entries the updated verifier now rejects, and consolidation preserves at least one representative per cluster. A monotone bounded sequence converges, so $\tau_{1}^{*}(D) = \lim_{t \to \infty} \tau_{1}(t)$ is well-defined.

Taking $t \to \infty$ in Equation~\ref{eq:halluc-decomp}: the first term tends to $(1 - \tau_{1}^{*}(D)) \cdot \varepsilon_{\text{arch}}$, the second to $\tau_{1}^{*}(D) \cdot 0 = 0$. Equation~\ref{eq:asymptotic-elimination} follows, and the bounded- and open-domain special cases follow by plugging in $\tau_{1}^{*}(D) = 1$ or $\tau_{1}^{*}(D) < 1$ respectively.
\end{proof}

Three properties of Theorem~\ref{thm:asymptotic-elimination} deserve explicit mention. First, $\varepsilon_{\text{arch}}$ is a property of the \emph{architecture}, not of the base generator's parameter count: the four gates score a three-billion-parameter generator's outputs and a seventy-billion-parameter generator's outputs through the same nine signals on the same (question, answer, passages) triple. Second, the bound is \emph{domain-independent} in that the same equation applies to claim verification, factoid QA, multi-hop QA, open-ended QA, and truthfulness benchmarks alike --- the query distribution $D$ affects only the numerical value of $\tau_{1}^{*}(D)$, not the functional form of the limit. Third, the bound structure is \emph{multiplicative}: residual hallucination is the product of the coverage gap and the architectural ceiling, so closing either factor to zero closes the limit. This is why CAEM invests simultaneously in retrieval growth (closing the coverage gap via memory accumulation, retroverify, and consolidation) and verifier quality (lowering $\varepsilon_{\text{arch}}$ via the cascaded gates) rather than relying on one alone.

\begin{corollary}[Joint optimum saturation]
\label{cor:joint-optimum}
Fix a query distribution $D$, a base generator with parameter count $|\theta|$, and a retrieval corpus $C$. Let $x(|\theta|, C) \in [0, 1]$ denote the base generator's asymptotic correctness ceiling on Tier-2/3 queries under retrieval corpus $C$. At equilibrium,
\begin{equation}
\lim_{t \to \infty} H_{t} \;=\; \bigl(1 - \tau_{1}^{*}(D, C)\bigr) \cdot \bigl(1 - x(|\theta|, C)\bigr) \;+\; O(\varepsilon_{\text{arch}}).
\label{eq:joint-optimum}
\end{equation}
The two residual factors are \emph{external} constraints on CAEM: $1 - \tau_{1}^{*}(D, C)$ is determined by the overlap between the query distribution and the retrieval corpus, and $1 - x(|\theta|, C)$ is determined by the base generator's parameter capacity. Neither factor can be reduced by architectural changes inside CAEM; the architecture attains the joint optimum given its inputs.
\end{corollary}

\begin{proof}[Proof sketch]
The Tier-1 hit rate $\tau_{1}^{*}(D, C)$ is saturated by the three cycle-boundary memory passes of Section~\ref{sec:self-improvement}: retroverify preserves committed episodes that the updated verifier still endorses, deferred-entry reconsideration promotes previously-borderline correct episodes once the verifier has improved, and consolidation collapses near-duplicates into a single representative whose $u_{\text{stored}}$ is the cluster maximum. Together these maintain maximal correct-answer coverage per unit of memory capacity, so the residual coverage gap $1 - \tau_{1}^{*}(D, C)$ cannot be reduced by further architectural work on the memory alone. The Tier-2/3 ceiling $x(|\theta|, C)$ is saturated by the verifier-reject path: the four cascaded gates reject a hallucinated generation with joint false-negative rate at most $\varepsilon_{\text{arch}}$, so among surviving Tier-2/3 answers the user-facing correctness rate is bounded above by the base generator's own ceiling $x(|\theta|, C)$ under retrieval. Combining the two and absorbing the $\varepsilon_{\text{arch}}$ contribution into a big-$O$ term yields Equation~\ref{eq:joint-optimum}.
\end{proof}

Corollary~\ref{cor:joint-optimum} is an interpretive result rather than an operational one. It reframes the residual hallucination at equilibrium from a CAEM shortcoming into a property of CAEM's external inputs: given a fixed base generator and a fixed corpus, no architectural change inside CAEM can further reduce $H_{t}$ at equilibrium. Reductions require either enlarging the corpus so that more queries fall within its coverage (shrinking $1 - \tau_{1}^{*}$) or training a more capable base generator (shrinking $1 - x(|\theta|, C)$). Both are outside CAEM's architectural scope.

\begin{corollary}[Base-model-conditional convergence rate]
\label{cor:base-conditional-rate}
Let $H_{\text{base}}(\theta) = 1 - x(|\theta|, C)$ denote the base generator's intrinsic Tier-2/3 hallucination rate and let $C_{\infty}(\theta)$ denote the asymptotic Tier-1 coverage CAEM attains under base generator $\theta$ and corpus $C$. Assume the cycle trajectory $H_{t}$ decays exponentially toward its equilibrium floor with rate $k > 0$. Then the number of cycles $c^{*}$ required to reach an $\epsilon$-neighbourhood of equilibrium satisfies
\begin{equation}
c^{*}(\theta, D) \;\approx\; \tfrac{1}{k} \,\ln\!\left(\tfrac{H_{\text{base}}(\theta)}{\epsilon \cdot C_{\infty}(\theta)}\right)
\;\propto\; \ln\!\left(\tfrac{H_{\text{base}}(\theta)}{C_{\infty}(\theta)}\right).
\label{eq:convergence-rate}
\end{equation}
\end{corollary}

\begin{proof}[Proof sketch]
Under the exponential-saturation form of Theorem~\ref{thm:convergence}, the gap to equilibrium decays as $\lvert H_{t} - H_{\infty}\rvert \sim A \cdot e^{-kt}$ with initial gap $A \propto H_{\text{base}}(\theta)$ (a stronger base generator starts closer to equilibrium) and asymptotic floor $H_{\infty} \propto 1 / C_{\infty}(\theta)$ (weaker coverage raises the floor). Solving $A \cdot e^{-k c^{*}} = \epsilon \cdot H_{\infty}$ and taking logarithms yields Equation~\ref{eq:convergence-rate}.
\end{proof}

The log-ratio form has an immediate engineering reading. A strong base generator (small $H_{\text{base}}$, large $C_{\infty}$) gives a small log ratio and therefore converges in few cycles; a weak base generator needs many cycles and plateaus higher. This matches the scaling pattern reported in recent studies that compare base, mid-trained, and post-trained generators on self-improvement-only loops; CAEM inherits the same rate behaviour because the memory-plus-verifier composition preserves the exponential-saturation form. Chapter~5 reports the fitted $(A, k, c^{*})$ values extracted from the observed CES trajectory via the equilibrium-fit diagnostic and positions the Qwen-2.5-3B-Instruct backbone within the saturation curve.

\subsection{Epistemic Floor and Self-Correction}
\label{sec:epistemic-floor}

The final two results frame CAEM's equilibrium residual as a universal epistemic bound rather than a design gap, and prove that hallucinated episodes which briefly evade the storage gates do not persist indefinitely.

\begin{corollary}[Corpus-bounded hallucination floor]
\label{cor:corpus-floor}
Let $K(C)$ denote the set of facts representable within the retrieval corpus $C$. For any grounded reasoning system $S$ that asserts only claims derivable from $C$,
\begin{equation}
H(S, D) \;\geq\; \Pr_{q \sim D}\!\bigl[\text{answer}(q) \notin K(C)\bigr] \;=:\; P_{\text{gap}}(D, C).
\label{eq:corpus-floor}
\end{equation}
CAEM attains this lower bound at equilibrium, with
\begin{equation*}
\lim_{t \to \infty} H_{t} \;=\; P_{\text{gap}}(D, C) \;+\; O(\varepsilon_{\text{arch}}).
\end{equation*}
\end{corollary}

\begin{proof}[Proof sketch]
The lower bound is a direct consequence of the grounded-reasoning restriction: a system that only asserts claims in $K(C)$ fails whenever the true answer is not in $K(C)$, regardless of how its internal weights, gates, or architecture are designed. The attainment statement combines Theorem~\ref{thm:asymptotic-elimination} and Corollary~\ref{cor:joint-optimum}: CAEM's equilibrium residual is the coincidence of missing memory coverage ($1 - \tau_{1}^{*}$) and missing generator capability over $C$ ($1 - x(|\theta|, C)$), which under the grounded-reasoning convention is within $O(\varepsilon_{\text{arch}})$ of the rate at which the true answer lies outside $K(C)$.
\end{proof}

Equation~\ref{eq:corpus-floor} describes a universal epistemic floor shared by every grounded reasoning system. A human historian in 2016 cannot answer a question whose answer first became public in 2024; a peer reviewer cannot cite work published after their submission deadline; a clinician cannot diagnose a disease characterised only after their training. CAEM's residual hallucination rate at equilibrium is of this kind: it is not an architectural deficiency but a knowledge-coverage limit. Reductions below this floor require expansion of the corpus $C$ or admission of ungrounded extrapolation, not further changes to CAEM's decision rules.

\begin{corollary}[Self-correction under parameter drift]
\label{cor:self-correction}
Fix a hallucinated episode $e_{h}$ that enters memory at Cycle $N$ with $u_{\text{stored}}(e_{h};\, \theta_{N}) \geq \tau_{\text{store}}$. Assume the Cycle-$N$ training-pool purity satisfies $P_{N} \geq \tfrac{1}{2}$ (the $\alpha > \tfrac{1}{2}$ condition of Theorem~\ref{thm:purity}) and that fine-tuning is monotone in the pool's correctness rate (Theorem~\ref{thm:monotonicity}). Then there exists a finite $k \geq 1$ such that
\begin{equation}
u_{\text{stored}}(e_{h};\, \theta_{N+k}) \;<\; \tau_{\text{retro}}, \qquad \tau_{\text{retro}} = 0.50,
\label{eq:self-correction}
\end{equation}
so $e_{h}$ is pruned from memory by the retroverify pass of Section~\ref{sec:retroverify} at Cycle $N + k$.
\end{corollary}

\begin{proof}[Proof sketch]
Because $P_{N} \geq \tfrac{1}{2}$, the fine-tune at Cycle $N$ shifts $\theta_{N}$ in the direction of correctness on the pool's distribution by Theorem~\ref{thm:monotonicity}, yielding $p_{N+1} \geq p_{N}$. The updated verifier $V^{(N+1)}$ therefore scores $e_{h}$ under parameters that are closer to the true-answer manifold on questions similar to $e_{h}$'s question, and the nine-signal composite shifts downward: the three grounding signals $p_{\text{ground,mean/max/atomic}}$ drop as the improved model now detects support mismatch; the entailment signal $p_{\text{entail}}$ and the semantic-coherence signal $s_{\text{avg}}$ drop as the reasoning chain no longer coheres; the internal-confidence signal shifts on an answer the model no longer endorses; the question-answer relevance signal $q\_a\_\text{relevance}$ tightens on an answer the improved model can now distinguish from a topical near-miss. The composite $u_{\text{stored}}$ is a convex combination with strictly positive weights of these signals and therefore decreases strictly as long as at least one signal moves. Under the monotonicity of Theorem~\ref{thm:monotonicity} across cycles, the sequence $u_{\text{stored}}(e_{h};\, \theta_{N+k})$ is non-increasing and crosses $\tau_{\text{retro}} = 0.50$ in finite cycles, at which point the retroverify pass of Section~\ref{sec:retroverify} removes $e_{h}$ from memory.
\end{proof}

Corollary~\ref{cor:self-correction} is strictly stronger than a static-gate argument. A static reading of CAEM would say: if a hallucination evades all four storage gates at Cycle $N$, it persists forever. The self-correcting reading says: evasion at Cycle $N$ is not terminal, because the very fine-tuning that the stored pool drives makes the Cycle-$(N+1)$ verifier measurably better at catching exactly the kind of mismatch $e_{h}$ represents. The result closes Theorem~\ref{thm:asymptotic-elimination}'s $H_{\text{mem}}(t) \to 0$ premise: hallucinated episodes do not accumulate in memory across cycles, so the Tier-1 contribution to user-facing hallucination asymptotically vanishes.

Taken together, the five results yield the chapter's central positioning. Theorems~\ref{thm:purity}--\ref{thm:convergence} prove that the training pool is strictly purifying at any single cycle, that purity rises monotonically across cycles as long as the verifier stays above chance, and that the resulting recurrence converges. Theorem~\ref{thm:asymptotic-elimination} together with Corollaries~\ref{cor:joint-optimum} and \ref{cor:base-conditional-rate} lifts these statements to the user-facing hallucination rate and names its log-ratio convergence speed. Corollary~\ref{cor:corpus-floor} identifies the asymptotic floor with the universal epistemic bound imposed by corpus coverage and base-generator capacity, and Corollary~\ref{cor:self-correction} closes the loop by proving that hallucinated episodes do not persist in memory. Chapter~5's trajectory figures and the purity, hallucination, and retroverify-disposition tables audit each result independently: a trajectory that does not rise (Theorem~\ref{thm:monotonicity}), rises indefinitely without diminishing returns (Theorem~\ref{thm:convergence}), does not saturate near the coverage-and-capacity floor (Corollary~\ref{cor:corpus-floor}), or accumulates rather than prunes hallucinated episodes (Corollary~\ref{cor:self-correction}) would falsify the corresponding result and require the implementation to be inspected for a mechanism operating outside the model the theorems assume.
% TODO(phase-3-coherence): cross-chapter integration hits for this section:
%   - Ch1 motivation: adopt the universal-epistemic-floor framing (C9) with
%     the human-historian / peer-reviewer / clinician analogies to frame
%     hallucination as a retrieval-bounded constraint rather than a model flaw.
%   - Ch5: add per-cycle tau_1(c) column to the retroverify-disposition table
%     (validates T4's Tier-1 saturation premise), a user-facing H(CAEM, B, c)
%     per-cycle table on the six evaluation benchmarks (validates T4's
%     monotone-decrease prediction), and a trace of the Hour-5 audit
%     hallucinations across Cycles 0..c* showing they get pruned within
%     2-3 cycles (validates C10).
%   - Ch6 discussion: integrate the three-limit separation (T4 benchmark-
%     fixed architectural, C7 parameter-bounded, double-limit composition)
%     as the thesis-positioning paragraph. Tracked under task #67.

\section{Implementation}
\label{sec:implementation}

The preceding sections describe CAEM as a set of decision rules, threshold gates, and update procedures. This section records the software and hardware substrate on which those rules are realised, so that any numerical result reported in Chapter~5 can be traced back to a specific library version, a specific model checkpoint, and a specific configuration value. Three organising principles govern the implementation: every numeric constant lives in a single configuration object so that no module silently hard-codes a threshold, concerns are separated by pipeline stage so that swapping a verifier or a router does not require touching unrelated components, and random seeds together with benchmark manifests are pinned so that a cycle run at month~$m$ can be repeated at month~$m + k$ with the same inputs.

The implementation is written in Python~$3.10$ and runs on PyTorch with the HuggingFace Transformers and Datasets libraries for model loading and tokenisation. Sentence-level embeddings are produced through the \texttt{sentence-transformers} package, nearest-neighbour search is delegated to \texttt{faiss-cpu}, and numerical utilities (temperature calibration fits, exponentially weighted updates, agglomerative clustering for semantic entropy) come from \texttt{numpy}, \texttt{scipy}, and \texttt{scikit-learn}. No custom CUDA kernels are used; every component runs through the public PyTorch API so that the system is portable across any CUDA-capable device that meets the memory requirement below.

CAEM composes three pretrained models that are treated as fixed assets for the duration of an experimental run. Only the Qwen-2.5-3B-Instruct backbone receives gradient updates in the self-improvement loop of Section~\ref{sec:self-improvement}; the sentence encoder and the verifier backend are frozen throughout. Table~\ref{tab:pretrained-components} summarises the three checkpoints, their parameter counts, and their role in the pipeline.

\begin{table}[htbp]
\centering
\small
\caption{Pretrained components and their role in the CAEM pipeline. Parameter counts are for inference; only the Qwen-2.5-3B-Instruct backbone is trainable during the self-improvement loop.}
\label{tab:pretrained-components}
\begin{tabular}{@{}p{3.4cm}p{2.2cm}p{2.0cm}p{5.0cm}@{}}
\toprule
\textbf{Component} & \textbf{Checkpoint} & \textbf{Parameters} & \textbf{Role} \\
\midrule
Base model & Qwen-2.5-3B-Instruct~\citep{qwen25} & $3.09$B & Decoder-only generation (Tier~2, Tier~3), MC~Dropout, self-consistency, semantic-entropy sampling; trainable via 8-bit AdamW full fine-tune with $L_{2}$ anchor. Legacy Flan-T5-Large \citep{JMLR:v25:23-0870} is retained as an ablation backbone for the cross-model C8 scaling check. \\
Sentence encoder & all-mpnet-base-v2 \citep{reimers-gurevych-2019-sentence} & $110$M & $768$-dim embeddings for memory keys, retrieval similarity, and semantic-entropy clustering; frozen. \\
Verifier backend & MiniCheck-Flan-T5-Large \citep{tang2024minicheck} & $770$M & Entailment and contradiction probabilities for $s_{\mathrm{avg}}$, $p_{\mathrm{entail}}$, $p_{\mathrm{ground}}$, and $p_{\mathrm{contra}}$; frozen. Legacy \texttt{roberta-large-mnli} \citep{liu2019roberta} is retained as an ablation backend; see Remark~\ref{rem:alpha-empirical-scope}. \\
\bottomrule
\end{tabular}
\end{table}

The hardware envelope targets a single consumer-grade NVIDIA accelerator with at most $32$~GB of VRAM paired with a Ryzen-class host CPU and $64$~GB of system RAM; the experiments reported in Chapter~5 were run on a rented RTX~5090 ($32$~GB), while the same configuration runs unchanged on an RTX~4090 ($24$~GB) by reducing the deferred-buffer capacity from $10^{4}$ to $5 \times 10^{3}$ and the static-batch size from $32$ to $16$. Qwen-2.5-3B-Instruct is loaded at \texttt{bfloat16} precision, and the self-improvement cycle uses three stacked techniques to fit a full-parameter fine-tune of a three-billion-parameter backbone into the $32$\,GB envelope alongside a frozen MiniCheck-Flan-T5-Large verifier: \texttt{bitsandbytes}'s 8-bit AdamW halves the optimiser-state footprint, gradient checkpointing trades activation memory for a second backward pass ($\approx 40\%$ savings for $\approx 20\%$ extra compute), and the $L_{2}$ anchor $\theta_{\text{prev}}$ (Eq.~\ref{eq:finetune-loss}) is moved to the same device as the training parameters at the start of the cycle so the anchor penalty is computed without host-device transfers per step. The FAISS index for episodic memory is held on the GPU for vector search and spilled to CPU only during the periodic save; the Wikipedia passage index used by Tier~3 RAG lives on CPU throughout because its size ($21$M vectors; see \Cref{sec:router}) exceeds what can be co-resident with the generator on a single-GPU budget. Attention throughput is recovered through a cuDNN-SDPA priority-order kernel selection (\texttt{torch.nn.attention.sdpa\_kernel([CUDNN, FLASH, EFFICIENT])}), measured at $67$~tok/s on Qwen-2.5-3B \texttt{bfloat16} versus $45$~tok/s under eager attention ($1.37\times$) on the target hardware.

\paragraph{Generator prompt design.} Every Tier~2 and Tier~3 generation call uses a uniform few-shot scaffolded chain-of-thought prompt with a forced \texttt{Reasoning:} prefill, so that the verifier downstream sees a substantive claim rather than a degenerate label token. The prompt template carries (i)~a one-line instruction to reason step by step before answering, (ii)~a single worked few-shot example in the exact \texttt{Reasoning: \ldots Answer: \ldots} format, (iii)~the retrieved Context block on Tier~3 calls (absent on Tier~2), (iv)~the task-specific instruction line with its constrained answer-format line (\texttt{supports | refutes | not enough info} for FEVER; \texttt{yes | no} for StrategyQA; \texttt{A | B | C | D} for ARC; \texttt{<concise factual answer>} for open-ended QA), and (v)~the response-format scaffold \texttt{Reasoning:} / \texttt{Answer:} that the decoder is asked to fill in. On the Qwen-2.5-3B-Instruct backbone the template is wrapped in Qwen's ChatML envelope: the system and user turns are constructed by the tokenizer's \texttt{apply\_chat\_template} call inside the shared \texttt{build\_tier2\_prompt} builder, and the forced \texttt{Reasoning:} prefix is concatenated onto the chat-wrapped string as a prefill rather than passed through a dedicated \texttt{decoder\_input\_ids} argument (the latter is a T5-specific API with no analogue on decoder-only backbones). The same builder is used at inference and at training time, so the Stage-8 fine-tune sees the identical scaffold the Tier~2/3 path emits; Section~\ref{sec:self-improvement} describes the training-side label masking that fires cross-entropy only on the reasoning-chain continuation.

This design is motivated by a pre-launch compliance smoke test on the Flan-T5-Large pilot (Branch~A/B, pre-Branch-C): under zero-shot scaffolding alone the pilot emitted the \texttt{Reasoning:} header in $0\%$ of queries across the six pilot-era benchmarks (the instruction-tuning prior overrode the abstract scaffold), adding a single few-shot example lifted compliance to $\sim\!57\%$, and additionally forcing the prefix lifted compliance to $\sim\!87\%$ averaged across benchmarks, with FEVER and Natural Questions reaching $100\%$ and ARC-Challenge remaining the only benchmark below $80\%$ (multiple-choice formatting pressures the decoder toward short letter answers even after the prefix). Branch~C's port to Qwen-2.5-3B-Instruct preserves all three layers (few-shot example, scaffold, forced prefill) and extends the evaluation panel to seven benchmarks with ASQA added; Chapter~5 re-measures the per-benchmark compliance profile on the new backbone and reports it alongside the Stage-6 disposition counts. Outputs flow through the standard \texttt{extract\_cot\_answer} helper used by the evaluation harness, which returns the text after the final \texttt{Answer:} marker (or the natural-language \texttt{the answer is X} suffix when present), so the user-facing display field in the pipeline result is the short answer even though the raw generation and the stored memory entry carry the full scaffolded text. The full scaffolded text is what the verifier sees, which is the whole point of the design: classification benchmarks (FEVER, StrategyQA, ARC) whose model predictions were previously one-word labels are now substantive propositional claims that the three external-grounding signals ($p_{\text{ground}}^{\text{max}}$, $p_{\text{ground}}^{\text{mean}}$, $p_{\text{ground}}^{\text{atomic}}$) can meaningfully score.

The code is organised by pipeline stage rather than by utility type. Eight submodules live under a top-level \texttt{caem} package: \texttt{memory} (episodic store, FAISS wrapper, novelty filter, pruner, consolidation pass, deferred-entry buffer), \texttt{confidence} (pre-routing estimator and the nine-signal post-generation estimator, each implemented as a separate class), \texttt{routing} (the adaptive router with its two-mechanism dispatch), \texttt{retrieval} (Tier~1 memory retrieval and Tier~3 passage retrieval, sharing the FAISS wrapper), \texttt{verification} (the \texttt{AdaptiveNLIJudge} ensemble composing MiniCheck, the frozen Qwen judge, and Platt calibration, plus the directional grounding, negation, and multi-choice scorers that feed the nine signals into the decision tree), \texttt{training} (self-improvement cycle, $L_{2}$ anchor, MMLU retention probe, and the repetitive-loop filter used during training-pool curation), \texttt{eval} (the benchmark harness that drives the Chapter~5 evaluation and ablation grid), and \texttt{ablation} (flags that disable individual mechanisms for the three-variant Branch-C panel reported in Chapter~5 and the broader, dormant flag registry that predates Branch~C). A top-level \texttt{pipeline.py} sequences the eight per-query stages under the serial contract; two sibling entry points \texttt{pipeline\_batch.py} and \texttt{pipeline\_batch\_prefetch.py} run the same stages in static batches for the Chapter~5 throughput profile described later in this section. A top-level \texttt{prompts.py} holds the shared ChatML builder \texttt{build\_tier2\_prompt} used by both the generator and the self-improvement training dataset, \texttt{model\_loader.py} handles backbone loading and the stale-CUDA-process reap, \texttt{benchmark\_splits.py} names the ID and transfer benchmarks consumed by the training-pool ID/OOD gate, and \texttt{config.py} holds the single configuration dataclass that every submodule reads from.

Every numeric constant in the system is annotated with one of three categories documenting its provenance. \textbf{Literature-derived constants} ([LIT]) are fixed from the cited paper and are not varied: examples are the MC-Dropout sample count $K = 5$ \citep{pmlr-v48-gal16}, the self-consistency chain count $M = 3$ \citep{wang2023selfconsistency}, the semantic-entropy sample count $K = 10$ at $T = 1.0$ \citep{farquhar2024semantic}, and the $768$-dimensional output of \texttt{all-mpnet-base-v2}. \textbf{Design-choice constants} ([DES]) are principled defaults chosen with reference to the thesis requirements and changed only with ablation evidence: examples are the Tier~1 and Tier~2 similarity thresholds, the novelty threshold $0.95$, and the composite weights for $u_{\text{stored}}$. \textbf{Calibration constants} ([CAL]) carry an initial value used for the warm-up cycles and are refitted on the Chapter~5 calibration set once the nine signals have accumulated enough observations for the fit to be meaningful; these include the temperature scalar for $u_{\text{pre}}$ calibration \citep{pmlr-v70-guo17a} and the post-generation composite weights if the calibration fit recommends deviating from the equal-weight initialisation. Table~\ref{tab:hparams} lists the full set of constants together with their category tags; any cell not traceable to this table is a bug in the implementation, and the configuration dataclass asserts on load that every field is present.

\begin{table}[p]
\centering
\scriptsize
\setlength{\tabcolsep}{4pt}
\caption{Consolidated hyperparameter reference for CAEM. Category tags follow Section~\ref{sec:implementation}: [LIT] = literature-fixed, [DES] = design choice, [CAL] = empirically calibrated. Initial values shown for calibration parameters; Chapter~5 reports the post-calibration values.}
\label{tab:hparams}
\begin{tabular}{@{}p{4.6cm}p{1.5cm}p{0.9cm}p{4.8cm}@{}}
\toprule
\textbf{Parameter} & \textbf{Value} & \textbf{Cat.} & \textbf{Role} \\
\midrule
\multicolumn{4}{@{}l}{\emph{Embedding and memory}} \\
SBERT model output dim & $768$ & [LIT] & \texttt{all-mpnet-base-v2} native dim \\
Memory capacity $M_{\max}$ & $10^6$ & [DES] & Full-scale run cap (\S\ref{sec:memory}) \\
FAISS \texttt{nlist} / \texttt{nprobe} & $4096 / 32$ & [DES] & IVF-PQ memory index \\
FAISS PQ $m$ / bits & $64 / 8$ & [DES] & IVF-PQ memory index \\
IVF-PQ training threshold $N_{\text{train}}$ & $20{,}000$ & [DES] & Promotion from flat to IVF-PQ \\
Novelty threshold $\tau_{\text{nov}}$ & $0.95$ & [DES] & Skip-storage cosine gate \\
Pruning trigger / amount & $0.95 / 0.20$ & [DES] & Value-score eviction \\
Recency decay $\lambda$ & $0.01$ & [DES] & $R = e^{-\lambda\,\Delta t}$ \\
Value weights $I / R$ & $0.70 / 0.30$ & [DES] & $V = 0.7\,I + 0.3\,R$ \\
\midrule
\multicolumn{4}{@{}l}{\emph{Pre-routing confidence}} \\
$u_{\text{pre}}$ weights (token / $C_{\text{conv}}$) & $0.60 / 0.40$ & [DES] & Two-signal estimator \\
Temperature $T$ & $1.0$ init & [CAL] & Single-parameter temperature scaling \citep{pmlr-v70-guo17a}; refit per cycle on the disjoint calibration slice \citep{ovadia2019trust,thulasidasan2019mixup} \\
\midrule
\multicolumn{4}{@{}l}{\emph{Adaptive router}} \\
Routing weight $\lambda$ & $0.70$ & [DES] & $S = \lambda\,s + (1-\lambda)\,\hat u$ \\
Tier~1 threshold $\tau_1$ & $0.90$ & [DES] & Route to retrieval \\
Tier~2 similarity floor $\tau_{\text{sim}}$ & $0.75$ & [DES] & Route to zero-shot \\
Safety floor $\tau_{\text{safe}}$ & $0.60$ & [DES] & OR-condition $\to$ Tier~3 \\
\midrule
\multicolumn{4}{@{}l}{\emph{Post-generation signals (Tier~2)}} \\
MC-Dropout passes $K$ & $5$ & [LIT] & \citet{pmlr-v48-gal16} \\
Self-consistency chains $M$ & $3$ & [LIT] & \citet{wang2023selfconsistency} \\
Semantic-entropy samples $K$ & $10$ at $T{=}1.0$ & [LIT] & \citet{farquhar2024semantic} \\
MC-Dropout rate & $0.1$ & [LIT] & \citet{pmlr-v48-gal16} \\
$\hat u$ accept threshold & $0.60$ & [DES] & Else escalate to Tier~3 \\
\midrule
\multicolumn{4}{@{}l}{\emph{$u_{\text{stored}}$ composite, 8 weights (sum to $1.0$)}} \\
$w_{\text{pground-mean}}$ & $0.22$ & [DES] & External grounding (mean over passages) \\
$w_{\text{pground-max}}$ & $0.18$ & [DES] & External grounding (max over passages; Branch~C) \\
$w_{\text{pground-atomic}}$ & $0.06$ & [DES] & Atomic-claim grounding, conservative floor \\
$w_{\text{q\_a\_rel}}$ & $0.20$ & [DES] & Question--answer relevance (Branch~C) \\
$w_{\text{entail}}$ & $0.14$ & [DES] & $p_{\text{entail}}$ (chain $\to$ answer) \\
$w_{\text{sc}}$ & $0.10$ & [DES] & $s_{\mathrm{avg}}$ pairwise SBERT \\
$w_{\text{uinternal}}$ & $0.08$ & [DES] & $\tfrac12 u_{\text{token}} + \tfrac12(1 - u_{\text{dropout}})$ \\
$w_{\text{se}}$ & $0.02$ & [DES] & $1 - h_{\text{norm}}$ (normalised semantic entropy) \\
\midrule
\multicolumn{4}{@{}l}{\emph{Verifier decision tree}} \\
Early-exit $u_{\text{internal}}$ floor & $0.70$ & [DES] & Confabulation gate \\
Early-exit $p_{\text{ground-max}}$ ceiling & $0.20$ & [DES] & Confabulation gate \\
Store threshold $\tau_{\text{store}}$ & $0.65$ & [DES] & \textsc{store} outcome \\
Defer threshold $\tau_{\text{defer}}$ & $0.45$ & [DES] & \textsc{deferred} outcome \\
Abstain ceiling $\tau_{\text{abs}}$ & $0.20$ & [DES] & \textsc{abstain} outcome \\
NLI retrieve / rerank $k$ & $20 / 3$ & [DES] & Grounding pipeline \\
\midrule
\multicolumn{4}{@{}l}{\emph{Tier~3 RAG}} \\
Passages retrieved $k$ & $5$ & [DES] & DPR convention \citep{karpukhin-etal-2020-dense} \\
Context budget (tokens) & $384$ & [DES] & DPR context cap for Tier~3 prompt \\
RAG / CoT max new tokens & $512 / 512$ & [DES] & Generation cap; raised from $256$ to accommodate the few-shot scaffolded CoT output (reasoning + answer) under the forced decoder prefix (\S~\ref{sec:implementation} \textit{Generator prompt design}) \\
\midrule
\multicolumn{4}{@{}l}{\emph{Self-improvement}} \\
Train threshold $\tau_{\text{train}}$ & $0.75$ & [DES] & Stricter than $\tau_{\text{store}}$ \\
General mix ratio & $0.10$ & [DES] & $90{:}10$ curated/general \\
L$_2$ anchor $\lambda$ & $0.01$ & [LIT] & EWC approx.\ \citep{doi:10.1073/pnas.1611835114} \\
Learning rate & $1{\times}10^{-5}$ & [DES] & AdamW \\
Batch size & $16$ & [DES] & Per-device \\
Epochs per cycle $E$ & $3$ & [DES] & Early-stopping-free \\
Warmup steps & $500$ & [DES] & Linear schedule \\
MMLU retention floor $\rho_{\min}$ & $0.93$ & [DES] & Abort \& rollback \\
Retroverify prune $\tau_{\text{retro}}$ & $0.50$ & [DES] & \S\ref{sec:retroverify} \\
Deferred-buffer capacity $N_{\text{buf}}$ & $10{,}000$ & [DES] & Bounded FIFO (\S\ref{sec:deferred-reconsider}) \\
Deferred-buffer TTL $T$ & $2$ cycles & [DES] & Max age before drop (\S\ref{sec:deferred-reconsider}) \\
Consolidation similarity $\tau_{\text{cons}}$ & $0.92$ & [DES] & SBERT-cosine merge gate, Step~8c (\S\ref{sec:self-improvement}) \\
Consolidation $u_{\text{stored}}$ spread cap & $0.15$ & [DES] & Max $\Delta u_{\text{stored}}$ within a merged cluster \\
8-bit AdamW & on & [DES] & \texttt{bitsandbytes} optimiser, primary Full-FT path \\
Gradient checkpointing & on & [DES] & Activation-memory saving for Full-FT \\
Feedback-loop $\eta$ & $0.01$ & [DES] & Exponential smoothing \\
Feedback min retrievals & $5$ & [DES] & Statistical stability \\
Cycles $C$ & $10$ & [DES] & Length of self-improvement run \\
Questions per benchmark per cycle & $3{,}000$ & [DES] & Chapter~5 protocol; $3{,}000 \times 3\;\text{ID benchmarks} = 9{,}000$ SIL-stream samples/cycle \\
\bottomrule
\end{tabular}
\end{table}

\paragraph{Performance engineering for experimental throughput.}
\label{par:perf-engineering}
The Chapter~\ref{ch:results} protocol is bounded by per-query wall-clock latency on the rented RTX~5090. At the per-benchmark per-cycle budget fixed in Chapter~5 ($n_{\text{questions}} = 3{,}000$), the ten-cycle SIL stream across three training benchmarks totals $3 \times 3{,}000 \times 10 = 90{,}000$ pipeline queries; the per-cycle seven-benchmark evaluation panel contributes $7 \times 500 \times 11 = 38{,}500$ queries (Cycles~$0$--$10$); and the MMLU retention probe contributes $200 \times 11 = 2{,}200$ queries --- roughly $130{,}000$ pipeline queries in total over the main run. A naive serial execution of the eight-stage pipeline cost $\approx 13$~s per query on the baseline profile before any batching (Profile~v3 on FEVER, see \Cref{tab:perf-profile-sequence}), which would extrapolate to $\approx 470$ GPU-hours at \$$0.80$/h --- prohibitively above the self-funded compute envelope. This paragraph records the \emph{static batching} layer that reduces per-query latency by $\approx 21\%$ to $\approx 10.2$~s / query and, once Tier-1 memory hits amortise across cycles, to an effective $\approx 7$--$8$~s / query averaged over the run --- without changing the decision rules or verifier signals described in \Cref{sec:verifier}. An additional
$\approx 12$-percentage-point speedup appeared achievable on paper by
pooling two further verifier stages ($u_{\text{token}}$
teacher-forcing and $u_{\text{dropout}}$ MC-sampling) across samples,
but live measurement on the target hardware showed the gain was
erased by a thermal regression in the cross-encoder rerank kernel
when the GPU ran at sustained $\sim\!100\%$ utilisation; those two
pools were therefore reverted, and their implementation is retained
in-file as infrastructure for future hardware that does not throttle
under the same workload. A third candidate --- atomic-decomposition
pooling --- was rejected earlier in the investigation for a
systematic correctness bias (see
\Cref{par:perf-atomic-rejection}).

\paragraph{Static batching wrapper.}
A wrapper class \texttt{BatchPipeline}
(\texttt{caem/pipeline\_batch.py}) processes $N$ queries in lockstep
through each of the eight pipeline stages. The serial
\texttt{CAEMPipeline.answer} entry point is preserved;
\texttt{answer\_batch(samples)} is a new entry point that runs
tier-dispatch per sample, buckets by tier, issues one padded
\texttt{model.generate} per tier bucket for Tier~2 and Tier~3
samples, runs one batched verifier pass over all Tier~2 and Tier~3
outputs, and finally commits each sample's decision through the
serial \texttt{answer} helper so memory-store writes happen in
submission order. Default batch size is $N = 32$, chosen to fit
Qwen-2.5-3B at \texttt{bfloat16} plus MiniCheck-Flan-T5-Large plus
the cross-encoder within the $32$~GB VRAM envelope. On an RTX~4090
($24$~GB) the same path runs unchanged at $N = 16$.

\paragraph{Within-verifier pooling (shipped).}
Two GPU-heavy stages inside a single \texttt{verify\_batch} call are
pooled across samples in the shipped configuration:
\begin{enumerate}
\item \textbf{M-chain self-consistency sampling}
  ($M = 3$ chains per sample at $T = 0.7$) and \textbf{semantic-entropy
  sampling} ($K = 10$ samples per sample at $T_{\text{SE}}$) pool
  $N \cdot M$ and $N \cdot K$ rows into one padded
  \texttt{model.generate(num\_return\_sequences=\cdot)} each,
  dispatched on two CUDA streams concurrently to overlap their
  kernel launches (\texttt{\_pooled\_sample\_dual}).
\item \textbf{Cross-encoder passage rerank} pools $N \cdot 20 = 640$
  (query+answer, passage) pairs at $N = 32$ into one
  \texttt{CrossEncoder.predict} call, eliminating the per-sample
  tokenisation and kernel-launch overhead
  (\texttt{\_retrieve\_and\_rerank\_batch}).
\end{enumerate}
The remaining two GPU-heavy stages --- $u_{\text{token}}$
teacher-forcing and $u_{\text{dropout}}$ MC-Dropout sampling --- run
per sample inside the serial \texttt{verify} loop from
\texttt{verify\_batch}. Cross-sample pooled implementations
(\texttt{\_pool\_u\_token\_batch}, \texttt{\_pool\_u\_dropout\_batch})
were written, correctness-validated by the diagnostic below, and then
reverted after live profiling on the target hardware (see
\Cref{par:perf-u-pools-thermal}). A third candidate --- pooling the
atomic-decomposition decoder (\texttt{\_score\_atomic\_batch}) ---
was rejected during equivalence validation for a systematic bias; see
\Cref{par:perf-atomic-rejection}.

\paragraph{Batching correctness: equivalence diagnostic.}
\label{par:perf-equivalence}
Each pool was validated through a per-signal equivalence diagnostic
(\texttt{scripts/diff\_verify\_serial\_vs\_batch.py}) that runs both
the serial \texttt{verify} and the batched \texttt{verify\_batch}
paths on the same deterministic (query, answer) pairs and diffs every
field of the resulting \texttt{UnifiedVerifierOutput}. The pass
criterion is two-part:
$|\text{mean}\ \Delta u_{\text{stored}}| < 0.01$ across the test set
AND no uni-directional sign pattern (e.g., $+N / -0$) on any
stage-isolable signal, because a balanced sign pattern indicates
\texttt{bfloat16} matmul reassociation noise whereas a uni-directional
pattern indicates a systematic algorithmic bias.
\Cref{tab:perf-pool-validation} summarises the per-pool outcome at
$N = 32$ on FEVER.

\begin{table}[tbp]
\centering
\small
\caption{Per-pool validation on the FEVER diagnostic at $N = 32$.
  Sign pattern counts samples with $\Delta > +0.005$, $\Delta < -0.005$,
  and $|\Delta| \le 0.005$ respectively. A pool is accepted when
  $|\text{mean}\ \Delta u_{\text{stored}}| < 0.01$ AND no signal
  exhibits a uni-directional sign pattern at the per-sample level.}
\label{tab:perf-pool-validation}
\begin{tabular}{@{}p{4.1cm}p{2.5cm}rrp{2.6cm}@{}}
\toprule
\textbf{Pool} & \textbf{Status} & \textbf{mean Δ} & \textbf{max |Δ|} & \textbf{Sign pattern} \\
\midrule
M-chain + SE (CUDA stream)        & shipped                                 & $s_{\text{avg}}\ +0.081$ & $0.65$ & $12 / 13 / 7$ balanced \\
Cross-encoder rerank              & shipped                                 & $\approx 0$ & $\approx 0$ & identical passages \\
$u_{\text{token}}$ (right-pad)    & correct but reverted\textsuperscript{*} & $+0.0004$   & $0.004$     & $0 / 0 / 32$ clean \\
$u_{\text{dropout}}$ (MC pool)    & correct but reverted\textsuperscript{*} & $-0.038$    & $0.40$      & $1 / 5 / 26$ conservative \\
Atomic decomposition              & \textbf{rejected}\textsuperscript{\dag} & $+0.229$    & $0.72$      & $\mathbf{4 / 0 / 0}$ systematic \\
\bottomrule
\end{tabular}\\[2pt]
\textsuperscript{*}Passes correctness bar but was reverted after live
profiling showed no wall-clock gain on the RTX~5090; see
\Cref{par:perf-u-pools-thermal}.\\
\textsuperscript{\dag}Rejected for systematic correctness bias; see
\Cref{par:perf-atomic-rejection}.
\end{table}

\paragraph{Reverted pools: $u_{\text{token}}$ and $u_{\text{dropout}}$ pooling.}
\label{par:perf-u-pools-thermal}
Pooled implementations of both internal-calibration stages passed the
equivalence diagnostic at $N = 32$ --- $|\text{mean}\ \Delta|$
entries of $0.0004$ and $0.038$ on $u_{\text{token}}$ and
$u_{\text{dropout}}$ respectively, with no uni-directional sign
pattern on either --- so the correctness bar was satisfied. Live
profiling on the RTX~5090 nonetheless showed the two pools did not
net-save wall-clock time: the $\approx 35$~s they save per batch of
$32$ samples in the per-sample verifier loop is offset by a
$\approx 50$~s regression in the cross-encoder rerank stage. The
regression appears to be thermal: when the rerank, M-chain / SE, and
two $u$-pool stages all run back-to-back at sustained $\sim\!100\%$
GPU utilisation without the per-sample Python orchestration gaps of
the serial path, the sustained power draw pushes the $5090$ into its
throttling envelope and the rerank kernel (which depends on the
longest-input row in the pooled batch) loses throughput before the
other stages feel it. The net per-query latency (Profile~v8,
\Cref{tab:perf-profile-sequence}) was $\approx 11.5$~s vs
Profile~v4's $10.2$~s --- a $\approx 13\%$ regression. Both pools
were therefore reverted; the serial \texttt{\_compute\_u\_token} and
\texttt{\_compute\_u\_dropout} paths run per sample from inside
\texttt{verify\_batch} as before. The helper functions remain in-tree
with dead-code markers so a future hardware platform with different
thermal behaviour can re-enable the pools without further code
change.

\paragraph{Rejected pool: atomic decomposition.}
\label{par:perf-atomic-rejection}
A pooled variant of \texttt{\_score\_atomic} would have saved
$\approx 40$~s per $N = 32$ batch by batching the greedy
atomic-decomposition generate and pooling all (passage, fact) NLI
pairs across samples. The diagnostic rejected it. Under batched
greedy decoding on Qwen-2.5-3B at \texttt{bfloat16} with $N \ge 4$
left-padded inputs, the decomposition decoder terminates $2$--$3$
tokens earlier than the unpadded serial path, emitting only $4$--$7$
atomic facts per sample vs the serial path's $13$--$14$. With fewer,
coarser facts, the $\min(p_{\text{entail}})$ reduction over per-fact
grounding probabilities lands $+0.23$ higher on average, inflating
$p_{\text{ground}}^{\text{atomic}}$ uniformly across all test
samples. The shift propagates through the composite weight
$w_{p_{\text{ground}}^{\text{atomic}}} = 0.14$ to $+0.05$ on
$u_{\text{stored}}$, which is large enough to move STORE / DISCARD
boundary cases. The sign pattern was $+4 / -0 / {\sim}0$ at $N = 4$:
all four samples shifted in the same direction, ruling out bf16
matmul noise as the explanation. The pool was therefore reverted
with a sentinel guard (\texttt{atomic\_facts\_pool[i] = None}) so the
serial path runs per sample from inside \texttt{verify\_batch}; the
helper itself is retained in-tree so a future decoder framework that
fixes EOS timing under batched greedy decoding can re-enable it
without further code change.

\paragraph{Rejected optimisation: \texttt{torch.compile}.}
Empirical drift measurement via
\texttt{scripts/compile\_drift\_check.py} on the same hardware stack
(torch~$2.10.0$+cu$130$, CUDA~$13.1$, Blackwell sm\_$120$, bf$16$,
Qwen-2.5-3B-Instruct) recorded a maximum logit drift of
$\mathbf{5.156}$ between eager and compiled forward passes under both
\texttt{compile\_mode="reduce-overhead"} and \texttt{="default"},
or $\approx 50{,}000{\times}$ the documented $10^{-4}$ envelope.
Because a $5$-unit logit drift is large enough to flip greedy
decoder choices --- exactly the failure mode that rejected the
atomic-decomposition pool above --- \texttt{use\_torch\_compile}
remains \texttt{False} in the default configuration. The gate and
baseline are recorded in \texttt{outputs/perf\_log.csv}; the
$\approx 10$--$15\%$ generation speedup torch.compile would offer is
not commensurate with a correctness-boundary change of that
magnitude.

\paragraph{Other Goal-5 levers.}
Two additional optimisations landed without triggering the
equivalence diagnostic because they do not alter the numerical
inference path. \textbf{$8$-bit AdamW} (via \texttt{bitsandbytes})
for the self-improvement loop halves the optimiser-state footprint
and makes the L$_2$ anchor plus frozen copy of the base model
co-resident on GPU under the $32$~GB envelope.
\textbf{Tier-adaptive FAISS \texttt{nprobe}} (Tier~2 queries use
$\texttt{nprobe} = 16$, Tier~3 uses $\texttt{nprobe} = 64$)
allocates retrieval accuracy to where it matters: Tier~3 RAG answers
are verifier-gated while Tier~2 answers are accepted pre-verifier
once the self-consistency similarity clears $\tau_{\text{sim}}$. An
\textbf{automatic stale-CUDA-process reap} is wired into
\texttt{caem/model\_loader.py} (gated on
\texttt{CAEM\_FORCE\_GPU\_CLEANUP=1}) so that a crashed prior run
does not force a manual
\texttt{nvidia-smi --query-compute-apps \ldots{} | kill -9} before the
next pipeline load.

\begin{table}[tbp]
\centering
\small
\caption{Measured verifier-batch wall-clock on the FEVER$@64$ smoke
  test (RTX~5090, bs $= 32$). Profile~v3 is the pre-pool baseline;
  \textbf{v4 is the shipped configuration}; v7 shows the
  atomic-pool-contaminated regime that the equivalence diagnostic
  rejected; v8 adds the two further $u$-pools (correct but
  thermal-regressed, subsequently reverted). Per-query latency
  includes all generation and verification stages.}
\label{tab:perf-profile-sequence}
\begin{tabular}{@{}llrrrr@{}}
\toprule
\textbf{Profile} & \textbf{Config} & \textbf{verify\_batch} & \textbf{Batch wall} & \textbf{s / query} & \textbf{vs v3} \\
\midrule
v3          & no verifier pools                              & $\sim\!246$~s & $\sim\!415$~s & $13.0$ & -- \\
\textbf{v4} & \textbf{$+$ m-chain, SE, rerank (shipped)}     & $\mathbf{216}$~\textbf{s} & $\mathbf{327}$~\textbf{s} & $\mathbf{10.2}$ & $\mathbf{21\%}$ \\
v7          & $+$ $u_{\text{token}}$, $u_{\text{dropout}}$, atomic  & $121$~s & $220$~s & $6.9$ & $47\%$\textsuperscript{\dag} \\
v8          & v7 $-$ atomic pool                             & $229$~s & $367$~s & $11.5$ & $12\%$\textsuperscript{\ddag} \\
\bottomrule
\end{tabular}\\[2pt]
\textsuperscript{\dag}Rejected by equivalence diagnostic
($+0.05\ u_{\text{stored}}$ bias; see
\Cref{par:perf-atomic-rejection}).\\
\textsuperscript{\ddag}Correctness-clean but thermal-regressed on the
$5090$; $u_{\text{token}}$ and $u_{\text{dropout}}$ pools reverted
post-profile (\Cref{par:perf-u-pools-thermal}). The shipped
configuration is v4.
\end{table}

Reproducibility is handled at three layers. A single integer seed is passed at pipeline construction and drives every stochastic component: the MC-Dropout masks, the self-consistency and semantic-entropy sampling, the training-pool shuffle, and the $10\%$ general-domain sampler inside \texttt{\_mix}. The configuration dataclass is serialised to JSON alongside every checkpoint so that a cycle-$c$ model can be replayed with exactly the constants it was trained under, even if the repository defaults change. Benchmark manifests (the exact question indices used in each cycle's $3{,}000$-question-per-benchmark SIL draw and each $200$-question MMLU probe) are written next to the checkpoint. The one component that remains non-deterministic is low-level CUDA arithmetic, which produces floating-point results that can differ in the final decimal places across runs even under a fixed seed; this is a limitation shared with every system using the same stack, and the variance it induces is smaller than the cycle-over-cycle effects reported in Chapter~5.

\section{Chapter Summary}
\label{sec:ch4-summary}

This chapter has operationalised every functional requirement stated in Chapter~\ref{ch:requirements} into a concrete mechanism of the CAEM pipeline. Section~\ref{sec:system-overview} gave the architecture on a single page, fixed the notation used in the rest of the chapter, and traced a single fabricated query end to end through all eight stages. Section~\ref{sec:pre-route} then specified Stage~1, the two-signal pre-routing confidence estimator $\upre$ that fuses token-level geometric-mean probability with conversation-context similarity in a single encoder forward pass. Section~\ref{sec:memory} specified the episodic memory store that holds Stage~2's retrieval target: a FAISS IVF-PQ index over SBERT embeddings, paired with the $\ustored$ composite stored alongside each vector, the cosine novelty filter, the value-score pruner, and the retrieval feedback loop that updates $\ustored$ as entries prove useful or harmful in downstream queries. Section~\ref{sec:router} combined Stages~2 and~3 into the adaptive router that dispatches each query to Tier~1 (cached answer), Tier~2 (zero-shot generation), or Tier~3 (retrieval-augmented generation) via a blended score $S = \lambda\,s + (1-\lambda)\,\hat u$ with a safety override that forces Tier~3 whenever both inputs fall below a minimum floor.

Sections~\ref{sec:verifier} and~\ref{sec:decision} specified the verifier-and-decision block that gates every Tier~2 and Tier~3 output before it can enter memory. The nine-signal verifier composes internal uncertainty, self-consistency, semantic entropy, NLI-based grounding against retrieved passages, and internal entailment into a single composite $\ustored \in [0, 1]$; a confabulation early-exit gate fires when an answer presents as internally confident while failing grounding. Section~\ref{sec:decision} then partitions the $\ustored$ axis into four decisions, \textsc{Store}, \textsc{Deferred}, \textsc{Abstain}, and \textsc{Discard}, and stated the per-cycle distribution over these outcomes as a primary Chapter~5 diagnostic. Sections~\ref{sec:retroverify} and~\ref{sec:deferred-reconsider} are the cycle-boundary counterparts of the verifier: the first re-scores every stored entry under the fine-tuned verifier and prunes, promotes, or keeps it via a three-way cascade on $\tau_{\text{retro}}$; the second re-scores the bounded FIFO of \textsc{Deferred} entries held from the previous cycle and promotes, drops (TTL-aged), or keeps each via a symmetric three-way cascade. Section~\ref{sec:self-improvement} closed the outer loop: a training pool is curated from retrieved-and-useful episodes that clear $\tau_{\text{train}}$ under a multi-source ID/OOD gate and a repetitive-loop chain filter, the Qwen-2.5-3B-Instruct backbone is fine-tuned for $E$~epochs under an $L_{2}$ anchor approximating EWC with an 8-bit AdamW optimiser and gradient checkpointing, and the cycle is aborted with a weight rollback whenever a held-out MMLU retention probe falls below $\rho_{\min}$ against the pristine Cycle-0 baseline. A three-tier training-path cascade (Full-FT $\to$ LoRA-16 $\to$ memory-only) protects the pipeline when any cycle's Full-FT hits an out-of-memory, retention, or convergence failure; all three cycle-boundary memory passes (retroverify, deferred-entry reconsideration, and consolidation) are skipped on aborted cycles. Section~\ref{sec:implementation} then recorded the software substrate (Python~$3.10$, PyTorch, three pretrained components, single-GPU envelope), the code organisation (seven stage-aligned submodules plus a single configuration dataclass), the hyperparameter provenance conventions ([LIT]/[DES]/[CAL]), and the reproducibility layers (fixed seed, serialised config, pinned benchmark manifests). Table~\ref{tab:hparams} consolidates every numeric constant in the system with its category tag.

Chapter~5 reports the experimental trajectory of the system specified here against the acceptance criteria of Chapter~\ref{ch:requirements} (Section~\ref{sec:nfr}) and the Composite Evaluation Score of Section~\ref{sec:ces}, together with the three-variant Branch-C ablation panel of Section~\ref{sec:ablation-design}. Chapter~6 then discusses those results in light of the catastrophic failure modes enumerated in Section~\ref{sec:risks} and the ethical considerations of Section~\ref{sec:impact}.
